\documentclass[11pt]{ctexart}
\usepackage[a4paper,margin=1in]{geometry}
\usepackage{graphicx}
\usepackage{xcolor}
\usepackage{hyperref}
\definecolor{refgray}{gray}{0.45}
\title{\textbf{市场最新观点汇总}\\[0.4em]{\large 全球市场主线：金融条件收紧与结构性分化并行，AI资本开支周期扩散，中国政策转向制度优化}}
\author{KC桌面 自动生成}
\date{260609}
\begin{document}
\maketitle
\tableofcontents
\newpage
\section{一页摘要}
\begin{itemize}
  \item 美国金融条件指数（FCI）因就业数据超预期而收紧至俄乌冲突前水平，但收紧主要由股市和汇率驱动，而非央行加息预期，市场自身在主动重新定价风险。
  \item 全球信用市场呈现明显的“质量分化”：欧洲高收益信用利差大幅收窄，而美国高收益走阔，资金流向与价格走势出现错位。
  \item AI资本开支周期正从算力硬件向更广泛的设备领域扩散，传统服务器市场预期大幅上修，PCB和CCL成为新的供应链瓶颈。
  \item 中国政策重心从需求刺激转向制度优化：公积金改革扩围超预期，利率新规利好银行净息差，对外投资收紧重在堵监管漏洞而非资本管制。
  \item 新能源车渗透率创新高至63\%，但终端折扣同步扩大，订单动能回归常态，行业正从规模扩张切换至盈利韧性考验。
\end{itemize}
\section{宏观与利率：金融条件收紧，增长分化加剧}
\textbf{美国金融条件收紧并非央行主导，而是市场自身在重新定价风险；全球增长动能分化，美国强于预期，日本放缓，新兴市场整体稳健。}

\begin{itemize}
  \item 美国FCI上周收紧22.9个基点，其中股票市场贡献13.7个基点，汇率贡献6.4个基点，短期利率贡献几乎为零，表明收紧是市场自发行为。
  \item GS将美联储降息时间点推迟至2027年6月和12月，并预测美国10年期收益率年底升至4.4\%。
  \item 全球CAI仍高于潜在增速，5月为+3.0\%（年化），美国5月CAI为+3.1\%，中国+5.4\%，印度+7.4\%，日本CAI下滑0.8个百分点至+1.0\%。
  \item GS上调印度2026年GDP预测0.3个百分点，但全球整体2026年GDP预测下调0.15个百分点。
  \item 美国、欧元区、英国的工资跟踪指标均从高点回落，就业-工人缺口在主要经济体持续收窄，核心通胀仍在逐步正常化。
  \item GS的FCI脉冲模型显示，金融条件收紧将对未来4个季度的GDP增长产生温和拖累，但幅度可控。
\end{itemize}
\begin{figure}[htbp]
\centering
\includegraphics[width=0.88\linewidth]{figures/fig_055.jpg}
\caption{Exhibit 1 - Exhibit 1: Our US FCI Tightened Back to Its Pre-War Level Following Better Than Expected Employment Data Last Friday}
\end{figure}
\begin{figure}[htbp]
\centering
\includegraphics[width=0.88\linewidth]{figures/fig_056.jpg}
\caption{Exhibit 2 - Exhibit 2: The Global ex Russia FCI Tightened by +3.3bps Last Week Primarily on Equities}
\end{figure}
\begin{figure}[htbp]
\centering
\includegraphics[width=0.88\linewidth]{figures/fig_058.jpg}
\caption{Exhibit 4 - Exhibit 5: Our Global CAI Remains Above Potential}
\end{figure}
\begin{figure}[htbp]
\centering
\includegraphics[width=0.88\linewidth]{figures/fig_059.jpg}
\caption{Exhibit 6 - Exhibit 6: GS Wage Trackers and Inflation Measures}
\end{figure}
\begin{figure}[htbp]
\centering
\includegraphics[width=0.88\linewidth]{figures/fig_060.jpg}
\caption{Exhibit 7 - Exhibit 7: GS Jobs-Workers Gaps \#\# Detailed Indicators Update \#\# Financial Conditions Index (FCI) Exhibit 8: GS Global ex Russia FCI Level (Left) and Weekly Change With Contributions (Right)}
\end{figure}
\begin{figure}[htbp]
\centering
\includegraphics[width=0.88\linewidth]{figures/fig_067.jpg}
\caption{Exhibit 14 - Exhibit 14: FCI Impulses Over the Next 4 Quarters (Left) and in the Euro Area, UK, and US (Right)}
\end{figure}
{\small\color{refgray}\textbf{References:} [R002] GS：市场真正低估的不是IPO供给洪峰，而是企业回购与并购需求的结构性支撑; [R003] GS：金融条件收紧的定价信号远比市场想象得更复杂}

\section{AI/算力：资本开支周期扩散，供应链瓶颈转移}
\textbf{AI资本开支周期正从GPU向更广泛的设备、服务器和存储领域扩散，传统服务器预期大幅上修，PCB和CCL成为新的供给瓶颈。}

\begin{itemize}
  \item GS将2026-2030年传统服务器市场年均预测上调约31\%，预计2030年市场规模达1640亿美元，增速预期修正幅度是AI服务器的两倍。
  \item Dell在传统服务器市场份额从20\%跃升至30\%，单季营收同比增长85\%，AI服务器份额从5\%升至17\%。
  \item Citi调研显示，CCL产能扩张速度系统性落后于PCB，导致CCL供应持续吃紧，涨价已成行业共识且客户基本接受。
  \item NVIDIA将Vera Rubin的SoCAMM2容量削减50\%，主因是DRAM产能瓶颈（供应商仅能满足60\%需求），而非需求疲软，需求将向CMX技术迁移。
  \item 摩根斯坦利指出，SLC NAND、高密度NOR等“旧存储”在AI推理时代面临结构性短缺，数据中心NAND需求2025-2031年CAGR达34\%。
  \item AI设备供应链调研显示，几乎所有相关供应商对2026-2027年订单和收入增长都给出极为乐观的展望，产品升级带来更高单价和利润率。
\end{itemize}
\begin{figure}[htbp]
\centering
\includegraphics[width=0.88\linewidth]{figures/fig_227.jpg}
\caption{Exhibit 2 - Global data center server by vertical and AI/traditional (\textbackslash{}\$, mn) Industry AI server estimate revisions (\textbackslash{}\$, mn)}
\end{figure}
\begin{figure}[htbp]
\centering
\includegraphics[width=0.88\linewidth]{figures/fig_228.jpg}
\caption{Exhibit 3 - Industry traditional server estimate revisions (\textbackslash{}\$, mn)}
\end{figure}
\begin{figure}[htbp]
\centering
\includegraphics[width=0.88\linewidth]{figures/fig_229.jpg}
\caption{Exhibit 4 - Exhibit 4: Nvidia leads the AI server market with 44\% share as of 1Q26, followed by whitebox players (22\%), Dell (17\%), and Super Micro (11\%) AI server revenue market share (\%)}
\end{figure}
\begin{figure}[htbp]
\centering
\includegraphics[width=0.88\linewidth]{figures/fig_232.jpg}
\caption{Exhibit 7 - Exhibit 7: Dell lead the traditional server market in 1Q26 (30\% share), followed by white-box (24\%), HPE (11\%), Lenovo (7\%), and IBM (6\%). Traditional server revenue market share (\%)}
\end{figure}
\begin{figure}[htbp]
\centering
\includegraphics[width=0.88\linewidth]{figures/fig_258.jpg}
\caption{Exhibit 1 - Exhibit 1: Nvidia Vera CPU specs NVIDIA Vera CPU for Agents NVIDIA-Custom Olympus Core 88-Core / 176-Threads}
\end{figure}
\begin{figure}[htbp]
\centering
\includegraphics[width=0.88\linewidth]{figures/fig_267.jpg}
\caption{Exhibit 10 - Exhibit 10: Key directions for SSD products 3 Key Directions for SSD products Super High IOPS GPU Direct Storage High Performance}
\end{figure}
{\small\color{refgray}\textbf{References:} [R007] Citi：NVIDIA削减Rubin内存的真正信号，不是减产，而是存储需求的二次重塑; [R014] GS：服务器市场的真正拐点不是AI，而是传统服务器被重新定价; [R020] 摩根斯坦利：市场正在低估这一轮资本开支周期的宽度和持续性; [R021] 摩根斯坦利：市场低估的不是AI需求，而是“旧存储”的供给再定价; [R030] Citi：AI供应链的瓶颈正在从GPU转移到PCB和铜箔基板}

\section{能源与大宗：供给成本结构性重置，金属需求分化}
\textbf{大宗商品投资者面临的风险已从需求侧周期性放缓转向供给侧持续性成本膨胀，霍尔木兹海峡关闭通过航运、能源和耗材三个渠道施压。}

\begin{itemize}
  \item JPM数据显示，中国铜消费连续第六周趋弱，库存约220kt，比去年同期高40kt；铝去库加速，锌库存攀升至2022年以来同期最高水平。
  \item 霍尔木兹海峡持续关闭推高铁矿石运费，从澳大利亚、巴西、南非到中国的运费自冲突以来均涨超50\%。
  \item JPM系统性降低其对EMEA金属与采矿板块的敞口，仅对挪威海德鲁维持超配评级。
  \item 中国5月出口同比增长19.4\%，进口增长27.5\%，贸易顺差扩大至1054亿美元，但进口增长更多由价格驱动而非数量。
  \item 半导体出口同比暴涨110.9\%，但出口量仅增长2.1\%，主要是价格拉动；原油进口额+15.3\%，进口量却大跌29\%。
  \item GS指出，美国进口韧性并非简单的“解放日效应”消化，而是由不确定性驱动的补库模式正在形成，进口商将关税内化为供应链规划中的常态化变量。
\end{itemize}
\begin{figure}[htbp]
\centering
\includegraphics[width=0.88\linewidth]{figures/fig_208.jpg}
\caption{Figure 1 - Figure 1: Weekly change in China visible copper inventory (SHFE + Bonded). Net inventory movement broadly flat in week ended 5 Jun'26 x-axis = weeks post Chinese New Year (week 0 = week closest to start of CNY), y-axis = kt of cop}
\end{figure}
\begin{figure}[htbp]
\centering
\includegraphics[width=0.88\linewidth]{figures/fig_213.jpg}
\caption{Figure 6 - Figure 6: Aluminum inventories movements in 2026 – weekly change in China visible aluminium inventory (SHFE + Bonded). Aluminum de-stocking stronger at -26kt last week x-axis = weeks post Chinese New Year (week 0 = week closest to}
\end{figure}
\begin{figure}[htbp]
\centering
\includegraphics[width=0.88\linewidth]{figures/fig_216.jpg}
\caption{Figure 9 - Figure 9: Total China visible zinc inventory for week ended 5 Jun'26. Total inventory (264kt) is at the highest seasonal level since 2022 x-axis = weeks post Chinese New Year (week 0 = week closest to start of CNY), y-axis = kt of}
\end{figure}
\begin{figure}[htbp]
\centering
\includegraphics[width=0.88\linewidth]{figures/fig_218.jpg}
\caption{Figure 11 - Figure 11: China steel mill margins extend losses driven by higher coking coal prices}
\end{figure}
\begin{figure}[htbp]
\centering
\includegraphics[width=0.88\linewidth]{figures/fig_159.jpg}
\caption{Exhibit 1 - Exhibit 1: Year-over-year nominal trade growth accelerated in May}
\end{figure}
\begin{figure}[htbp]
\centering
\includegraphics[width=0.88\linewidth]{figures/fig_160.jpg}
\caption{Exhibit 2 - Exhibit 2: China's trade surplus increased to US\textbackslash{}\$105.4bn in May}
\end{figure}
{\small\color{refgray}\textbf{References:} [R006] GS：市场低估的不是短期补库，而是关税博弈正从“一次性冲击”变为“结构性重塑”; [R008] GS：5月中国贸易数据超预期，但真正值得关注的不是增速，而是价格信号; [R013] JPM：金属市场的真正拐点不在需求，而在供给成本的结构性重置}

\section{地产与消费：制度优化与结构分化}
\textbf{中国房地产政策从需求刺激转向制度优化，公积金改革超预期；美国奢侈品消费韧性被低估，但结构分化是真正故事。}

\begin{itemize}
  \item GS指出，公积金改革在三个维度超预期：装修和物业费纳入提取、灵活就业人员纳入缴存（覆盖超2亿人）、全国范围异地互认互贷。
  \item 中国新房销售面积第23周环比下降13\%，但同比仍增长15\%；二手房环比-7\%，同比+31\%，市场情绪保持稳定。
  \item 香港地产股过去两周跑输恒指7\%，主因市场担忧内地资本外流管控和加息预期，但JPM维持全年房价涨幅10\%-15\%的预测。
  \item Citi信用卡数据显示，5月美国奢侈品消费同比增长3\%，连续第五个月正增长，两年叠加增速自2024年以来首次转正至+2\%。
  \item 奢侈品消费呈现“量缩价升”特征：客单价中双位数增长，但交易客户数同比下降约10\%，高端客群支撑增长。
  \item 手表零售大涨7\%，珠宝加速至6\%，皮具放缓至2\%；2026年多数品牌仅小幅提价1-3\%，定价策略趋于谨慎。
\end{itemize}
\begin{figure}[htbp]
\centering
\includegraphics[width=0.88\linewidth]{figures/fig_247.jpg}
\caption{Exhibit 1 - Exhibit 1: Primary GFA sold last week was -13\% wow and +15\% yoy in c.75 cities}
\end{figure}
\begin{figure}[htbp]
\centering
\includegraphics[width=0.88\linewidth]{figures/fig_249.jpg}
\caption{Exhibit 3 - Exhibit 3: Secondary GFA sold last week was -7\% wow and +31\% yoy in c.20 cities Average weekly volume of secondary property sales}
\end{figure}
\begin{figure}[htbp]
\centering
\includegraphics[width=0.88\linewidth]{figures/fig_251.jpg}
\caption{Exhibit 5 - Exhibit 5: Average CSI was +0.4pp wow and +5.9pp yoy Weekly Centraline Salesman Index (CSI) tracker in 4 cities}
\end{figure}
\begin{figure}[htbp]
\centering
\includegraphics[width=0.88\linewidth]{figures/fig_087.jpg}
\caption{Figure 5 - © 2026 Citi Inc. No redistribution without Citi's written permission. Figure 5. Luxury Top Brands, Total Luxury Brands, Total Luxury Market, Beauty/cosmetics and Sport brands/retailers – Monthly absolute credit card spend, 2019–2026 (US\textbackslash{}\$ in millions)}
\end{figure}
\begin{figure}[htbp]
\centering
\includegraphics[width=0.88\linewidth]{figures/fig_089.jpg}
\caption{Figure 7 - © 2026 Citi Inc. No redistribution without Citi's written permission. Figure 7. US credit card monthly spend for Luxury, Sport brands/retailers, and Beauty/cosmetics categories, 2022–2026 (YoY \% chg)}
\end{figure}
\begin{figure}[htbp]
\centering
\includegraphics[width=0.88\linewidth]{figures/fig_090.jpg}
\caption{Figure 8 - © 2026 Citi Inc. No redistribution without Citi's written permission. Figure 8. US credit card monthly spend for Luxury, Sport brands/retailers, and Beauty/cosmetics categories, 2022–2026 (2Y stack \% chg)}
\end{figure}
{\small\color{refgray}\textbf{References:} [R004] Citi：美国奢侈品消费的韧性被低估了，但结构分化才是真正的故事; [R018] GS：公积金改革范围超预期，但市场真正低估的是供给侧的再定价信号; [R027] JPM：香港地产市场真正的考验不是政策阴云，而是数据能否持续验证复苏}

\section{区域市场：亚洲结构性机会与分化}
\textbf{亚洲市场超额收益来自结构性的非共识定价，日本光通信与水冷模块、印度组织化消费企业、韩国正极材料分化是核心线索。}

\begin{itemize}
  \item 摩根斯坦利“三个可操作想法”组合累计获得超过9300个基点的相对超额收益，平均12个月总回报率16.0\%。
  \item 古河电工有望实现新的创纪录利润，光通信产品需求升级和水冷模块在数据中心散热领域的渗透加速是双轮驱动。
  \item 印度消费行业增长驱动力正从普惠式复苏切换为组织化企业的份额掠夺，组织化企业销量增速创四年新高。
  \item 韩国正极材料三巨头战略路径根本分化：Ecopro BM押注欧洲高镍，L\&F率先量产LFP正极，LG Chem出货回升。
  \item NOM指出，中国造船份额的V型反弹揭示的是统计幻觉而非竞争力衰退，中国在绿色船舶上已卡位，Q1 2026年80\%新订单是绿色燃料船。
  \item 日本计划2026年11月修改免税政策，可能影响代购活动，日元贬值使日本奢侈品价格比中国便宜20-25\%。
\end{itemize}
\begin{figure}[htbp]
\centering
\includegraphics[width=0.88\linewidth]{figures/fig_368.jpg}
\caption{Exhibit 2 - Exhibit 2: Cumulative outperformance}
\end{figure}
\begin{figure}[htbp]
\centering
\includegraphics[width=0.88\linewidth]{figures/fig_369.jpg}
\caption{Exhibit 3 - Exhibit 3: Weekly hit ratio}
\end{figure}
\begin{figure}[htbp]
\centering
\includegraphics[width=0.88\linewidth]{figures/fig_268.jpg}
\caption{Figure 5 - The flagship project is Hanwha's planned investment of around USD5bn to expand the Philadelphia shipyard towards twenty vessels a year, and substantive contracts have begun to materialise. In early 2026, Hanwha won a subcontract for concept-design work on the }
\end{figure}
\begin{figure}[htbp]
\centering
\includegraphics[width=0.88\linewidth]{figures/fig_170.jpg}
\caption{EXHIBIT 1 - EXHIBIT 1: European Defense stocks have been weak since Q2'26... European Defense Stocks YTD Performance Since Ukraine War, (base 100 from 1st Jan. 2026)}
\end{figure}
\begin{figure}[htbp]
\centering
\includegraphics[width=0.88\linewidth]{figures/fig_171.jpg}
\caption{EXHIBIT 2 - EXHIBIT 2: ... But valuations are only back to the levels of January 2026 EU Defense - 12m fwd EV/EBITDA}
\end{figure}
{\small\color{refgray}\textbf{References:} [R022] NOM：市场低估了中国造船的结构性护城河，真正的风险不在需求侧; [R032] 摩根斯坦利：亚洲市场真正的超额收益来自结构性的“非共识定价”; [R033] NOM：印度消费复苏并非全面开花，真正的结构性机会在组织化企业的份额掠夺; [R034] NOM：韩国正极材料三巨头的分化，才是理解全球电池供应链重构的真正线索}

\section{信用市场：质量分化而非方向性转折}
\textbf{全球信用市场经历质量分化，欧洲跑赢美国，高收益跑赢投资级，资金流向与价格走势出现错位。}

\begin{itemize}
  \item 美国投资级信用利差走阔1bp，欧洲投资级反而收窄3bp；美国高收益走阔7bp，欧洲高收益收窄13bp。
  \item 美国投资级基金上周仍有36亿美元净流入，年初至今累计690亿，但一级发行470亿美元，年内总供给达1.095万亿（同比+26\%）。
  \item 欧洲投资级基金净流入约7亿美元，ETF为主要流入渠道；欧洲高收益利差大幅收窄13bp，全面领先。
  \item 美国高收益市场中，贷款录得正总回报+0.05\%，跑赢高收益债（-0.4\%），电信和消费品表现垫底，能源和科技相对坚挺。
  \item 摩根斯坦利指出，信用与权益的联动重新变得紧密，权益市场回调直接拖累了信用市场。
\end{itemize}
{\small\color{refgray}\textbf{References:} [R005] 摩根斯坦利：全球信用市场正在经历一场“质量分化”而非方向性转折}

\section{风险偏好与事件驱动：世界杯冲击有限，防务股定价权在预算}
\textbf{2026世界杯对美国经济的真实冲击远小于市场想象；欧洲防务股真正定价权不在乌克兰停火，而在各国预算的不可逆扩张。}

\begin{itemize}
  \item GS估算，世界杯对就业的净影响接近于零：6月非农就业获约4万人额外提振，8月转为约1.5万人拖累。
  \item 世界杯对GDP的拉动有限：Q2年化增速提高0.1个百分点，Q3提高0.05个百分点，Q4小幅回落。
  \item 主办城市酒店价格在比赛日比平时高出31\%-160\%，预计6月全国酒店CPI上涨约1\%。
  \item Bernstein指出，欧洲防务股当前12.9倍EV/EBITDA并不便宜，但如果预算扩张路径成立，这个估值水平提供了未来2-3年的安全边际。
  \item 德国将国防支出目标设定为GDP的3.7\%，法国、意大利、波兰等国的预算增长路径已写入中长期财政框架，停火反而可能强化预算扩张的合法性。
  \item 美国2027财年预算提案1.45万亿美元创纪录，但国会可能不会全批，BAE Systems（45\%收入来自美国）和Leonardo（25\%）受益最大。
\end{itemize}
\begin{figure}[htbp]
\centering
\includegraphics[width=0.88\linewidth]{figures/fig_163.jpg}
\caption{Exhibit 2 - Exhibit 2: Payroll Employment Rose More Than 80k Above Trend in Host Cities During the 1994 World Cup, With Job Gains Largely Concentrated in the Leisure \& Hospitality, Trade, and Business Services Sectors Deviation of Payrolls f}
\end{figure}
\begin{figure}[htbp]
\centering
\includegraphics[width=0.88\linewidth]{figures/fig_165.jpg}
\caption{Exhibit 4 - Exhibit 4: We Estimate That Spending by Foreign and Domestic World Cup Fans Should Provide a Peak Boost to Retail Sales Growth of 0.3pp in June and to Consumer Spending of 0.05pp in Q2 Real Consumption Expenditure by Foreign Trav}
\end{figure}
\begin{figure}[htbp]
\centering
\includegraphics[width=0.88\linewidth]{figures/fig_167.jpg}
\caption{Exhibit 6 - Exhibit 6: Sharply Higher Hotel Prices in Host Cities Around Match Days Imply a World Cup Boost to Average US Hotel Prices of About 1\% in June Hotel Room Price Difference from Regular Prices on Match Nights in Host Cities in June}
\end{figure}
\begin{figure}[htbp]
\centering
\includegraphics[width=0.88\linewidth]{figures/fig_170.jpg}
\caption{EXHIBIT 1 - EXHIBIT 1: European Defense stocks have been weak since Q2'26... European Defense Stocks YTD Performance Since Ukraine War, (base 100 from 1st Jan. 2026)}
\end{figure}
\begin{figure}[htbp]
\centering
\includegraphics[width=0.88\linewidth]{figures/fig_171.jpg}
\caption{EXHIBIT 2 - EXHIBIT 2: ... But valuations are only back to the levels of January 2026 EU Defense - 12m fwd EV/EBITDA}
\end{figure}
\begin{figure}[htbp]
\centering
\includegraphics[width=0.88\linewidth]{figures/fig_176.jpg}
\caption{EXHIBIT 7 - EXHIBIT 7: Germany is the largest Defense market in Europe ex. Russia}
\end{figure}
{\small\color{refgray}\textbf{References:} [R009] GS：2026世界杯对美国经济的真实冲击远小于市场想象; [R010] Bernstein：欧洲防务股的真正定价权不在乌克兰停火，而在各国预算的不可逆扩张}

\section{中国政策与金融：制度重塑进行时}
\textbf{中国政策重心从短期刺激转向制度优化，对外投资收紧重在堵漏洞而非资本管制，利率新规利好银行净息差，人民币定价模型释放复杂信号。}

\begin{itemize}
  \item Citi指出，本轮对外投资收紧的核心不是堵资本外流，而是修补监管漏洞、强化国内政策有效性，宏观影响有限。
  \item 摩根斯坦利认为，利率新规对银行净息差是明确利好，同时作用于资产端收益率和负债端成本，堵死了高息揽储的制度漏洞。
  \item NOM的人民币中间价预测模型显示，模型误差从2025年初接近-1800个基点收窄至约600个基点，定价效率结构性提升。
  \item 最新模型预测（不含逆周期因子）为6.7921，较前值大幅调低236个基点，欧元和韩元是主要贡献因素。
  \item GS指出，5月中国贸易数据超预期，但进口增长驱动力正从量转向价，出口韧性背后结构性分化比总量重要。
  \item 下半年重要事件窗口：7月底政治局会议、11月深圳APEC峰会、12月中央经济工作会议、年底中美高层会晤。
\end{itemize}
\begin{figure}[htbp]
\centering
\includegraphics[width=0.88\linewidth]{figures/fig_001.jpg}
\caption{Figure 1 - Capital controls undermine RMB internationalization? — Curbs on illicit cross-border flows need not undermine the "golden window" for RMB internationalization, in our view. Flows associated with RMB internationalization run through official channels, including}
\end{figure}
\begin{figure}[htbp]
\centering
\includegraphics[width=0.88\linewidth]{figures/fig_004.jpg}
\caption{Figure 2 - © 2026 Citi Inc. No redistribution without Citi's written permission. Figure 2. The capital outflow pressure could be smaller compared with the episode in 2015}
\end{figure}
\begin{figure}[htbp]
\centering
\includegraphics[width=0.88\linewidth]{figures/fig_008.jpg}
\caption{Figure 8 - © 2026 Citi Inc. No redistribution without Citi's written permission. Figure 8. The expansion of official channels is one thing to watch, with Stock Connect showing need to diversify}
\end{figure}
\begin{figure}[htbp]
\centering
\includegraphics[width=0.88\linewidth]{figures/fig_159.jpg}
\caption{Exhibit 1 - Exhibit 1: Year-over-year nominal trade growth accelerated in May}
\end{figure}
\begin{figure}[htbp]
\centering
\includegraphics[width=0.88\linewidth]{figures/fig_160.jpg}
\caption{Exhibit 2 - Exhibit 2: China's trade surplus increased to US\textbackslash{}\$105.4bn in May}
\end{figure}
{\small\color{refgray}\textbf{References:} [R001] Citi：市场低估了本轮对外投资收紧的结构性含义，而非资本外流风险; [R008] GS：5月中国贸易数据超预期，但真正值得关注的不是增速，而是价格信号; [R016] 摩根斯坦利：市场低估了这份利率新规对银行净息差的修复力度; [R024] NOM：人民币中间价定价模型正在释放比市场预期更复杂的信号; [R025] NOM：人民币中间价模型正在释放一个被市场忽略的定价信号}

\section{新能源车：渗透率新高下的盈利考验}
\textbf{新能源车市场正从规模扩张竞赛切换进入盈利韧性考验，渗透率创新高但折扣同步扩大，订单动能回归常态。}

\begin{itemize}
  \item 5月中国新能源车零售渗透率达63\%，批发渗透率61.1\%，但主要品牌订单环比下降15\%，新车上市脉冲效应消退。
  \item 新能源经销商平均折扣7.79\%，比上周继续扩大；燃油车折扣19.54\%，仍在加码。
  \item 碳酸锂价格跌至16.85万元/吨，周降6.1\%，电池成本继续友好。
  \item 品牌分化明显：蔚来YTD订单同比+85\%最猛，问界+29\%，特斯拉+3\%，小米、理想、小鹏YTD同比仍在负区间。
  \item 重点新车日历：6/11蔚来乐道L60改款，6/17比亚迪大唐上市，6月底理想L8改款。
  \item GS指出，市场真正低估的不是销量，而是定价权的转移，终端折扣扩大正在侵蚀盈利空间。
\end{itemize}
\begin{figure}[htbp]
\centering
\includegraphics[width=0.88\linewidth]{figures/fig_245.jpg}
\caption{Exhibit 2 - Exhibit 2: NEV dealer discount trend}
\end{figure}
\begin{figure}[htbp]
\centering
\includegraphics[width=0.88\linewidth]{figures/fig_246.jpg}
\caption{Exhibit 3 - Exhibit 3: ICE dealer discount trend}
\end{figure}
{\small\color{refgray}\textbf{References:} [R017] GS：新能源车市场真正的压力测试，不是销量，而是定价权的转移}

\section{工业与制造业：资本开支周期超预期扩散}
\textbf{AI驱动的资本开支周期正从算力硬件向更广泛的设备领域传导，工程机械景气周期比市场想象的更强。}

\begin{itemize}
  \item 5月中国挖掘机总销量同比增长36\%，国内增长39\%，出口增长34\%，均显著超出市场预期。
  \item 国内需求韧性来自地方政府专项债提前发行和电动化替换需求，而非地产复苏。
  \item 三一重工预计二季度营收同比增长约20\%，中联重科认为外汇拖累减弱后海外增长会加速。
  \item 摩根斯坦利调研显示，AI设备供应链中几乎所有供应商对2026-2027年订单和收入增长都给出极为乐观的展望。
  \item 大族激光用于AI服务器的超快激光钻孔设备单价500-700万，GKG精密AI服务器焊膏设备毛利率高达65\%。
  \item 苹果供应链2026-27年创新周期支撑设备需求，苹果相关项目毛利率通常在40\%以上。
\end{itemize}
\begin{figure}[htbp]
\centering
\includegraphics[width=0.88\linewidth]{figures/fig_257.jpg}
\caption{Exhibit 1 - Exhibit 1: China excavator sales volume Exhibit 2: China's monthly excavator sales volume vs. Sany Heavy's share price}
\end{figure}
{\small\color{refgray}\textbf{References:} [R019] 摩根斯坦利：工程机械的景气周期比市场想象的更强; [R020] 摩根斯坦利：市场正在低估这一轮资本开支周期的宽度和持续性}

\section{科技与软件：AI分选机效应加速分化}
\textbf{AI不是潮水而是一台分选机，正在加速重塑软件和支付行业的商业模式与定价权，赢家与输家鸿沟扩大。}

\begin{itemize}
  \item Bernstein指出，AI推动企业软件从记录系统进化为行动系统，货币化能力成为新的分水岭。
  \item 约60\%的企业客户计划在2028年前用AI服务替代传统人力交付，按人头收费模式正在崩塌。
  \item 欧洲软件股中，SAP、Dassault Systèmes等具备AI经济学捕获能力的平台型企业处于正向循环。
  \item 全球支付市场价值增长正从量的扩张转向质的变现，交易笔数增速（9\%）高于金额增速（8\%）。
  \item 全球仍有约11万亿美元现金/支票待转化，欧洲、亚太、拉美是主要机会区域。
  \item 代币化、增值服务渗透、新支付流是未来收入增长的关键，即使卡交易额增速放缓，结合VAS仍可维持低双位数增长。
\end{itemize}
\begin{figure}[htbp]
\centering
\includegraphics[width=0.88\linewidth]{figures/fig_321.jpg}
\caption{EXHIBIT 1 - EXHIBIT 1: We expect card penetration to approach \textbackslash{}\textasciitilde{}71\% by 2030, up from 64\% in 2025 Card Penetration of Purchase PCE (Nominal and Adj. ex-China)}
\end{figure}
\begin{figure}[htbp]
\centering
\includegraphics[width=0.88\linewidth]{figures/fig_323.jpg}
\caption{EXHIBIT 3 - EXHIBIT 3: On a cc basis, card volumes grew at \textbackslash{}\textasciitilde{}11\% between 2019-22, \textbackslash{}\textasciitilde{}9\% between 2022-25 and we forecast \textbackslash{}\textasciitilde{}8\% growth between 2025-30. Global Card Purchase Volume Growth Drivers (ex. China)}
\end{figure}
\begin{figure}[htbp]
\centering
\includegraphics[width=0.88\linewidth]{figures/fig_330.jpg}
\caption{Exhibit 10 - EXHIBIT 10: Although it sure feels like cards are everywhere, there is still a surprising amount of cash/check in existence globally Global Cash/Other volumes by region (2025E, Total \textbackslash{}\$11T)}
\end{figure}
\begin{figure}[htbp]
\centering
\includegraphics[width=0.88\linewidth]{figures/fig_357.jpg}
\caption{Exhibit 36 - EXHIBIT 36: Value-added services and new flows, which are \textbackslash{}\textasciitilde{}40/\textbackslash{}\textasciitilde{}50\% of V/MA's revenues and growing 1.5-2x the rate of core growth, will be incremental to growth. VAS and new flows can sustain MA's revenue growth in the LDD range ev}
\end{figure}
\begin{figure}[htbp]
\centering
\includegraphics[width=0.88\linewidth]{figures/fig_358.jpg}
\caption{EXHIBIT 37 - EXHIBIT 37: Value-added services make up \textbackslash{}\textasciitilde{}40\% of Mastercard's and \textbackslash{}\textasciitilde{}30\% of Visa's net revenue Services share of net revenue for V/MA (\%)}
\end{figure}
{\small\color{refgray}\textbf{References:} [R011] Bernstein：AI浪潮不会托起所有欧洲软件股，分化已从预期走向现实; [R012] Bernstein：广告市场真正被低估的不是总量增长，而是结构性再分配中的定价权迁移; [R031] Bernstein：市场低估了支付网络的“量价分离”能力}

\section{医疗与创新药：化疗增敏价值重估}
\textbf{中国创新药竞争逻辑正从同类最佳转向化疗增敏，ASCO数据揭示化疗联合疗法的风险收益比重构。}

\begin{itemize}
  \item 摩根斯坦利重点分析AK112联合脂质体伊立替康在2L SCLC中的数据：ORR 61.7\%，mPFS 8.1个月，6个月PFS率69.1\%。
  \item 对比同类方案（BNT327联合紫杉醇）ORR仅37.2\%，mPFS 5.4个月，但AK112组难治患者比例更低（36.7\% vs 50\%），基线差异影响横向对比。
  \item 脂质体伊立替康单药历史数据ORR 44.1\%、mPFS 4.0个月，化疗本身建立了疗效底线，AK112的关键加分项是持久性。
  \item Mesutoclax在MDS/AML双线开花：HR-MDS初治队列ORR 100\%，compCR 90\%；AML初治队列cCR 81.8\%，MRD阴性率86.5\%。
  \item Bernstein指出，Medicare Advantage利润率已在2023-2025年触底，预计2026年回升约1个百分点，此后每年扩张约0.5个百分点。
  \item MA市场将趋于公用事业化，最终由4-5家规模化玩家主导，稳态利润率维持在2-3\%之间。
\end{itemize}
{\small\color{refgray}\textbf{References:} [R028] Bernstein：Medicare Advantage的利润率正在触底，但市场低估了“公用事业化”的长期定价逻辑; [R029] 摩根斯坦利：ASCO数据揭示中国创新药进入“化疗增敏”价值重估期}

\section{结语}
综合来看，当前市场主线是金融条件收紧与结构性分化并行，AI资本开支周期扩散至更广泛领域，中国政策转向制度优化。投资者需关注供给侧的再定价信号，而非仅盯着需求端波动。
\newpage
\section{报告覆盖清单}
\begin{itemize}
  \item [R001] 中国政策与金融：制度重塑进行时 - Citi：市场低估了本轮对外投资收紧的结构性含义，而非资本外流风险
  \item [R002] 宏观与利率：金融条件收紧，增长分化加剧 - GS：市场真正低估的不是IPO供给洪峰，而是企业回购与并购需求的结构性支撑
  \item [R003] 宏观与利率：金融条件收紧，增长分化加剧 - GS：金融条件收紧的定价信号远比市场想象得更复杂
  \item [R004] 地产与消费：制度优化与结构分化 - Citi：美国奢侈品消费的韧性被低估了，但结构分化才是真正的故事
  \item [R005] 信用市场：质量分化而非方向性转折 - 摩根斯坦利：全球信用市场正在经历一场“质量分化”而非方向性转折
  \item [R006] 能源与大宗：供给成本结构性重置，金属需求分化 - GS：市场低估的不是短期补库，而是关税博弈正从“一次性冲击”变为“结构性重塑”
  \item [R007] AI/算力：资本开支周期扩散，供应链瓶颈转移 - Citi：NVIDIA削减Rubin内存的真正信号，不是减产，而是存储需求的二次重塑
  \item [R008] 能源与大宗：供给成本结构性重置，金属需求分化 / 中国政策与金融：制度重塑进行时 - GS：5月中国贸易数据超预期，但真正值得关注的不是增速，而是价格信号
  \item [R009] 风险偏好与事件驱动：世界杯冲击有限，防务股定价权在预算 - GS：2026世界杯对美国经济的真实冲击远小于市场想象
  \item [R010] 风险偏好与事件驱动：世界杯冲击有限，防务股定价权在预算 - Bernstein：欧洲防务股的真正定价权不在乌克兰停火，而在各国预算的不可逆扩张
  \item [R011] 科技与软件：AI分选机效应加速分化 - Bernstein：AI浪潮不会托起所有欧洲软件股，分化已从预期走向现实
  \item [R012] 科技与软件：AI分选机效应加速分化 - Bernstein：广告市场真正被低估的不是总量增长，而是结构性再分配中的定价权迁移
  \item [R013] 能源与大宗：供给成本结构性重置，金属需求分化 - JPM：金属市场的真正拐点不在需求，而在供给成本的结构性重置
  \item [R014] AI/算力：资本开支周期扩散，供应链瓶颈转移 - GS：服务器市场的真正拐点不是AI，而是传统服务器被重新定价
  \item [R015] 未被主题摘要引用 - 摩根斯坦利：面板价格周期的高点已至，市场定价的乐观情绪正在透支未来
  \item [R016] 中国政策与金融：制度重塑进行时 - 摩根斯坦利：市场低估了这份利率新规对银行净息差的修复力度
  \item [R017] 新能源车：渗透率新高下的盈利考验 - GS：新能源车市场真正的压力测试，不是销量，而是定价权的转移
  \item [R018] 地产与消费：制度优化与结构分化 - GS：公积金改革范围超预期，但市场真正低估的是供给侧的再定价信号
  \item [R019] 工业与制造业：资本开支周期超预期扩散 - 摩根斯坦利：工程机械的景气周期比市场想象的更强
  \item [R020] AI/算力：资本开支周期扩散，供应链瓶颈转移 / 工业与制造业：资本开支周期超预期扩散 - 摩根斯坦利：市场正在低估这一轮资本开支周期的宽度和持续性
  \item [R021] AI/算力：资本开支周期扩散，供应链瓶颈转移 - 摩根斯坦利：市场低估的不是AI需求，而是“旧存储”的供给再定价
  \item [R022] 区域市场：亚洲结构性机会与分化 - NOM：市场低估了中国造船的结构性护城河，真正的风险不在需求侧
  \item [R023] 未被主题摘要引用 - NOM：工业AI的货币化，市场低估了2-3年过渡期的结构性摩擦
  \item [R024] 中国政策与金融：制度重塑进行时 - NOM：人民币中间价定价模型正在释放比市场预期更复杂的信号
  \item [R025] 中国政策与金融：制度重塑进行时 - NOM：人民币中间价模型正在释放一个被市场忽略的定价信号
  \item [R026] 未被主题摘要引用 - GS：美国供应链的“新常态”比市场预期的更早到来
  \item [R027] 地产与消费：制度优化与结构分化 - JPM：香港地产市场真正的考验不是政策阴云，而是数据能否持续验证复苏
  \item [R028] 医疗与创新药：化疗增敏价值重估 - Bernstein：Medicare Advantage的利润率正在触底，但市场低估了“公用事业化”的长期定价逻辑
  \item [R029] 医疗与创新药：化疗增敏价值重估 - 摩根斯坦利：ASCO数据揭示中国创新药进入“化疗增敏”价值重估期
  \item [R030] AI/算力：资本开支周期扩散，供应链瓶颈转移 - Citi：AI供应链的瓶颈正在从GPU转移到PCB和铜箔基板
  \item [R031] 科技与软件：AI分选机效应加速分化 - Bernstein：市场低估了支付网络的“量价分离”能力
  \item [R032] 区域市场：亚洲结构性机会与分化 - 摩根斯坦利：亚洲市场真正的超额收益来自结构性的“非共识定价”
  \item [R033] 区域市场：亚洲结构性机会与分化 - NOM：印度消费复苏并非全面开花，真正的结构性机会在组织化企业的份额掠夺
  \item [R034] 区域市场：亚洲结构性机会与分化 - NOM：韩国正极材料三巨头的分化，才是理解全球电池供应链重构的真正线索
\end{itemize}
\section{逐篇报告摘录}
\subsection{[R001] Citi：市场低估了本轮对外投资收紧的结构性含义，而非资本外流风险}
{\small\color{refgray}归类：中国政策与金融：制度重塑进行时}

Citi：市场低估了本轮对外投资收紧的结构性含义，而非资本外流风险

过去一个月，中国多部委密集出台了一系列针对对外投资的收紧措施。从国务院发布新的境外投资条例，到证监会点名三家券商并处罚其境内业务，再到香港金管局与证监会同步发布开户指引，市场的第一反应往往是“资本管制又回来了”。但Citi近日发布的一份研报给出了一个完全不同的判断：本轮收紧的核心不是堵资本外流，而是修补监管漏洞、强化国内政策有效性，其宏观影响有限，但对行业格局和资产定价的结构性含义，市场可能尚未充分定价。

这份报告的价值在于，它没有停留在“政策收紧”的表层叙事，而是区分了两种完全不同的政策逻辑——ODI（对外直接投资）层面是国家安全考量上升，组合投资层面则是执法行动而非新规出台。这两种逻辑的宏观影响不同，对企业和投资者的含义也截然不同。

以下是我们从这份研报中提炼出的五个核心洞察，以及一个报告尚未完全回答的关键问题。

1. 这一轮收紧与2015年的本质区别在于汇率背景和资本外流压力完全不可同日而语

市场对资本管制的记忆，很大程度上仍停留在 年。当时人民币贬值预期强烈，外汇储备大幅下降，资本外流压力巨大。但Citi研报明确指出，当前宏观环境与此截然不同。

报告提供了两个关键数据支撑：第一，目前人民币处于升值周期，外汇结售汇数据呈现净流入态势；第二，从国际收支平衡表测算的“无法解释的资本外流”规模，远低于2015年时期。Citi的计算显示，当前资本外流的上限是可控的——更重要的是，中东资金正在流入香港市场，这在一定程度上可以对冲收紧带来的负面冲击。

这意味着，本轮收紧的政策出发点并非应对汇率危机，而是主动的结构性调整。对于投资者而言，这是一个重要的信号：不要用2015年的框架去理解2026年的政策。如果市场仍然按照“资本管制”的旧剧本定价，就可能错失真正的结构性机会。

2. ODI收紧真正考验的是企业能否将海外扩张转化为国家安全合规能力

Citi研报指出，ODI领域的收紧核心是国家安全考量上升，尤其是在AI竞赛的背景下。新的监管框架将技术、数据控制整合到一个统一体系，并将监管范围扩大至个人。这听起来像是增加了企业的合规成本，但Citi的判断是，ODI本身的增长趋势不会因此改变。

报告中的图表显示，中国非金融类ODI从2016年的约1100亿美元上升至2025年的约1400亿美元，且仍处于结构性上升通道。中国企业的海外扩张，本质上是寻找增长引擎的必然选择。政策收紧不会逆转这一趋势，但会改变竞争格局。

这里的关键洞察是：合规能力本身正在成为企业海外扩张的核心竞争力。那些能够率先建立国家安全合规体系的企业，将获得政策层面的先发优势，而依赖“灰色通道”的企业将面临出清。对于产业决策者而言，这意味着海外投资的战略规划需要前置合规成本，将其从“费用项”重新定义为“竞争壁垒”。

3. 券商罚单不是新规，而是迟到多年的执法——真正的冲击要看执行范围是否扩大

Citi研报特别强调，本轮对券商组合投资业务的收紧并非新监管出台，而是执法行动。被点名的三家券商早在2016年就被证监会“标记”，2022年已有过一轮收紧。这意味着，对于合规经营的主体而言，政策风险是可控的。

但报告也提出了一个重要的“二阶问题”：执行范围是否会扩大？Citi指出，被点名的三家券商持有约2500亿港元的内地客户资产。如果执法行动仅限于这三家，对市场的影响是有限的、可量化的。但如果后续扩大至其他券商、保险、理财产品、私募股权甚至香港物业，影响将显著放大。

这才是市场真正需要关注的风险点。Citi研报的判断是“扩大不会令我们意外”，但并未给出具体的概率判断。这里需要读者结合自身的投资敞口和风险偏好做出判断——尤其是那些通过非官方渠道配置香港资产的主体，需要重新审视合规风

[... middle omitted ...]

高净值个人而言，海外资产的税务合规风险正在上升。Citi的研报没有展开讨论具体的税务影响，但这恰恰是读者最需要深入理解的部分：哪些资金回流会触发税务申报？历史未申报的海外所得是否面临追溯风险？这些问题的答案将直接影响资产配置决策。

5. 人民币国际化的“黄金窗口”并未被关闭，但官方渠道的扩容节奏决定了下半场的胜负

市场另一个常见的担忧是：资本管制会削弱人民币国际化的进程。Citi研报对此给出了明确的否定判断：堵住非法跨境流动，不会影响人民币国际化的官方渠道。QFII、债券通、CIBM Direct、股票通等官方渠道不受影响。香港作为连接器的角色，在人民币升值和持续的中东资金流入支撑下，也不会被实质性削弱。

更值得注意的是，Citi研报提出了一个“后门关闭、前门敞开”的政策逻辑。报告指出，官方对外投资渠道的扩容大概率会跟进，尽管节奏可能较为渐进。可能的选项包括：提高QDII额度（目前为1760亿美元）、扩大理财通试点范围（从大湾区推广至全国）、以及启动保险通试点。

\subsection{[R002] GS：市场真正低估的不是IPO供给洪峰，而是企业回购与并购需求的结构性支撑}
{\small\color{refgray}归类：宏观与利率：金融条件收紧，增长分化加剧}

GS：市场真正低估的不是IPO供给洪峰，而是企业回购与并购需求的结构性支撑

上周全球股市下跌超过2\%，美国市场领跌，动量因子跑输大盘。表面看，这是美国就业数据超预期后利率路径重新定价引发的调整。但GS在最新一期《Global Weekly Kickstart》中，围绕一个看似矛盾的命题给出了一个值得深思的判断：2026年美股IPO规模将创历史纪录，但企业股权需求仍将超过供给。这不是一个简单的“利好出尽”或“供给冲击”的故事。真正值得关注的，不是IPO数量本身，而是支撑需求端的结构性力量——回购、并购与资金流入——正在如何重塑市场的定价逻辑。

这份报告的洞察力在于，它没有停留在“IPO多了会抽血”的直觉层面，而是将股权供需拆解为三个维度：历史规模的相对值、企业行为的结构性变化、以及投资者情绪的真实状态。当多数人盯着IPO的绝对金额时，GS提醒我们：规模本身不是风险，供需失衡的方向才是。而当前，需求端的力量远比市场预期的更持久、更系统。

1. 2026年IPO规模创纪录，但放在市值中只相当于历史平均水平的2/3

GS美国团队预计，2026年将成为美股历史上美元计价的股权发行最高年份。但这份报告同时给出了三个被市场忽视的对比数据。

第一，IPO数量本身并不夸张。今年预计的IPO数量大约在100家左右，仅与历史平均水平持平。这意味着单笔交易的规模更大，而不是市场在无差别地接纳大量小盘股。第二，也是更关键的一点：当发行规模与美股总市值相比时，2026年的发行额仅相当于市值的1.0\%，而历史平均水平是1.5\%。换句话说，即便创纪录，供给的相对强度仍然低于长期均值。第三，需求端的支撑并非空谈——回购、并购和投资者资金流入构成了一个远比IPO更庞大的资金池。

这三个事实合在一起的含义是：市场对“IPO洪峰”的担忧，本质上是混淆了绝对规模与相对压力。1.0\%的供给率意味着，只要需求端不出现系统性萎缩，市场完全有能力消化。真正需要追问的，不是“IPO会不会太多”，而是“需求端凭什么能够持续”。

GS明确将需求端拆解为三个支柱：企业回购、并购交易以及投资者净流入。这三者的共同特征是——它们都不是短期行为，而是由企业资产负债表和战略决策驱动的结构性力量。

回购方面，美国企业近年来的现金储备和盈利能力维持在高位，股票回购已经成为一种制度化的资本回报手段。即便利率环境发生变化，已公布的回购计划也很难大幅削减。并购方面，产业整合的逻辑在多个行业同时展开，尤其是科技、能源和医疗领域，现金加股票的交易结构本身就会产生对股权的结构性需求。

更重要的是，这三者之间还存在正反馈。回购提升每股收益，支撑股价，进而为并购提供更有利的估值环境；并购活动活跃又会吸引更多资本流入。GS的判断——IPO活动回升本身反映了投资者需求强劲——正是建立在这种正反馈之上。企业只有在确信市场有能力吸收新供给时，才会选择上市。因此，IPO的增多不是需求疲软的信号，而是需求旺盛的结果。

上周市场下跌的直接导火索是美国就业数据走强，推动市场重新定价利率路径。GS经济学家已将美联储最后两次降息的时间点推迟至2027年6月和12月，利率策略师预计美国10年期收益率将在年底达到4.4\%，高于此前预期的4.1\%。

但这里需要区分两个概念：利率路径的“延后”与“逆转”。GS并未改变降息周期的方向判断，只是拉长了时间轴。4.4\%的10年期收益率虽然高于之前的预测，但相较 年的高位仍属正常化区间。更重要的是，利率预期对股权需求的影响并非线性。回购决策更多取决于企业现金流和战略判断，而非短期利率波动；并购交易的驱动因素则是产业逻辑和估值差异，利率只是其中的一个变量。

市场真正需要警惕的，不是利率微调本身，而是利率预期变化是否会在情绪层面引发连锁反应。GS的风险偏好指标目

[... middle omitted ...]

医疗板块相对表现较好的背景——前者受益于大宗商品价格，后者则具备典型的防御属性。

对于投资者而言，这一信号的含义是：短期内不应过度追逐趋势性交易，而应重新审视持仓的防御质量。GS的Bull/Bear市场指标目前处于68\%的百分位，虽然尚未进入极端区间，但已经高于历史中位数。这意味着市场定价中已经包含了相当程度的乐观预期，一旦利率或盈利预期出现边际恶化，回调空间是存在的。

5. 报告尚未完全回答的关键问题：需求端假设是否过于依赖企业行为惯性？

GS的判断逻辑高度依赖于一个隐含假设：企业回购和并购行为具有惯性，不会因利率路径的边际变化而大幅调整。这一假设在过去几年确实成立，但并非没有风险。

第一个风险在于监管。美国对大型科技公司反垄断审查的加强，可能抑制部分行业的并购活动。如果并购交易减少，股权需求的一个核心支柱就会削弱。第二个风险在于企业自身的资本配置优先级变化。如果利率长期维持在高位，部分企业可能优先偿还债务而非回购股票。GS报告中并未对这些情景进行压力测试。

\subsection{[R003] GS：金融条件收紧的定价信号远比市场想象得更复杂}
{\small\color{refgray}归类：宏观与利率：金融条件收紧，增长分化加剧}

市场正在重新定价一个被许多人忽略的信号：金融条件正在收紧，而且速度比大多数人预期的要快。这份GS最新发布的全球经济指标更新，提供了一个关键观察窗口——美国金融条件指数（FCI）在就业数据超预期后，已经收紧至俄乌冲突前的水平。这不是一个简单的“加息预期升温”故事，而是一个关于增长、通胀和资产定价之间微妙平衡的重置。

为什么现在重要？因为市场此前一直沉浸在“软着陆”叙事中，认为金融条件会维持宽松以支撑风险资产。但GS的数据显示，过去一周，全球（除俄罗斯）FCI收紧3.3个基点，美国FCI更大幅收紧22.9个基点，其中股票市场贡献了13.7个基点，汇率贡献了6.4个基点。这意味着，市场自身的力量——而非央行政策——正在成为收紧的主要驱动力。对于投资者而言，理解这一动态，比猜测下一次降息时点更为紧迫。

这份报告提供了一个独特的数据框架：它不仅有FCI的周度变化，还有FCI脉冲（FCI Impulse）对GDP增长的前瞻影响，以及当前活动指标（CAI）对经济现状的实时追踪。但真正有价值的，不是这些数字本身，而是它们组合在一起所揭示的“增长-通胀-金融条件”三角关系正在发生结构性转变。以下是我们基于这份研报提炼出的五个核心洞察。

1. 美国金融条件收紧并非央行主导，而是市场自身在重新定价风险

GS的FCI数据显示，美国FCI上周收紧22.9个基点，其中股票市场贡献了13.7个基点，汇率贡献了6.4个基点，而短期利率的贡献几乎为零。这是一个值得深思的构成：收紧的主要驱动力来自风险资产价格下跌和美元走强，而非市场对美联储加息路径的重新定价。

这意味着什么？意味着市场正在主动“收紧”自己，而不是被动等待央行的信号。这种自发的收紧往往比政策驱动的收紧更具破坏性，因为它直接作用于资产价格和信心渠道。当股票市场下跌导致金融条件收紧，而收紧的金融条件又可能进一步抑制经济活动，从而形成负反馈循环。GS的数据显示，美国FCI水平已经回到2022年初的水平，那正是俄乌冲突爆发前的状态。但当前的经济基本面与当时截然不同——通胀虽在回落，但仍高于目标，劳动力市场依然紧张。这种“借力”背景下的收紧，对资产定价的含义更为复杂。

对投资者的含义：不要单纯将金融条件收紧解读为“加息预期上升”。当前收紧的根源是风险溢价的重估，这意味着防御性资产和低波动策略可能比单纯做空利率更有价值。

2. FCI脉冲显示：美国2026年的增长阻力正在累积，但市场尚未充分定价

GS报告中最具前瞻性的部分可能是FCI脉冲数据。FCI脉冲衡量的是金融条件变化对未来四个季度GDP增长的拖累或提振效应。数据显示，美国2026年的FCI脉冲为正0.40个百分点（以年率计），这意味着金融条件收紧将对2026年的经济增长形成显著拖累。相比之下，2025年的脉冲仅为0.15个百分点，而欧元区2026年的脉冲为负0.10个百分点，日本则为正0.50个百分点。

这些数字的MECE意义在于：美国当前看似稳健的增长（GS5月CAI为+3.1\%），正在为2026年的放缓埋下伏笔。金融条件的收紧不是即时的，它通过信贷渠道、财富效应和汇率渠道滞后传导至实体经济。GS的脉冲模型捕捉到了这一点，但市场定价可能尚未完全反映这一滞后效应。

这里仍需验证的关键问题是：FCI脉冲模型的传导机制是否在当前环境下依然有效？因为疫情后，企业资产负债表相对健康，居民部门杠杆率较低，金融条件收紧对实体经济的冲击可能弱于历史均值。GS报告没有完全展开这一点，但投资者需要密切关注这一假设的可靠性。

3. 全球增长分化加剧：印度上修、中国持平、欧洲疲弱，但背后是同一套逻辑

GS的5月CAI数据揭示了一个清晰的增长分化格局。印度以+7.4\%的CAI增速领跑新兴市场，中国为+5.4\%，而德国则录得-0.9\%的负

[... middle omitted ...]

映了欧洲制造业的持续疲弱和能源转型阵痛。

这些分化背后有一个共同逻辑：金融条件收紧对不同经济体的传导速度不同。美国收紧最快，但其经济韧性最强；欧洲收紧较慢，但结构性脆弱性更高。这意味着，全球投资者需要从“统一做多风险资产”转向“精选国家和板块”的策略。

GS的工资追踪器显示，美国、欧元区和英国的工资增速仍在高位，但通胀指标已明显回落。美国PCE通胀从2023年的5.5\%降至2025年的3.0\%附近，而工资增速仍维持在4.5\%左右。这种“工资涨、通胀降”的组合看似矛盾，实则揭示了劳动力市场与通胀之间的传统关系正在发生结构性变化。

报告中的“Jobs-Workers Gap”图表（Exhibit 7）提供了一个关键视角：劳动力供需缺口正在收窄，但速度慢于市场预期。这意味着，工资增长可能不会像通胀那样快速回落，从而对服务业通胀形成持续支撑。GS的数据显示，美国工资增速与PCE通胀之间的差距正在扩大，这暗示企业利润空间可能被压缩，或者企业将成本转嫁给消费者的能力在减弱。

\subsection{[R004] Citi：美国奢侈品消费的韧性被低估了，但结构分化才是真正的故事}
{\small\color{refgray}归类：地产与消费：制度优化与结构分化}

Citi：美国奢侈品消费的韧性被低估了，但结构分化才是真正的故事

美国消费者还在买奢侈品。5月Citi信用卡数据显示，奢侈品品牌整体消费同比增长3\%，已是连续第五个月正增长。更值得关注的是，两年叠加增速自2024年以来首次转正，从4月的-3\%回升至+2\%。

这个数字本身并不惊人。但放在当前宏观背景下来看——中东冲突持续、通胀粘性未消、消费者信心指数连续三个月走弱——它传递的信号远比表面增速复杂。

市场对奢侈品板块的定价已经反映了相当悲观的预期。年初至今，奢侈品股票跑输欧洲市场约26个百分点，估值回调约20\%，当前约20倍远期市盈率较历史均值折价15-20\%。问题在于，这种折价究竟是反映了基本面实质恶化，还是市场对“消费降级”叙事的过度反应？

Citi这份基于超过1000万美国信用卡持卡人样本的报告，提供了一个关键视角：美国奢侈品消费并非铁板一块，高端客群与大众客群正在走向截然不同的消费路径。而真正决定未来12个月板块走向的，不是需求总量，而是供给侧的定价策略与客群结构能否匹配。

1. 消费总量企稳背后，是高端客群对中低端客群的“替代效应”

5月数据最核心的结构性特征是：客单价（ASP）仍保持中双位数增长，但交易客户数（交易量）同比下降约10\%。这意味着，支撑增长的并非更多人在买，而是更少的人在花更多的钱。

这种“量缩价升”的组合，在过去两年中已经成为美国奢侈品消费的常态。Citi报告指出，2023年至2025年期间，高收入人群的消费韧性持续优于中低收入人群，客单价增长始终跑赢客群数量增长。

这背后是一个正在固化的分层逻辑。疫情期间的大幅提价（部分品牌累计涨幅超过20\%）将大量中产消费者挤出了核心奢侈品购买圈层，尤其是入门级产品线——小皮具、迷你包、帆布包、运动鞋等。这些品类原本是品牌培养年轻客群和扩大渗透率的关键入口，但现在价格门槛已经显著提升。

真正有意思的是，这种“高端化”并非品牌主动选择的结果，而是在通胀、关税和地缘风险共同作用下被逼出来的。当品牌不得不持续提价以对冲成本压力时，它们事实上在被动筛选客群——留下的客群消费力更强、价格敏感度更低，但基数也在持续收缩。

这对行业意味着什么？如果交易量持续萎缩，即使客单价维持增长，总量增长的天花板也会越来越低。Citi报告中的两年叠加增速转正是一个积极信号，但若无法转化为交易量的实质性回升，这一改善可能只是价格效应的短期反映。

2. 品类轮动揭示了一个被忽视的变量：硬奢的定价权正在超越软奢

从品类维度看，5月的轮动方向值得深入追问。手表和珠宝品类表现最为亮眼：手表品牌及零售商同比增长7\%，较4月的-4\%明显改善；珠宝加速至+6\%。相比之下，皮具（+2\%）和成衣（+2\%）的增速明显温和。

两年叠加数据更能说明问题：手表品类两年叠加增速从4月的+6\%跃升至+23\%，珠宝从+3\%升至+16\%。而皮具的两年叠加增速仍为-7\%，成衣为-2\%。

这种分化的核心驱动力，在于硬奢品类（手表、珠宝）的定价逻辑与软奢（皮具、成衣）存在本质差异。硬奢的定价中，原材料成本（贵金属、宝石）占比更高，品牌在提价时拥有更充分的成本传导依据。同时，硬奢的消费场景更偏向资产配置和长期保值，客群对价格敏感度本身就低于软奢。

Citi报告在定价策略分析中提到，2026年至今，软奢品牌普遍仅执行极低个位数提价，而硬奢品牌提价幅度略高，反映了贵金属成本上升和瑞郎走强的影响。报告还特别指出，爱马仕和历峰集团旗下珠宝品牌（卡地亚、梵克雅宝、Buccellati）仍存在多年的追赶式提价空间。

这实际上指向一个更深层的判断：在通胀环境和关税不确定性持续的背景下，硬奢品类的定价权优势可能进一步扩大。软奢品牌面临一个两难困境——不提价则利润承压，提价则可能进一步压缩本已脆弱的交易量。

[... middle omitted ...]

更为明显——在这个样本中可能被稀释了。

同时，报告对全球宏观环境的分析提示了一个潜在风险：如果地缘政治紧张和通胀压力持续，实际收入负增长和储蓄侵蚀将最终传导至消费行为。当前高端客群的财富效应（股市上涨、房价温和增长）能否持续对冲这些压力，仍然是一个开放性问题。

Citi在报告中列出了对美出口依赖度最高的奢侈品牌——Tapestry、Capri、Watches of Switzerland、Birkenstock等——以及依赖度最低的品牌——Moncler、Prada、爱马仕等。这一排序本身就是一份风险地图。对于那些高度依赖美国市场的品牌，如果美国消费韧性出现松动，估值折价的修复逻辑将面临根本性挑战。

Citi报告对定价策略的梳理提供了一个有价值的分析框架。2023年品牌间定价策略高度分化（从低个位数到约10\%不等），2024年趋于收敛（平均3-4\%），2025年则因关税调整而出现区域性分化——美国市场因关税上升而提价幅度更高，日本市场则因汇率波动而出现价格洼地。

\subsection{[R005] 摩根斯坦利：全球信用市场正在经历一场“质量分化”而非方向性转折}
{\small\color{refgray}归类：信用市场：质量分化而非方向性转折}

摩根斯坦利：全球信用市场正在经历一场“质量分化”而非方向性转折

全球信用市场过去一周的表现，表面看是方向一致的波动，但仔细拆解后会发现，真正有意义的信息藏在结构里，而非方向里。

摩根斯坦利最新一期全球信用策略周报提供了这样一个信号：美国投资级信用利差走阔1bp，欧洲投资级反而收窄3bp；美国高收益走阔7bp，欧洲高收益却收窄13bp。同一类资产，两个市场，走势截然相反。这不是简单的风险偏好波动，而是一场由资金流向、供给结构和评级分化共同驱动的“质量分化”。

对于配置全球信用资产的决策者来说，理解这种分化的成因，比判断整体方向更重要。因为方向可能是噪音，结构才是信号。

这份周报的框架覆盖了美国、欧洲、亚洲三大市场的投资级、高收益、贷款和信用衍生品，并附带了摩根斯坦利在年中展望中确立的行业配置建议。但报告最值得关注的，不是哪一类资产涨了或跌了，而是“资金流向与价格走势之间的错位”以及“同一市场内部评级梯度的断裂”。

1. 美国信用市场正在承受“供给冲击”与“权益溢出”的双重压力，但资金仍在流入

美国投资级信用利差走阔1bp，超额收益为-0.1\%。表面看只是微幅调整，但结合以下两个数据，就能看出压力所在。

第一，美国投资级债券上周定价规模高达470亿美元，年初至今累计发行1.095万亿美元，同比增长26\%。这是供给端的历史性高水位。第二，权益市场回调直接拖累了信用市场，这是报告原话“the pullback in equities also weighed on credit”。信用与权益的联动在这一周重新变得紧密，而此前市场一度认为信用可以独立于权益走强。

但值得注意的是，美国投资级基金上周仍有36亿美元的净流入，年初至今累计净流入690亿美元。这意味着资金没有撤退，只是定价在重新谈判。供给在放量，需求仍在，但买方变得挑剔。

对于决策者而言，这意味着美国投资级市场的定价权正在从发行人向投资者倾斜。发行人可以趁窗口大量融资，但必须付出更高的利差成本。摩根斯坦利维持对科技/超大规模企业的低配，以及对银行和公用事业的高配，恰好呼应了这一判断——供给冲击之下，防御性和现金流稳定性重新成为定价核心。

2. 欧洲信用市场走出了独立行情，但驱动因素并非基本面改善，而是资金结构的变化

欧洲投资级信用利差收窄3bp，欧洲高收益收窄13bp，明显优于美国。但报告揭示了一个关键细节：资金流入高度集中于ETF，共同基金反而出现小幅净流出。

这是一个容易被忽视的信号。ETF的流入往往反映的是被动配置或战术性仓位调整，而共同基金的流出则代表主动管理人的谨慎。欧洲信用市场的“表面繁荣”可能更多来自流动性驱动而非深度价值认可。

此外，欧洲高收益的利差收窄几乎完全由BB级债券贡献，更低的评级表现则更加分化。这说明欧洲高收益市场的风险偏好是有边界的——资金愿意追逐相对安全的BB级，但对B级和CCC级保持距离。

对于配置欧洲信用的投资者，这份报告传递的信息是：欧洲信用并非没有机会，但机会集中在评级光谱的上端。摩根斯坦利在欧洲推荐高配银行、能源和防御性板块，低配周期性板块尤其是汽车，这一配置逻辑与“资金集中在ETF和BB级”的结构完全吻合。

3. 亚洲信用市场正在经历最剧烈的压缩，但这一压缩的可持续性需要验证

亚洲信用市场在过去一周整体利差收窄2bp，但内部结构分化极其显著。亚洲高收益收窄18bp，而亚洲投资级仅收窄2bp。中国高收益收窄18bp，中国投资级收窄3bp。

从历史分位数看，亚洲投资级当前利差为53bp，处于5年来的最窄水平（0\%分位）。这意味着亚洲投资级几乎已经没有额外的信用溢价。报告将亚洲信用市场描述为“低火、低雇佣”的市场，即交易活跃度低、新增配置意愿低。

亚洲高收益的收窄幅度很大，但报告也指出，亚洲

[... middle omitted ...]

得警惕：隐含波动率上升但信用滞后于权益

报告显示，CDX IG和iTraxx Main的利差分别走阔2bp和1bp，而CDX HY和Xover的利差走阔幅度更大。更关键的是，报告专门指出：“while implied volatility ticked up from the YTD lows, credit lagged the rise in equities vol。”

这句话的含义是：权益市场的隐含波动率已经从年内低点回升，但信用市场的波动率尚未跟上。信用市场通常滞后于权益波动率的变化，如果权益波动率继续上升，信用利差的走阔可能才刚刚开始。

这是一个典型的“先行指标”信号。对于持有信用多头仓位的投资者来说，这可能是需要重新审视风险敞口的时刻。信用衍生品市场正在用沉默发出警告——不是现在，但可能很快。

5. 报告尚未完全回答的关键问题：供给放量之后，需求能否持续消化

这份周报最有趣的地方在于，它提供了大量数据，但刻意回避了“需求能否持续”这个终极问题。

\subsection{[R006] GS：市场低估的不是短期补库，而是关税博弈正从“一次性冲击”变为“结构性重塑”}
{\small\color{refgray}归类：能源与大宗：供给成本结构性重置，金属需求分化}

GS：市场低估的不是短期补库，而是关税博弈正从“一次性冲击”变为“结构性重塑”

过去一周，从中国发往美国的重箱船舶数量环比微增1\%，同比仍为正增长3\%。洛杉矶港的TEU数据预测未来两周还将保持正增长——同比增幅分别达到21.5\%和34\%。这些数字本身并不令人意外，毕竟去年同期正是“解放日”关税生效的基数低点。

但真正值得关注的，不是这些短期数字的高低，而是数字背后的行为逻辑正在发生根本性转变。GS最新一期《美国关税影响追踪》报告揭示了一个关键判断：市场当前看到的进口韧性，并非简单的“解放日效应”被消化，而是一种新的、由不确定性驱动的补库模式正在形成。

这份报告的深层价值，不在于它告诉你下周的船运数据是涨是跌，而在于它提供了一个观察框架，帮助投资者把碎片化的周度数据拼成一幅可决策的图景。对于关注全球贸易、物流运输和制造业布局的决策者来说，理解这个框架，比预测任何单一数据点都更重要。

1. 补库行为正在从“一次性抢跑”转向“持续性的不确定性对冲”

GS的数据显示，在“解放日”周年之后，进口流量并未如部分市场预期的那样大幅回落。相反，洛杉矶港的TEU预测显示，未来两周的同比增幅仍将维持在20\%以上。如果仅仅是基数效应，这个数字应该会快速收窄。

报告对此给出了一个审慎但有力的解释：当前进口的韧性，部分来自“在不确定的地缘政治背景下，以较低的有效关税税率进行补库”。这句话翻译过来就是，进口商不再把关税看作一个需要一次性规避的风险事件，而是将其内化为供应链规划中的常态化变量。

这意味着什么？意味着过去那种“关税宣布-抢运-库存高企-需求断崖”的脉冲式周期，可能正在被一种更平滑、但持续时间更长的补库行为所取代。进口商不再赌关税会取消或降低，而是假设关税将长期存在，并据此调整库存水位和安全边际。

对于航运和物流公司而言，这种变化意味着需求波动率的下降，但同时也意味着“抢运溢价”的消失。对于投资者来说，需要重新评估那些依赖短期运价飙升来获得超额收益的交易逻辑。

2. 第122条关税的150天倒计时，才是影响中期运力的真正变量

GS在报告中明确指出了市场尚未充分定价的风险点：美国行政当局依据《第122条》宣布的新关税将在150天后到期，而政府尚未公布后续方案。这种政策路径的不确定性，正在对中长期的货运规划产生实质性影响。

这里需要区分两个时间维度。短期来看，那些有效税率相对较低的国家，可能会增加对美出口。GS特别指出，这种转移“需要时间才能在数据中体现出来”，尤其是中国制造业和航运在经历了较晚的春节后，产能恢复节奏本身就有滞后。

中期来看，150天后的政策走向才是真正的分水岭。如果关税被延续或扩大，当前的供应链转移将加速固化；如果关税出现松动，那些为了避险而建立的冗余库存和替代产能，将面临痛苦的消化过程。

报告没有给出明确的概率判断，但数据本身已经指向一个方向：不确定性本身正在成为驱动贸易流量的独立变量。对于做全球资产配置的投资者而言，这意味着地缘政治风险的定价因子，需要从“尾部风险”上调为“基准情景”。

报告对近期地缘冲突对运费的潜在影响进行了冷静的拆解。GS认为，当前冲突对集装箱运价的推升作用，需要区分“成本驱动”和“供需驱动”。霍尔木兹海峡并非集装箱航线的主要通道，因此不太可能像红海危机那样导致有效运力的大幅收缩。

真正需要关注的是燃料成本。如果地缘冲突推高燃油价格，航运公司会通过附加费的形式转嫁成本，这将直接推升全球运价，尤其是亚洲至美国航线。但报告同时指出，更高能源成本对全球货运需求的抑制效应，尤其是对美国消费者的影响，可能更为深远。

这是一个典型的“二阶效应”问题。市场容易看到运价上涨对航运公司利润的直接利好，却容易忽略需求端的反噬。如果能源成本持续高企，最终受损的将是整体货运量。对于投

[... middle omitted ...]

业投资的增加（苹果、英伟达、IBM、辉瑞、强生等公司已宣布相关计划）、以及“一个大的美丽法案”中恢复的奖金折旧政策。

这些因素叠加在一起，指向一个判断：即使关税博弈持续，美国国内的货运需求仍有可能在2026年下半年出现实质性改善。报告特别提到，卡车运输板块近期股价表现改善，可能正是市场在提前定价“量”的底部正在形成。

但这里有一个关键的隐含前提：这些利好因素能否兑现，取决于宏观环境是否配合。如果关税导致的经济不确定性持续抑制企业资本开支，或者消费者信心因通胀而走弱，那么“量”的拐点就可能被推迟。报告对此没有给出明确的时间表，但数据框架本身已经为投资者提供了跟踪的锚点。

5. 报告没有完全展开的关键问题：供应链转移的“真实成本”尚未被量化

这是阅读这份报告时最需要保持警惕的地方。GS提到了“中国+1或+2策略”可能带来的结构性贸易机会，也提到了美国制造业回流对国内货运需求的提振。但报告没有深入讨论的是，这种转移本身需要付出多大的成本，以及这些成本最终由谁来承担。

\subsection{[R007] Citi：NVIDIA削减Rubin内存的真正信号，不是减产，而是存储需求的二次重塑}
{\small\color{refgray}归类：AI/算力：资本开支周期扩散，供应链瓶颈转移}

Citi：NVIDIA削减Rubin内存的真正信号，不是减产，而是存储需求的二次重塑

当市场看到“NVIDIA将Vera Rubin的SoCAMM2容量削减近50\%”这条消息时，第一反应通常是“需求不行了”或“供应过剩了”。但Citi这份最新研报给出了一个完全相反的判断框架：削减容量的背后，是DRAM供应瓶颈倒逼的技术路线调整，而真正的需求不仅没有被削弱，反而在向一个新的存储层级——CMX技术——迁移。

这不是一个关于“减少”的故事，而是一个关于“转移”的故事。理解这个转移，才能看懂未来两年AI基础设施的存储竞争格局。

1. 削减SoCAMM2容量不是需求疲软，而是DRAM产能瓶颈的被动选择

Citi报告明确指出，NVIDIA将Vera Rubin NVL72服务器中每颗CPU的SoCAMM2容量从1,536GB（192GB x 8）降至768GB（96GB x 8），降幅达到50\%。但真正驱动这一决策的，不是NVIDIA对AI推理需求的下调，而是全球DRAM产能的硬约束。

报告估算，当前内存供应商只能满足SoCAMM2总需求的60\%。这是一个极为关键的数字：即使NVIDIA不主动削减，供应链也无法支撑原定的192GB模块方案。60\%的供应率意味着，即便需求不变，物理上的产能天花板也在迫使所有AI硬件厂商重新设计内存架构。

这里需要特别注意一个容易被忽略的逻辑：Citi认为SoCAMM2的需求总量不会因为单模块容量下降而减少。为什么？因为AI工作负载对内存容量的需求是刚性的——更大的模型、更长的上下文窗口、更高的KV Cache要求，这些都不会因为NVIDIA换用了更小容量的模块就自动消失。需求只是被“挤压”到了另一个技术路径上。

2. CMX技术才是这份报告真正的核心判断：G3.5内存层将重塑NAND需求曲线

Citi报告中最具前瞻性的判断，是提出Context Memory eXtension（CMX）技术将作为“G3.5内存”被引入，以应对不断增长的KV Cache需求。这个判断的含义远比字面看起来深远。

传统上，AI服务器的内存层级是清晰的：DRAM（HBM/SoCAMM）负责高速缓存，NAND（SSD）负责持久化存储。但KV Cache的爆炸式增长正在打破这个边界。当大模型的上下文窗口从几千token扩展到百万级别，DRAM的容量和成本都无法支撑全量KV Cache的驻留。CMX技术本质上是在DRAM和NAND之间插入一个“中间层”，用NAND来扩展KV Cache的容量。

Citi的结论是：CMX的采用将最终增加NAND需求。这个判断的潜台词是，AI基础设施对NAND的拉动，正在从传统的“存模型权重”转向“存推理状态”。这是一个需求性质的转变——模型权重存储是一次性的、静态的，而KV Cache存储是持续的、动态的、每轮推理都在产生新数据。如果这个判断成立，NAND的长期需求曲线将比市场当前定价的更为陡峭。

3. 市场对“容量削减”的定价可能是错误的：需要重新理解SoCAMM2的真实含义

当前市场对NVIDIA SoCAMM2容量削减的解读，大多停留在“成本优化”或“供应链管理”层面。Citi的分析暗示，这些解读可能低估了结构性变化的深度。

首先，SoCAMM2本身是NVIDIA为Vera Rubin设计的专用DRAM模块，而非通用内存。它的容量削减并不意味着整个AI服务器市场的DRAM需求下降，而是意味着DRAM的分配方式正在从“每服务器高容量”转向“更多服务器+中等容量+CMX补充”。总量逻辑没有变，分配逻辑变了。

其次，Citi强调“即使NVIDIA将SoCAMM2模块从192GB降至96GB，SoCAMM2需求不变”——这句话的隐含意思是，NVIDIA可能会通过增

[... middle omitted ...]

G3.5内存的框架，但有几个关键问题并未展开。

第一，CMX技术是NVIDIA自研，还是需要与NAND供应商联合开发？如果CMX需要新的接口标准和控制器协议，那么NAND供应商的竞争格局将发生分化——谁先与NVIDIA完成技术对接，谁就可能获得先发优势。

第二，CMX对NAND的寿命和性能要求是什么？KV Cache的读写模式与传统存储工作负载完全不同——高频率、小粒度、随机访问。这对NAND的耐久度（Program/Erase cycles）和延迟都是严峻挑战。现有企业级SSD是否能满足，还是需要开发专门的“KV Cache SSD”？这个问题的答案将决定哪些NAND厂商能真正受益。

第三，CMX的性价比是否真的优于继续扩大SoCAMM2容量？Citi认为DRAM产能是瓶颈，但NAND同样面临产能限制，尤其是高密度3D NAND的良率和资本开支。如果CMX所需的NAND产能同样紧张，那么“转移需求”可能只是把瓶颈从DRAM移到了NAND，而不是真正解决了问题。

\subsection{[R008] GS：5月中国贸易数据超预期，但真正值得关注的不是增速，而是价格信号}
{\small\color{refgray}归类：能源与大宗：供给成本结构性重置，金属需求分化 / 中国政策与金融：制度重塑进行时}

GS：5月中国贸易数据超预期，但真正值得关注的不是增速，而是价格信号

5月中国出口同比增长19.4\%，进口同比增长27.5\%，双双超出市场预期的15\%和26\%。贸易顺差从4月的848亿美元跃升至1054亿美元。这份数据表面上看是“量价齐升”的乐观图景，但GS这份报告揭示了一个更微妙的信号：进口增长的驱动力正在从“量”转向“价”，而出口的韧性背后，结构性分化远比总量数字重要。

对于关注中国资产定价和全球供应链布局的决策者而言，5月数据最值得追问的不是“增速能持续多久”，而是“价格信号在告诉我们什么”。GS的团队在报告中给出了几个关键线索，但一些二阶影响并未完全展开——这正是我们需要深入讨论的地方。

1. 出口超预期主要靠美国订单拉动，但低基数效应可能掩盖了真实动能

5月中国对美出口同比增速从4月的11.3\%飙升至35.4\%，这一数字看起来惊人，但GS明确指出，部分原因是去年同期的低基数。从环比看，经季节性调整后的对美出口仅增长1.5\%，远低于4月的4.2\%。这意味着，对美出口的“爆发式增长”更多是统计上的错觉，而非需求端的根本性改善。

更值得关注的是区域分化。对欧盟出口增速从13.4\%放缓至7.6\%，对拉丁美洲出口甚至出现环比下降。与此同时，对东盟出口同比增长24.3\%，对非洲和其他新兴市场保持18\%-22\%的增速。这个格局说明：中国出口的韧性并非来自单一市场的“补库存”，而是新兴市场多元化布局的持续深化。

但这里有一个尚未完全回答的问题：新兴市场的需求是否可持续？GS的报告没有深入分析东盟和非洲国家的进口能力与债务压力，而这些恰恰是判断下半年出口动能的关键变量。

GS在报告中反复强调一个现象：进口金额增速远高于进口数量增速。半导体进口额同比增长68.0\%，但进口量同比下降1.0\%；原油进口额增长15.3\%，但进口量暴跌29.0\%；天然气进口额增长11.0\%，但进口量基本持平；铜矿进口额增长34.0\%，但进口量下降1.4\%。

这不是偶然。这是全球大宗商品价格传导和中国国内需求结构变化的共同结果。进口“量价背离”意味着：中国正在为同样的实物支付更高的价格，而国内的实际需求可能并没有名义数据显示的那样强劲。这对投资者而言，是一个需要警惕的信号——如果价格驱动的进口增长不可持续，那么进口增速的回落可能会比预期更快。

GS没有进一步讨论的是：这种价格压力是否会传导至PPI和CPI，进而影响中国央行的货币政策空间。这是一个需要结合后续通胀数据才能回答的问题。

3. 半导体出口的“价格游戏”暴露了中国科技出口的真实竞争力

5月半导体出口额同比增长110.9\%，但这个数字同样需要拆解。GS的数据显示，半导体出口量仅增长2.1\%，几乎全部增长来自价格效应。这意味着，中国半导体出口的“爆发”并非产能或技术突破带来的量增，而是全球芯片价格上涨的被动受益。

这引出一个更深层次的问题：当全球芯片周期转向下行时，中国半导体出口还能保持这样的增速吗？报告没有给出判断，但我们可以合理推测：如果价格是主要驱动力，那么出口增速的波动性将远高于市场预期，这对于依赖半导体出口增长的地区（如广东、江苏）的经济预测，可能意味着更大的不确定性。

相比之下，铝材出口量的温和增长（同比38.5\%，其中部分也是价格驱动）和家电出口的稳定增长（同比9.2\%）显得更加“真实”。这些传统制造业的出口韧性，或许才是中国出口基本盘的真实写照。

5月1054亿美元的贸易顺差是2025年以来的最高值。从表面看，这是出口强劲、进口相对疲弱的结果。但结合前面的分析，进口疲弱并非需求不足，而是价格高企抑制了实际采购量。如果价格压力持续，进口额可能被动上升，从而压缩顺差空间。

另一个尚未被充分讨论的问题是：顺差的结构性来源是什么？GS的数据显示，对美顺差

[... middle omitted ...]

天然气进口价格的上涨，是否会加速中国对可再生能源的投入？半导体进口价格的上涨，是否会改变中国“自主可控”战略的节奏？这些问题的答案，不仅影响资产定价，也影响全球供应链的重新配置。GS的报告留下了这些伏笔，但没有给出判断。

6. 投资者需要建立的观察框架：从“总量思维”转向“结构思维”

第一层：追踪进口“量价拆分”。不要只看进口总额增速，要看数量增速和价格增速的背离程度。如果价格增速持续高于数量增速，说明国内需求可能被高估，而全球通胀压力正在通过贸易渠道传导至中国。

第二层：关注出口的区域结构变化。对美出口的低基数效应将在下半年消退，届时真正的出口动能将暴露出来。如果对东盟和非洲的出口能够持续增长，说明中国制造的多元化布局正在奏效；如果这些地区的增速也开始放缓，那么出口下行风险将显著上升。

第三层：警惕半导体出口的“价格泡沫”。半导体出口的“量平价增”是不可持续的。当全球芯片周期转向时，中国半导体出口数据可能会出现剧烈波动，这将对相关行业指数和市场情绪产生冲击。

\subsection{[R009] GS：2026世界杯对美国经济的真实冲击远小于市场想象}
{\small\color{refgray}归类：风险偏好与事件驱动：世界杯冲击有限，防务股定价权在预算}

一场在大半个美国土地上持续六周的体育盛事，5到6百万现场观众，11个覆盖美国GDP三分之一的都会区。当这些数字叠加，市场很容易产生一个直觉：2026年世界杯将给美国经济带来显著脉冲。但GS最新发布的这份研报给出了一个需要审慎对待的判断——世界杯的经济影响不仅规模有限，而且几乎会在赛事结束后完全逆转。

这份报告的价值不在于它确认了世界杯会带来增长，而在于它提供了一个对冲市场过度乐观的量化框架。GS基于1994年美国世界杯、过去20年的超级碗以及洛杉矶、亚特兰Daiwa盐湖城奥运会的经验数据，对2026年世界杯的就业、GDP和通胀影响做了跨场景的严谨测算。结论是：市场真正需要关注的不是世界杯带来的短期增量，而是这些增量如何被后续的回落所抵消，以及由此产生的数据噪音对趋势判断的干扰。

GS对就业影响的估计是整份报告中最具操作参考价值的部分。根据历史数据，1994年美国世界杯期间，主办城市的就业比趋势线高出约8万人，而历届夏季奥运会的平均提振幅度约为4万人。这些增量高度集中在休闲酒店、零售贸易和运输行业，而商业服务业的招聘则在赛前数月就已提前启动。

GS将此规模放大约30\%以匹配2026年世界杯更大的比赛场次和观众规模，并考虑已发生的提前招聘，最终估算：6月非农就业将获得约4万人的额外提振，7月降至约1万人，8月赛事结束后转为约1.5万人的拖累，随后数月持续缓慢回归趋势。

这意味着，世界杯对就业的净影响接近于零。市场如果仅看到6月非农数据超预期而解读为经济动能增强，很可能在随后的数据修正中被迫调整预期。对于宏观交易者而言，这份报告提供了一个清晰的信号：6-7月的就业数据需要剔除世界杯效应后再做趋势判断。

GS用两种方法交叉验证了世界杯对GDP的影响。第一种是自下而上的消费支出加总，包括外国游客、国内游客和本地居民的食品饮料消费，并叠加了BEA旅行与旅游卫星账户的1.5倍乘数效应。第二种是自上而下地利用超级碗对州级GDP的影响进行规模缩放。

两种方法得到的结论高度一致：世界杯对Q2年化GDP增长的提振约为0.1个百分点，Q3约为0.05个百分点，随后在Q4转为小幅拖累。这个量级在季度GDP波动中几乎可以忽略不计。

更值得关注的是报告提到的一个下行风险：与超级碗不同，许多世界杯比赛将在工作时间进行，可能降低分心员工的生产率。这一因素未被纳入主估算，但GS明确将其列为风险点。如果这一效应显著，世界杯对GDP的净影响甚至可能为负。这是一个市场尚未充分定价的视角。

酒店价格的暴涨是市场最容易捕捉到的世界杯效应。报告显示，堪萨斯城比赛夜的酒店价格比平时高出160\%，迈阿密和费城也分别达到75\%和73\%。GS据此估算，世界杯将使6月全美平均酒店价格上涨约1\%，随后在7月和8月逐步回落。

但问题的关键在于，酒店价格以外的通胀影响存在巨大不确定性。历史数据显示，大型体育赛事期间，主办城市的餐饮和交通价格确实会跳升，但赛事结束后的回落幅度并不完整。GS估算，世界杯对6月核心CPI的提振仅为0.03个百分点，对核心PCE的提振为0.04个百分点，7月再增加0.01个百分点，随后在8月及之后转为约0.01个百分点的拖累。

这个量级几乎无法被月度通胀数据所辨识。真正的不确定性在于，餐饮和交通价格的上涨是否会在赛事结束后以“更缓慢的价格上调”的方式被部分固化。如果是，那意味着世界杯可能对服务通胀产生一个微小但持久的上行贡献。报告没有给出明确结论，这恰恰是需要后续数据验证的关键变量。

4. 报告的隐含框架：大型体育赛事的经济影响是一个“均值回归”过程

GS这份报告最深刻的洞察不在于具体数字，而在于它所揭示的经济规律：大型体育赛事的经济效应本质上是短期的资源再配置，而非总量的净增加。就业在赛前和赛中上升，在赛后回吐；消费在赛事

[... middle omitted ...]

某些行业和地区的表现，却也可能因此忽视更宏大的结构性趋势。

5. 市场尚未充分定价的视角：世界杯效应可能加剧宏观数据的误判

报告在结尾处点到为止地提出了一个容易被市场忽视的问题：世界杯期间，许多比赛在工作时间进行，可能降低分心员工的生产率。这一因素未被纳入GDP估算，但GS明确将其列为下行风险。

这一风险的潜在影响可能被低估。如果生产率下降导致企业产出受损，而同时就业数据因临时性招聘而表现强劲，那么市场将面临一个“就业强、产出弱”的矛盾信号。这种信号组合在过去往往被解读为经济过热，从而强化加息预期，但这一次背后的驱动力却是体育赛事带来的短期扭曲。

对于利率市场而言，这意味着2026年下半年的宏观数据需要更长的观察窗口才能确认趋势。任何基于单月数据做出的政策预期调整，都可能在不完整的框架下产生偏差。

这份GS研报还提供了许多有价值的细节，包括各主办城市在就业、消费和价格方面的差异化影响，以及不同行业在赛事前后的雇佣节奏差异。但这些内容已经超出了本文的讨论范围。

\subsection{[R010] Bernstein：欧洲防务股的真正定价权不在乌克兰停火，而在各国预算的不可逆扩张}
{\small\color{refgray}归类：风险偏好与事件驱动：世界杯冲击有限，防务股定价权在预算}

Bernstein：欧洲防务股的真正定价权不在乌克兰停火，而在各国预算的不可逆扩张

2026年第二季度以来，欧洲防务股经历了显著回调。自年初高点算起，莱茵金属回撤超过40\%，泰雷兹跌幅超过20\%，莱昂纳多也在15\%左右。市场的主流叙事是：乌克兰停火预期升温、地缘溢价消退、资金正在撤出这个曾经最拥挤的板块。

但这份Bernstein研报提出了一个与市场直觉相反的判断：停火本身并不是欧洲防务股的真正风险，真正决定这些公司长期价值的，是各国军事预算的扩张路径是否具有结构性、不可逆的特征。而恰恰在这个问题上，市场目前的定价存在系统性低估。

报告的核心主张可以概括为一句话：欧洲防务股当前12.9倍的EV/EBITDA并不便宜，但如果预算扩张路径成立，这个估值水平反而提供了未来两到三年的安全边际。关键在于，投资者需要区分“冲突溢价”和“预算周期”——前者会随停火消退，后者不会。

1. 市场对停火的定价逻辑存在根本性误判：停火反而可能强化预算扩张的合法性

这是整份报告最具反直觉价值的判断。Bernstein并不否认，如果出现可持续的停火，欧洲防务股会出现显著估值收缩。但他们同时指出，停火对预算的影响并非单向的负面冲击。

关键逻辑在于：欧洲各国当前的扩军承诺，其政治基础已经超越了“应对乌克兰危机”这一单一事件。德国将国防支出目标设定为GDP的3.7\%，法国、意大利、波兰等国的预算增长路径也已写入中长期财政框架。这些承诺的合法性来源，已经从“应对紧急冲突”转向“重建欧洲自主防务能力”这一结构性命题。

换句话说，即便乌克兰实现停火，欧洲各国也很难在政治上逆转已经启动的装备采购和产能扩张计划。原因很简单：这些计划涉及到的就业、工业基础、供应链安全，已经与国内政治承诺深度绑定。Bernstein的判断是，欧洲的再武装周期将持续到2030年，即便乌克兰实现持久和平。

这个判断如果成立，意味着当前市场对防务股的定价中，隐含了太多的“冲突溢价消退”假设，而忽略了“预算周期延续”这一更持久的驱动力。

2. 德国正在成为欧洲防务市场的绝对核心：21\%的复合增长意味着什么

Bernstein的数据显示，德国已经超越英国，成为欧洲最大的防务市场（占欧洲除俄罗斯外军事预算的17\%），也是全球第四大防务支出国。更重要的是，德国的预算增长路径是所有欧洲大国中最具确定性的。

报告测算，德国国防预算将从2026年的约1080亿欧元增长到2030年的1800亿欧元，复合年化增长率约为12\%。但更关键的指标是“可寻址市场”的增长：由于德国将更高比例的预算投向装备采购（而非人员和基础设施），装备支出的复合年化增长率将达到21\%。

这意味着什么？如果这个增速能够兑现，德国防务装备市场将在五年内翻倍。这对于在该市场有深度布局的企业——尤其是莱茵金属、蒂森克虏伯海事系统等本土冠军——意味着收入增长的可预见性远超市场当前预期。

莱茵金属是其中最突出的标的。Bernstein预计其2025至2030年的收入复合年化增长率将达到29\%，是欧洲防务股平均增速（13\%）的两倍以上。这个增速的支撑来自弹药、陆地车辆和海军装备三个领域，而这三个领域恰恰是德国预算中增长最快的细分项。

3. 美国预算的不确定性被市场过度放大：即使“打折版”也意味着强劲增长

市场对美国防务预算的担忧主要集中在政治层面：1.5万亿美元的2027年预算提案能否通过国会？Bernstein的判断是，不可能原封不动通过。但他们同时指出，即使只通过“基础预算”部分（约1.1万亿美元），也意味着同比22\%的增长。

这个判断的意义在于，市场往往把政治博弈的噪音等同于实际支出的风险。但美国防务预算的历史规律是：基础预算具有极强的韧性，即便在政治分裂时期也极少出现实质性削减。2027年的基础

[... middle omitted ...]

辑和风险特征差异显著。

莱茵金属的核心逻辑是增长速度和市场份额。29\%的收入复合增长预期意味着，即便行业整体增速放缓，莱茵金属仍能通过德国本土市场的深度渗透和弹药产能的持续扩张维持高增长。但风险在于，市场对其短期波动容忍度极低——当前股价较年内高点回调超过40\%，部分反映了市场对其订单节奏和产能爬坡的担忧。

莱昂纳多的核心逻辑是盈利上修空间。Bernstein认为，在覆盖的欧洲防务股中，莱昂纳多拥有最大的盈利超预期潜力。近期CEO更替带来的股价回调，在他们看来是一个入场窗口。但这里需要指出一个尚未充分验证的假设：莱昂纳多能否在保持美国市场增长的同时，提升利润率？报告没有完全展开这一点的风险。

泰雷兹的核心逻辑是结构性改善。Bernstein认为，泰雷兹的网络安全业务目前被市场过度悲观定价，而核心防务业务（尤其是航空电子和太空领域）的强劲表现将在2026年第二季度开始兑现。这个判断的关键在于，网络安全业务的估值拖累是否会持续超出预期？报告对此没有给出确定性的答案。

\subsection{[R011] Bernstein：AI浪潮不会托起所有欧洲软件股，分化已从预期走向现实}
{\small\color{refgray}归类：科技与软件：AI分选机效应加速分化}

Bernstein：AI浪潮不会托起所有欧洲软件股，分化已从预期走向现实

市场正在犯一个典型的“beta思维错误”——把AI当作一个可以拉升整个软件板块的行业性催化剂。但Bernstein最新发布的欧洲软件与IT服务深度研报给出了一个截然不同的判断：AI不是潮水，而是一台分选机。它正在以比多数投资者预期的更快的速度，重塑商业模式、定价权和利润率结构，其结果不是整体水位上升，而是赢家与输家之间的鸿沟加速扩大。

这份报告整合了该机构2025上半年在软件和IT服务领域的研究成果，核心结论简洁但尖锐：AI已经从一个“生产力叙事”演变为“经济模型重置”。那些能够将AI嵌入核心工作流、转化为定价权和留存率的平台型企业，正在进入一个正向循环；而那些依然依赖劳动力套利或浅层功能叠加的商业模式，则面临结构性定价压力。Bernstein的评级表清晰地反映了这一判断——Outperform集中在SAP、Dassault Systèmes、Capgemini、Indra等具备“AI经济学”捕获能力的标的，而低护城河、高自动化替代风险的模型则被标记为谨慎。

这不仅仅是一份评级更新。它提供了一个完整的分析框架，用以理解AI如何重新定义软件和服务的价值分配。以下是我们从这份研报中提炼出的五个核心洞察，以及一个尚未被充分回答的关键问题。

1. 软件的价值正从“记录系统”迁移到“行动系统”，货币化能力成为新的分水岭

Bernstein提出一个关键概念转换：AI正在推动企业软件从传统的“systems of record”（记录系统）进化为“systems of action”（行动系统）。前者负责记录发生了什么，后者则直接驱动决策和行动。这个转变看似抽象，但其经济含义非常具体——价值不再由功能广度决定，而是由AI嵌入工作流后带来的ARPU提升、客户留存率和定价权决定。

报告的核心增量信息是，ERP领域的AI渗透率预计到2030年将接近90\%，但早期Copilot类产品已经出现“商品化”迹象。真正的差异化正在向“自主化工作流”迁移——那些能够端到端自动化财务、供应链和人力资源流程的AI代理，才是未来定价权的来源。这意味着，市场目前对ERP的估值逻辑仍停留在“稳定的SaaS复合增长”上，但Bernstein认为，它已经变成了一个“AI执行故事”。谁能在AI货币化上兑现，谁就能获得估值重估；谁不能，谁就可能面临利润率压缩和增长失速。

这个判断的“所以呢”是：投资者不能再把ERP当作一个同质化的beta板块来配置。SAP、Oracle、Workday之间的估值分化，将越来越多地由AI执行力的差异驱动，而非云迁移进度。

2. IT服务正在经历从“劳动力套利”到“服务即软件”的结构性重置

如果说软件端的变革是价值层的迁移，那么IT服务端的变革则是商业模式的根本性重构。Bernstein的调研数据显示，约60\%的企业客户计划在2028年前用AI替代人工交付模式。这个节奏比市场预期的更快。

这意味着传统的“按时计费”模式正在被“基于成果、AI赋能”的新模式取代。报告将这种新范式定义为“Services-as-Software”（服务即软件）。其核心特征是可规模化、可重复、附带专有IP。那些能够将咨询能力、AI技术和自有IP结合起来的服务商，将获得更大的TAM（总可寻址市场），同时压缩传统收入池。而那些依然依赖劳动密集型交付模式的玩家，则会面临越来越大的定价压力和利润率侵蚀。

Bernstein在这个领域最看好的标的是Capgemini和Indra。前者被定义为“AI transformer”，即在AI转型中具备结构性优势；后者则被列为“最佳创意”（Best Idea），因为AI、国防和航天三大主题的叠加正在加速其估值重估。这里有一个值得注

[... middle omitted ...]

ty）两个维度。不同收入结构的公司落在这个矩阵的不同位置，决定了它们面对AI是受益还是受损。

高防御性的平台（如ERP、数据库、嵌入式工作流）能够将AI作为附加价值层来货币化，保持韧性。而高自动化程度的细分领域（如分析、业务流程外包、应用开发）则面临结构性定价压力。报告明确指出，AI很少是“生死存亡”级别的威胁——它更多是在技术栈内部重新分配价值。但正是这种“再分配”，导致了“AI转型者”和“低护城河模型”之间的差距加速扩大。

这个框架的实用价值在于：它告诉投资者，不要被“AI概念股”的标签所迷惑。一家公司是否受益于AI，取决于它的收入中有多少来自可防御的、嵌入式的AI应用，而不是它是否宣布了AI战略。Bernstein的评分卡显示，大型平台（微软、Oracle、SAP、Workday）受益于统一数据、研发规模和生态优势，而细分领域的专业厂商（如Sage、Visma）则在特定领域拥有可防御的利基。两者都可能被市场奖励，但中间地带的“无差异化通用型”公司将面临压力。

\subsection{[R012] Bernstein：广告市场真正被低估的不是总量增长，而是结构性再分配中的定价权迁移}
{\small\color{refgray}归类：科技与软件：AI分选机效应加速分化}

Bernstein：广告市场真正被低估的不是总量增长，而是结构性再分配中的定价权迁移

全球广告市场正处在一个微妙的拐点。一方面，数字广告依然保持着两位数的增长，2024年全球数字广告支出预计将达到近7000亿美元。另一方面，传统媒体的广告收入增长几乎停滞，甚至在某些季度出现负增长。市场的主流叙事仍然聚焦于“数字化渗透率还有多少空间”这样的总量问题。

但Bernstein这份最新研报提供了一个更具穿透力的结构性判断：广告增长的核心驱动力并非消费总量的扩张，而是注意力的再分配。全球媒体消费时间已经稳定在每天约11小时，这意味着总库存的增长空间极为有限。真正发生的是价值在不同格式、不同平台之间的系统性迁移。而这场迁移的赢家，将不是那些拥有最多用户的平台，而是那些能够将规模优势转化为不可替代的定价权的玩家。

这份报告的价值不在于罗列数字广告的增长数据，而在于它提供了一个理解广告市场底层逻辑的分析框架。它揭示了两个截然不同的广告市场、两阶段增长模型，以及一个正在被市场低估的关键变量——碎片化带来的执行复杂性溢价。

1. 两个截然不同的广告市场，决定了完全不同的增长韧性与估值逻辑

Bernstein的核心洞见之一，是将全球广告市场拆分为两个本质上不同的市场。这不是一个简单的“线上vs线下”的二分法，而是一个基于客户结构、收入稳定性和增长驱动力的根本性区分。

市场一由传统媒体和广告代理构成。这个市场的特点是高度集中、强周期性，收入高度依赖数千个大客户。在经济下行周期，广告支出往往成为企业最先削减的成本项。这种结构决定了其收入的波动性极高，且与宏观消费信心指数高度相关。Bernstein的数据显示，这类市场的收入在宏观压力下可能出现两位数的季度波动。

市场二由数字平台主导。这个市场的增长引擎是数百万中小企业。Bernstein的数据揭示了一个被市场普遍低估的事实：中小企业贡献了主要数字平台约三分之二的广告收入。中小企业的广告支出决策逻辑与大企业截然不同——它们更关注可量化的投资回报率，而非品牌预算的周期性调整。这使得市场二的收入结构更具韧性，也更难通过传统的宏观指标进行预测。

这两个市场的分化意味着，简单地用“广告行业是周期性行业”这一笼统判断来给所有相关资产定价，很可能导致系统性的误判。数字平台的收入结构本质上更具防御性，而这种结构性优势尚未被市场充分定价。

2. 数字广告增长已进入第二阶段，零售媒体成为新的结构性增量来源

Bernstein将数字广告的增长划分为两个阶段，这一框架对于理解当前的投资机会至关重要。

第一阶段是过去十年的主导叙事：中小企业通过自助式广告工具大规模涌入数字平台。这一阶段的增长驱动力是“新客户 onboarding”。Bernstein强调，正是中小企业贡献了Meta和Alphabet约三分之二的广告收入，这解释了为什么这些平台能够在宏观逆风中依然保持增长韧性。

第二阶段始于2023年左右，增长驱动力发生了根本性转变。这一阶段的核心是媒体库存的扩张，以及随之而来的规模化个性化投放能力。零售媒体、电商平台、外卖应用、OTT平台等新兴广告库存的涌现，正在重新定义广告漏斗的边界。

零售媒体尤其值得关注。Bernstein将其描述为“结构性高利润率、近乎零成本的广告库存”。对于亚马逊这样的电商平台，或者DoorDash这样的配送平台，广告位几乎是一种纯增量收入——它们不需要额外生产内容，不需要新增带宽成本，只需要在现有用户界面中插入广告位。更重要的是，这些平台拥有购物意图数据，能够实现远超传统媒体的精准投放。

这意味着，原本被视为“非媒体公司”的零售商和应用平台，正在成为广告市场中最具增长潜力的新玩家。这一趋势的隐含含义是：数字广告市场的竞争格局正在从“平台vs平台”演变为“数据

[... middle omitted ...]

太小。它们占集团总收入的比例有限，无法从根本上改变整体增长轨迹和估值逻辑。与此同时，核心线性广告业务在美国等主要市场仍在持续下滑，DTC（直接面向消费者）业务的增长只能部分弥补这一缺口。

更深层的问题在于客户集中度和周期性暴露。传统媒体和广告代理的收入高度依赖少数大客户，这些客户的预算决策受宏观环境影响极大。即便数字化业务增长良好，只要核心业务仍然暴露于周期性风险，市场就不太可能给予这些公司高于行业平均的估值倍数。

这引出一个关键判断：传统媒体的数字化转型，在投资层面可能是一个“必要的防御动作”，而非“价值重估的催化剂”。除非数字化业务能够占到集团收入的50\%以上，否则市场对这类资产的定价逻辑不会发生根本性变化。

Bernstein的报告反复强调一个主题：广告市场的碎片化程度正在急剧上升。数字广告已经不再是一个单一的渠道，而是包含了搜索、社交、视频、零售媒体、CTV、程序化展示等数十个细分领域。每个领域都有不同的定价机制、不同的效果衡量标准、不同的投放优化逻辑。

\subsection{[R013] JPM：金属市场的真正拐点不在需求，而在供给成本的结构性重置}
{\small\color{refgray}归类：能源与大宗：供给成本结构性重置，金属需求分化}

JPM：金属市场的真正拐点不在需求，而在供给成本的结构性重置

中国金属消费数据正在传递一个比表面数字更复杂的信号。JPM最新一期的中国金属活动追踪报告显示，铜的消费连续第六周趋弱，铝的去库存却在加速，而锌的库存已攀升至2022年以来同期最高水平。如果只看到“需求疲软”，就会错过这份报告真正想要表达的核心判断：大宗商品投资者面临的风险，已经从需求侧的周期性放缓，转向供给侧的持续性成本膨胀——而后者正在从根本上改变金属和采矿行业的估值锚点。

这份于6月8日发布的周度追踪，覆盖截至6月5日的数据窗口。表面上看，它是一份常规的高频库存监测。但JPM的分析师团队在数据背后铺设了一条清晰的逻辑链：中国需求的确在放缓，但真正值得警惕的是，霍尔木兹海峡的持续关闭正在通过航运成本、能源价格和采矿耗材三个渠道，对金属供给端施加远超预期的压力。当需求侧的下行与供给侧的成本上行同时发生，市场的定价逻辑就必须被重新审视。

报告的核心结论是：JPM正在系统性降低其对EMEA（欧洲、中东和非洲）金属与采矿板块的敞口，仅对挪威海德鲁（Norsk Hydro）维持超配评级，原因在于后者的铝业务和区域溢价暴露。这不是一个简单的周期判断，而是一个关于成本结构如何重塑竞争格局的战略论断。

1. 铜的“弱”不是问题，问题的关键在于它为什么没有变得更弱

铜库存数据确实指向需求走软。截至6月5日当周，中国铜显性库存（SHFE加保税区）净变动基本持平，这是连续第六周库存去化动力不足。与3至4月份异常强劲的消费相比，当前确实是一个明显的减速。但JPM分析师明确指出，这种减速符合正常的季节性规律，并预计这种趋势可能还会持续数周。

真正值得关注的是库存的绝对水平。目前中国铜库存约为218千吨，处于五年季节性区间的下沿，但比2025年同期高出约40千吨。这个数据组合意味着什么？需求在放缓，但库存并没有出现爆炸式堆积。这暗示供给端可能也在同步收紧，或者下游用户正在谨慎地维持低库存策略。

更关键的信号来自洋山铜溢价。报告显示，该溢价已回落至约65美元/吨，反映出中国铜采购意愿下降。溢价的收窄与库存去化放缓形成了一组相互印证的关系——中国买家不再愿意为现货支付额外溢价，因为他们并不急于补库。

但这里需要追问一个JPM没有完全展开的问题：如果霍尔木兹海峡的关闭持续超出预期，铜的供给端是否会面临比需求端更剧烈的冲击？报告提到了航运成本的大幅攀升——从澳大利亚、巴西和南非到中国的铁矿石散货运费均已上涨超过50\%。铜精矿的运输同样依赖这些航线。溢价下降可能掩盖的是供给端真正的物理瓶颈。

与铜的疲弱形成鲜明对比的是铝。截至6月5日当周，中国铝显性库存净减少26千吨，这一去库存速度略强于同期的季节性平均水平。铝库存总量约为140万吨，虽然处于2019年以来同期最高水平，但去库存的加速说明下游需求正在吸收存量。

这种分化本身就值得深究。铜和铝虽然同属基本金属，但下游应用结构存在显著差异。铜更多与电力、建筑和制造业资本支出相关，而铝在交通运输、包装和轻型化材料领域的需求韧性更强。报告提到5月PMI数据中制造业活动的疲弱对金属需求形成拖累，但铝的表现说明，并非所有金属都在承受同样的压力。

JPM对挪威海德鲁的超配评级，正是建立在这种结构性差异之上。海德鲁的核心优势在于其铝业务和区域溢价暴露。当全球航运成本飙升、区域间价差扩大时，拥有本地化产能和区域市场定价权的生产商将获得竞争优势。这不是一个关于铝价本身涨跌的判断，而是一个关于“谁能在成本分裂的市场中保持利润”的筛选逻辑。

对于投资者而言，这意味着铝行业内部正在出现明显的分化。拥有低成本水电产能、靠近终端市场、且能够通过区域溢价对冲全球成本波动的企业，其估值逻辑正在与那些依赖进口原料和远距离运输的企业脱钩。JPM对整

[... middle omitted ...]

的金属指标之一。如果锌库存继续累积，它将对其他基本金属的需求预期形成向下修正的压力。

同时，锌的供给端同样面临成本压力。采矿所需的炸药、钢球、轮胎等耗材价格正在上涨，而能源成本——尤其是柴油和电力——也在攀升。锌价能否覆盖这些成本上升，取决于冶炼厂的议价能力和下游需求的弹性。从当前库存水平看，需求侧显然没有提供足够的支撑。

这份报告中最具冲击力的数据可能不是金属库存，而是铁矿石航运成本。从西澳大利亚到青岛的运费已从年初的约8美元/吨上涨至16美元/吨，从巴西图巴朗到青岛的运费从18美元/吨涨至38美元/吨，从南非萨尔达尼亚到北仑的运费从12美元/吨涨至28美元/吨。所有主要航线的运费涨幅均超过50\%。

这意味着什么？铁矿石的CFR（到岸价）目前约为107美元/吨，而FOB（离岸价）约为91美元/吨，价差达到16美元/吨。这个价差中，相当一部分是由运费驱动的。对于矿企而言，如果以FOB计价，他们并没有直接受益于运费的上涨；但如果以CFR计价，运费成本则会侵蚀利润。

\subsection{[R014] GS：服务器市场的真正拐点不是AI，而是传统服务器被重新定价}
{\small\color{refgray}归类：AI/算力：资本开支周期扩散，供应链瓶颈转移}

GS：服务器市场的真正拐点不是AI，而是传统服务器被重新定价

当市场目光几乎完全锁定在AI服务器的爆发式增长时，一份来自GS的最新研报揭示了一个更具结构性的信号：传统服务器市场正在经历一次被严重低估的重新定价。

这不是一个关于“AI还能涨多少”的故事。这是一个关于“非AI部分正在发生什么”的故事。GS援引650 Group的最新数据，将 年传统服务器市场的年均预测上调了约31\%，预计到2030年市场规模将达到1640亿美元。相比之下，AI服务器市场的同期预测上调幅度为18\%。数字本身已经说明问题：传统服务器的增速预期修正幅度，几乎是AI服务器的两倍。

为什么这个信号重要？因为在过去两年里，几乎所有关于数据基础设施的讨论都被AI服务器垄断。传统服务器被视为“存量业务”甚至“夕阳业务”。但这份报告的核心判断是：市场可能低估了企业端和超大规模云厂商在非AI计算资源上的系统性投入。这不是周期性补货，而是由工作负载结构调整、安全合规需求以及AI推理链条中传统算力的“伴生需求”共同推动的结构性变化。

报告发布的时间窗口也值得注意。C1Q26的数据显示，Dell Technologies在传统服务器市场的份额从一年前的20\%跃升至30\%，单季营收同比增长85\%。这不仅仅是市场份额的转移，更意味着整个行业的需求基础正在发生质变。

以下是我们从这份GS研报中提炼出的五个核心洞察，以及一个尚未被完全回答的关键问题。

1. 传统服务器市场预期的大幅上修，揭示的不是补库存，而是企业IT支出的结构性重置

GS在本次更新中，将 年传统服务器市场的年均预测上调了约31\%，单位出货量预期上调19\%，平均售价预期上调10\%。这个修正幅度在过去的服务器周期中非常罕见。传统服务器通常被视为成熟、低增长的市场，年均增速长期在个位数徘徊。但这次的预期路径显示，该市场将以13\%的五年复合年增长率增长，到2030年达到1640亿美元。

驱动这一修正的核心力量来自两个垂直领域：企业端和超大规模云厂商。650 Group对这两个细分市场的营收预期分别上调了约48\%和47\%。企业端的单位出货量预期上调30\%，意味着这不是单纯的价格上涨，而是实实在在的物理部署增加。

这意味着什么？企业正在同时进行两件事：一是将AI推理工作负载部署到经过优化的传统服务器架构上，二是在AI驱动的自动化、安全监控和数据分析等场景中，产生了对通用计算资源的额外需求。这两股力量叠加，使得传统服务器从“维持性支出”转变为“增长性支出”。对于服务器厂商而言，这意味着一轮可持续的需求周期，而不是一次性的库存回补。

2. Dell在传统服务器市场的份额跃升，说明规模优势和供应链执行力正在转化为定价权

这份报告中最引人注目的单个数据点是Dell在传统服务器市场的份额变化。C1Q26，Dell的传统服务器营收同比增长85\%，其中单位出货量增长24\%，平均售价增长49\%。市场份额从一年前的20\%跃升至30\%。更值得关注的是，增长是均衡的：面向新云厂商（neocloud）的营收增长102\%，面向企业端的营收增长91\%。

Dell的ASP增长49\%，远高于单位出货量增长24\%，说明该公司不仅卖出了更多的服务器，而且卖出了更贵、配置更高的服务器。这背后是Dell在供应链管理和客户关系上的双重优势。当客户从AI实验阶段进入规模化部署阶段时，他们倾向于选择能够提供端到端交付、全球服务网络和稳定供货能力的供应商。Dell在这方面的积累，正在转化为实实在在的定价权和市场份额。

对于投资者而言，这个信号的含义是：在服务器硬件这个看似商品化的市场，头部企业的竞争优势正在扩大。规模不再是成本优势的附属品，而是客户信任和议价能力的直接来源。

3. AI服务器市场的增长预期仍在加速，但竞争格局正在从

[... middle omitted ...]

AI服务器的供应生态正在从GPU芯片的单一垄断，向系统集成和交付能力的多极竞争演进。Dell和Super Micro的崛起说明，客户对AI服务器的需求已经从“只要能拿到GPU就行”转向了“谁能提供最优化、最可靠、最及时的整体系统”。这为具有大规模制造和全球交付能力的硬件厂商打开了新的增长空间。

4. HPE在AI服务器市场的份额流失，警示了缺乏规模优势的硬件厂商可能被挤出核心赛道

HPE在C1Q26的传统服务器营收同比增长20\%，但AI服务器营收仅增长11\%，市场份额从一年前的30\%降至17\%。这个数据点值得所有关注硬件赛道的人警惕。

在传统服务器市场，HPE仍维持了11\%的份额，表现稳定。但在AI服务器这个高增长、高投入的领域，HPE的份额正在被Dell和Super Micro蚕食。根本原因在于AI服务器的定制化程度高、交付周期紧、资本占用大，供应商需要具备强大的供应链管理能力和客户关系深度。HPE在这些维度上与Dell和Super Micro的差距正在拉大。

\subsection{[R015] 摩根斯坦利：面板价格周期的高点已至，市场定价的乐观情绪正在透支未来}
{\small\color{refgray}归类：未被主题摘要引用}

摩根斯坦利：面板价格周期的高点已至，市场定价的乐观情绪正在透支未来

这份摩根斯坦利在2026年6月发布的TFT-LCD面板价格展望报告，传递了一个清晰且克制的信号：面板价格在经历了一轮温和上涨后，正在接近本轮周期的顶部。对于关注面板行业的投资者而言，真正需要警惕的不是6月价格的短期波动，而是当前估值已经反映了太多尚未兑现的乐观预期。

报告的核心判断可以概括为一句话：TV面板价格将从三季度开始转跌，IT面板价格将从四季度开始走弱，而市场对“新兴业务”的期待——尤其是玻璃基板在先进封装中的应用——短期内难以贡献实质性的业绩增量。这意味着，当前台湾面板厂商的股价，已经跑在了基本面之前。

1. TV面板价格的红利期仅剩一个季度，供给纪律正在被需求走弱所抵消

摩根斯坦利预计，6月主流TV面板尺寸的价格将环比持平。这一判断本身并不令人意外，真正值得注意的是报告对后续趋势的定性：从三季度开始，TV面板价格将进入下行通道。

支撑这一判断的逻辑链条是清晰的。需求端，年初以来品牌厂商的提前备货正在退潮，面板订单动能明显放缓。供给端，尽管面板厂仍然展现出良好的产出控制纪律，但上半年强劲的出货量意味着下半年的出货量可能低于季节性水平，这将进一步削弱面板厂的议价能力。

报告隐含的一个关键洞察是：供给纪律虽然能够为价格提供下行支撑，但无法阻止价格拐点的到来。在需求走弱和出货量低于季节性的双重压力下，面板厂“控产保价”的策略效果将逐步递减。这本质上是周期力量对人为干预的修正。

从季度平均价格来看，报告预计2Q26 TV面板价格仍可实现低个位数的环比增长，但这已经是本轮上涨的尾声。对于投资者而言，三季度开始的价格下行预期，意味着当前估值中隐含的持续涨价假设需要被修正。

2. IT面板价格的相对稳定更像是一种被动平衡，而非主动强势

与TV面板相比，IT面板价格的走势更为平淡。报告预计6月显示器面板价格环比微涨0.2\%，笔记本面板价格环比持平，且这种相对稳定的状态有望在后续季度延续。

但“相对稳定”并不等同于“基本面强劲”。IT面板市场的供需格局更像是双方在低水平上达成的一种平衡。终端需求缺乏爆发性增长的动力，而面板厂在TV面板价格走弱后，可能将更多产能转向IT面板，从而压制IT面板价格的上涨空间。

报告对IT面板的表述值得细读：IT面板价格相对于TV面板“更具稳定性”。这种比较本身就是在暗示，IT面板并非一个具有独立上行逻辑的市场，而是一个跟随整体面板周期、波动幅度更小的细分领域。对投资者而言，IT面板的稳定只能提供防御性，而非进攻性。

3. 面板厂的结构性变化被市场过度定价，估值已透支中周期均值

报告在讨论行业结构性变化时，承认了积极的一面：中国面板厂已经从追求规模转向追求利润，政府补贴的减少强化了供给纪律，“以销定产”的商业模式正在被更多投资者接受，行业整合（如LGD出售广州工厂给TCL、Sharp关闭堺工厂）也在改善供给格局。

但摩根斯坦利同时指出，这些结构性变化“并没有被公平地定价”。这句话的潜台词是：市场已经将这些积极因素充分甚至过度地计入了股价，而忽略了价格周期本身的下行风险。

以台湾面板厂为例，友达光电的市净率已经达到1.4倍，群创光电达到1.9倍，均高于过去几年的周期峰值。即使考虑到新兴业务的潜在价值，摩根斯坦利认为当前估值提供的风险回报比已经不再吸引人。相比之下，报告更偏好京东方，理由是规模优势、更多元化的收入组合以及相对便宜的估值。

这一判断的核心含义是：结构性变化是真实的，但它并不能免疫于周期性下行。当周期向下时，即使是结构性改善后的行业，其盈利能力和估值也会受到压力。当前市场对结构性变化的定价，可能已经透支了未来2-3年的改善空间。

4. 玻璃基板在先进封装中的应用：方向正确，但 年不是

[... middle omitted ...]

间维度上存在显著的不确定性。更重要的是，面板厂并非唯一试图进入这一市场的参与者。当半导体材料和设备厂商也在积极布局时，面板厂能否在竞争中获得显著份额，本身就是一个需要验证的问题。

报告没有完全否定这一方向，但明确表示“有意义的贡献需要时间”。对于投资者而言，如果当前估值中已经包含了玻璃基板带来的增长溢价，那么这一溢价的存在本身就是一种风险。

尽管报告对价格方向给出了明确的判断，但有两个关键问题并未完全展开。

第一，价格下行的幅度会有多大？报告提到“跌幅不会显著”，理由是供给纪律提供了下行支撑。但“不显著”是一个相对概念，如果需求超预期走弱，或者品牌厂商大幅削减库存，供给纪律能否真正抵御价格下行压力，仍然存在不确定性。

第二，价格下行的节奏是否可能提前？报告预计TV面板从三季度开始下跌，IT面板从四季度开始走弱。但这一判断建立在当前需求逐步放缓的基准假设之上。如果中国经济复苏不及预期，或者全球消费电子需求在通胀压力下进一步萎缩，价格下行的起点可能比报告预期更早。

\subsection{[R016] 摩根斯坦利：市场低估了这份利率新规对银行净息差的修复力度}
{\small\color{refgray}归类：中国政策与金融：制度重塑进行时}

6月5日，中国人民银行发布了《人民币存款和贷款利率管理规则》的征求意见稿。这份文件看似技术性极强，聚焦于利率计算、罚息规则和机构职责界定，但摩根斯坦利在第一时间发布的研报中给出了一个清晰且值得重视的判断：这是对银行净息差趋势的明确利好，且影响同时作用于资产端收益率和负债端成本。

这不是一次简单的规则更新。1999年的旧框架运行了超过25年，期间中国利率市场化改革经历了贷款基础利率(LPR)形成机制、存款利率自律上限调整、以及近年来的多次降息周期。旧规则已经无法覆盖当前市场环境下的复杂定价行为。新规的核心意图，不是微调技术细节，而是为市场化的利率定价建立一套清晰的规则书，并明确传递出打击不公平、非法竞争的信号。

对于关注中国银行业资产定价逻辑的投资者而言，这份征求意见稿的意义不在于它改变了什么短期数字，而在于它从制度层面为银行的收入恢复和利润稳定增长提供了支撑。市场当前的讨论多停留在“政策意图”层面，但摩根斯坦利的分析框架提示我们，真正需要被定价的，是这一规则对银行间竞争格局的结构性重塑。

存款市场的无序竞争长期以来是中国银行业净息差压力的重要来源之一。部分银行，尤其是中小银行，为了争夺存款规模，采取了多种变相高息揽储的手段。新规明确禁止了这些行为，包括非法手工补息、突破自律机制上限、以及存贷挂钩等操作。

这不仅仅是监管态度的重申。旧框架下，这些行为的界定和处罚缺乏明确的规则依据，导致监管执行存在灰色地带。新规将禁止性条款写入正式的《利率管理规定》，意味着从此有了清晰的执法标尺。摩根斯坦利在报告中指出，这对银行净息差在负债端成本上的改善是积极的。

更深一层的含义在于，高息揽储的禁令将加速存款市场的分化。过去依赖价格战获取存款的银行，将不得不转向通过服务、品牌和产品竞争力来留住客户。那些拥有稳定核心存款基础、客户粘性强的银行，其负债成本优势将进一步巩固。这相当于从制度层面为优质银行的低成本负债护城河加了一道墙。

2. 贷款利率透明化正在重塑资产端的定价纪律，而非简单的“降息”

新规要求贷款人必须明确展示年化利率，且贷款成本必须包含与贷款直接相关的利息和费用，同时明确罚息年化利率的计算方式。这一条被市场普遍解读为保护借款人的措施，但摩根斯坦利的分析视角更为犀利：它实际上在重塑银行资产端的定价纪律。

过去，部分银行通过隐藏费用、复杂计息方式或模糊的罚息条款来实现隐性加价。这种不透明的定价模式，一方面损害了借款人利益，另一方面也使得银行间的价格竞争变得混乱和不可比。新规要求所有成本显性化，意味着银行无法再通过“暗门”来调整实际收益率。

这带来的直接后果是，资产端的定价竞争将从“谁更会隐藏费用”转向“谁的风险定价能力更强”。银行必须更加精准地评估借款人的信用风险，并据此设定透明的利率。那些依赖信息不对称和收费模糊性来获取超额收益的银行，其资产端收益率可能会面临下行压力。但与此同时，整个行业的定价效率会提升，恶性价格战的空间被压缩。摩根斯坦利将此视为对银行收入恢复的长期正面因素，因为它降低了系统性的定价风险。

新规将罚息从过去的“加收一定比例”改为“协商确定罚息利率”。这一变动看似技术性调整，但其对银行损益表的影响不容忽视。旧框架下，罚息利率往往是固定的、标准化的，银行缺乏弹性空间。新规允许银行与借款人在合同中协商罚息利率，这实际上赋予了银行在风险发生后更大的定价自主权。

对于银行而言，这意味着两件事。第一，当借款人出现违约时，银行可以通过协商获得更高的罚息收入，这直接增加了非利息收入。第二，更灵活的罚息机制可能促使银行在贷款发放时更加审慎地评估风险，因为罚息条款本身也成为了风险定价的一部分。摩根斯坦利的分析框架中，这一变化对银行利润稳定增长的支撑是结构性的，而非一次性的。

当然，这一规则的实际效果高

[... middle omitted ...]

以根据人民银行规则和商业原则自主设定存贷款利率，但同时也强化了人民银行及其分支机构的监督职责，并指导市场利率自律机制发挥作用。摩根斯坦利在报告中强调，这为市场化的定价建立了清晰的规则书。

这意味着，未来的银行竞争将从“监管套利”转向“管理红利”。过去，谁能更巧妙地绕开监管限制，谁就能获得短期优势。但在新规则下，银行的竞争焦点将回归到核心能力：风险管理、成本控制、客户服务和产品创新。那些能够将规模优势转化为议价权、将风控能力转化为定价优势的银行，将在新的规则下持续受益。

5. 报告尚未完全回答的关键问题：执行力度与非银机构的传导效应

摩根斯坦利的分析逻辑清晰，但有两个关键变量在报告中并未完全展开，值得读者高度关注。

第一个是执行力度。征求意见稿转化为正式规则后，各地人民银行分支机构的执法力度是否一致？历史上，中国金融监管在执行层面存在地区差异，尤其是在存款市场竞争激烈的区域，高息揽储的变种形式可能以新的面目出现。如果执行打折扣，新规对净息差的改善效果将大打折扣。

\subsection{[R017] GS：新能源车市场真正的压力测试，不是销量，而是定价权的转移}
{\small\color{refgray}归类：新能源车：渗透率新高下的盈利考验}

GS：新能源车市场真正的压力测试，不是销量，而是定价权的转移

2026年5月，中国新能源车零售渗透率达到63\%，批发渗透率突破61\%。这两个数字看起来像是行业高歌猛进的注脚。但GS最新发布的周度图册（China New Energy Vehicle Weekly Chartbook, Week 23）揭示了一个更值得警惕的图景：在渗透率创新高的同时，终端折扣正在同步扩大，部分品牌的订单环比出现两位数下滑，而锂价在6月初一周内下跌超过6\%。

这不是一个简单的“需求好”或“需求差”的故事。GS这份报告的核心价值，在于它把订单、折扣、库存、原材料价格和事件日历编织在一起，指向一个正在成形的行业拐点：新能源车市场正在从“规模扩张竞赛”切换进入“盈利韧性考验”。而市场对这一切换的定价，可能仍然不够充分。

5月新能源车零售渗透率63\%，批发渗透率61.1\%，相比4月的62.8\%和57.3\%均有提升。单看这个数字，很容易得出“电动化势不可挡”的判断。但GS跟踪的周度订单数据给出了一个更精细的叙事。

2026年第23周（6月1日至6月7日），主要新能源品牌合计订单同比增长7\%，但环比下降了15\%。GS明确指出，这一环比下滑是“在新车型上市初期脉冲消退之后”发生的。这意味着，5月下旬新车发布带来的订单集中释放效应正在衰减，市场正在回归到更真实的周度需求水平。

更值得关注的不是总量，而是结构。在环比维度上，吉利、特斯拉、比亚迪表现出“防御性增长”——环比分别增长12\%、1\%和下降2\%。这三家头部企业的订单韧性，与行业内其他品牌形成了鲜明对比。而在累计同比维度上，蔚来以85\%的同比增速领跑，问界（HIMA）和特斯拉分别增长29\%和3\%。

这意味着两件事：第一，头部集中效应仍在加剧，但“头部”的定义正在变化——不再仅仅是销量规模，而是订单获取的稳定性；第二，蔚来和问界的高增速背后，是新车型周期的强力驱动，但这种驱动能否持续，取决于后续订单转化率和交付节奏。

对于投资者而言，渗透率数据已经不再是区分胜负的关键变量。真正需要跟踪的，是每个品牌在“脉冲消退期”的订单回落斜率。斜率越陡，说明该品牌对新车型的依赖程度越高，产品矩阵的防御性越弱。

2. 新能源车与燃油车的折扣同步走阔，价格战正在从“战术”演变为“结构性压力”

GS报告中最令人不安的图表之一，是新能源车和燃油车终端折扣的同步变化。

截至6月6日，新能源车平均经销商折扣为7.79\%，高于一周前的7.75\%。燃油车的折扣更深，达到19.54\%，同样环比走阔。这意味着，无论是电动化阵营还是传统阵营，终端都在以更大的价格让步换取销量。

这里的关键洞察不在于折扣本身，而在于两个趋势的叠加。新能源车的渗透率在上升，但折扣也在扩大。这打破了“渗透率提升=定价权增强”的简单线性逻辑。实际上，当新能源车渗透率突破60\%后，它面临的竞争格局不再是“替代燃油车”，而是“新能源品牌之间的内卷”。后者的定价逻辑，更接近存量博弈。

GS报告中的分品牌折扣数据进一步佐证了这一点。比亚迪的平均经销商折扣为5.14\%，与一周前持平，但高于去年同期。这意味着，即使是行业成本控制最强的玩家，也无法完全避免价格侵蚀。而其他品牌——尤其是那些依赖单一爆款或缺乏规模效应的品牌——面临的折扣压力只会更大。

对于产业决策者而言，这意味着一个残酷的现实：在渗透率突破60\%之后，价格战不会因为“电动化成功”而结束，反而会因为“同质化竞争加剧”而成为常态。企业的核心竞争力，正在从“能否造出电动车”转向“能否在不依赖大幅折扣的前提下维持市场份额”。

3. 锂价加速下跌，上游成本红利正在向中下游传导，但受益并不均匀

报告中的另一个关键数据点是上游电池定价动态：电池级碳酸锂价格降至每吨16.85万元，周度环比下

[... middle omitted ...]

个信号：上游产能过剩的格局尚未出清。这意味着，即使车企在未来几个季度看到电池采购成本下降，这种下降的可持续性也存疑。如果上游产能出清加速，锂价可能触底反弹，届时成本红利的窗口期可能比市场预期的更短。

对于投资者而言，这不是一个简单的“成本下降利好整车厂”的故事。真正需要判断的是：哪些车企能够在这一轮成本红利窗口期内，将节省的成本转化为利润，而非重新投入价格战。如果一家车企在锂价下跌的同时，终端折扣反而扩大，那么成本红利就被竞争完全吞噬了。

4. 燃油车经销商的折扣深度揭示了“替代”的终点，而不是起点

GS报告对燃油车折扣的持续跟踪，提供了一个常被忽视的视角。截至6月6日，燃油车平均经销商折扣高达19.54\%，虽然略低于去年同期的23.03\%，但仍处于历史高位。

这个数字的含义非常直接：燃油车经销商正在以接近八折的价格抛售库存。这不是一个短期的清仓行为，而是一个结构性趋势——燃油车的需求基础正在被新能源车不可逆地侵蚀，经销商不得不通过持续的大幅折扣来维持现金流。

\subsection{[R018] GS：公积金改革范围超预期，但市场真正低估的是供给侧的再定价信号}
{\small\color{refgray}归类：地产与消费：制度优化与结构分化}

GS：公积金改革范围超预期，但市场真正低估的是供给侧的再定价信号

中国房地产市场的焦点正在从“成交量能走多远”转向“制度变量如何重塑行业定价逻辑”。GS在最新一期周报中给出了一个值得反复推敲的判断：公积金改革的范围超出了此前预期，但市场对这一变化的定价可能仍然偏保守。周度交易量虽然环比回落，但同比仍在扩张，更重要的是，情绪指标保持稳定——这不是典型的政策刺激后的快速衰减，而更像是一种新的稳态正在形成。

这份报告的核心价值不在于周度数据本身，而在于它提示了一个正在发生的结构性变化：政策工具正在从单纯的“需求刺激”转向“制度优化+供给调控”的组合。公积金改革、新开工与竣工的持续分化、以及库存去化周期的微妙变化，这些信号叠加在一起，指向一个更深刻的命题——中国房地产市场的再平衡，可能比多数投资者预期的更早到来，但路径也将更加分化。

1. 公积金改革的制度红利被低估，其影响将超越短期的需求刺激

GS在报告中明确指出，住建部最新发布的《住房公积金管理条例》修订征求意见稿，在三个维度上超出了市场预期。最值得关注的是两点：一是将自住住房装修和物业管理费纳入提取范围，二是明确将灵活就业人员纳入缴存体系。后者的潜在覆盖人群据估算超过2亿人。

这不仅仅是扩大公积金的使用场景，更是一次制度性的“扩围”。过去市场对公积金改革的预期主要集中在贷款利率下调和使用门槛降低上，但这次修订触及了公积金作为“住房金融基础设施”的核心定位。当灵活就业人员能够合法、合规地参与公积金体系，意味着这部分此前被排除在正规住房金融之外的购买力，有了进入市场的通道。

从时点上看，这一改革发生在市场情绪已经企稳、但交易量尚未显著回升的阶段。GS的数据显示，尽管第23周新房销售面积环比下降13\%，但同比仍增长15\%，二手房更是同比增长31\%。情绪指标——经纪人预期指数（CSI）环比上升0.4个百分点——表明市场并非在“降温”，而是在“换挡”。公积金改革的制度红利，将在换挡期提供额外的购买力支撑。

但报告也提示了一个关键未解问题：公积金贷款的利率是否会进一步下调？以及，大量沉淀的公积金存款如何更有效地投放出去？这两个问题将决定制度红利的实际释放力度。目前来看，改革框架已经搭建，但具体的“弹药”尚未完全就位。

2. 交易量环比回落不是风险信号，而是市场从脉冲式修复转向稳态运行的标志

GS第23周的数据显示，新房销售面积环比下降13\%，二手房下降7\%，看房和认购量也环比下降约5\%。乍看之下，这似乎是市场动能在减弱。但深入分析会发现，环比回落是在前一周同比大幅扩张（新房+9\%、二手房-2\%）之后的正常修正，且同比表现依然强劲——新房+15\%、二手房+31\%。

更值得关注的是情绪指标的稳定性。新房搜索活动环比持平，新增挂牌节奏与5月相当，经纪人预期指数（CSI）环比上升，卖房者报价指数（CAI）也保持稳定。这说明市场参与者对价格走势的判断没有发生方向性变化。

GS还提供了一个关键参照系：2025年以来的累计数据显示，新房销售面积同比下降13\%，但二手房累计同比增长1\%。这一“新弱二强”的格局并非首次出现，但报告将其置于更长的历史维度中观察——当前二手房交易量已较2024年全年平均水平高出22\%，较2023年高出8\%。这意味着二手房市场正在成为真正的“价格发现”场所，而新房市场则更多受到供给端约束的影响。

对于投资者而言，这意味着“量”的波动不应被过度解读为趋势逆转。真正需要关注的是“价”的稳定性——只要价格预期没有出现系统性下行，交易量的正常波动就只是市场在寻找均衡点的过程。

GS通过其独有的GSPC（GS房地产完工指数）模型，对竣工数据进行了前瞻性估算。报告指出，基于浮法玻璃供需模型，5月竣工面积可能录得高双位数的同比降幅（与此前4月NBS公布的

[... middle omitted ...]

的现金流和资产负债表；另一方面，新增供给的减少将限制未来可售资源的规模，从而对开发商的收入增长形成约束。

GS在估值部分提供了一个值得深思的数据：当前其覆盖的开发商股价平均较2026年底的净资产价值（NAV）折价28\%（离岸）和27\%（在岸），对应的2026年预期市净率分别为0.5倍和0.4倍。这一估值水平已经低于历史上的多个周期底部（如2008年下半年、2011年下半年、2014年上半年）。但报告没有直接回答的问题是：在“去增量”的供给格局下，NAV本身是否需要被重新定价？如果可开发土地的价值系统性下降，那么基于历史NAV的折价可能并不具有真正的“安全边际”。

4. 开发商之间的分化正在加剧，但市场的定价尚未充分反映这一趋势

GS在周报中详细列出了不同类别开发商的股价表现。第23周，其覆盖的“较强国企开发商”股价平均下跌7\%，其中华润置地（1109.HK）和中国海外发展（0688.HK）表现相对较好，仅下跌3\%；而民营开发商和其他国企开发商分别下跌7\%和8\%。

\subsection{[R019] 摩根斯坦利：工程机械的景气周期比市场想象的更强}
{\small\color{refgray}归类：工业与制造业：资本开支周期超预期扩散}

市场对5月中国挖掘机销量的预期是“增速放缓”，但实际数据给出了完全不同的信号。摩根斯坦利最新发布的5月挖掘机销售追踪报告显示，当月总销量同比增长36\%，其中国内销量增长39\%，出口增长34\%。这三个数字都显著超出市场此前的悲观预期。

更值得关注的是，这份增长不是靠低基数或政策脉冲的短期刺激。报告明确指出，5M26累计销量同比增长25\%，国内和出口分别录得18\%和33\%的累计增速。这意味着，从年初至今的景气度是可持续的，而非单月波动。

为什么这份报告值得产业决策者和投资者仔细阅读？因为它揭示了一个正在发生的结构性变化：中国工程机械行业正在经历从“国内周期驱动”向“全球竞争力驱动”的转变。摩根斯坦利在报告中不仅更新了数据，还提供了与三一重工、中联重科管理层交流后的最新判断——两家公司对2Q26的展望都偏向积极。

以下是我们从这份报告中提炼出的四个关键洞察，以及一个报告尚未完全回答的核心问题。

1. 国内市场的韧性来自“提前发债”与“电动化替换”的双重支撑，而非地产复苏

市场对中国工程机械国内需求的最大担忧，是房地产投资持续疲弱。但5月国内销量39\%的同比增长，说明推动力来自其他方向。

摩根斯坦利在报告中给出了两个核心驱动因素。第一是地方政府专项债券（LGSB）的提前发行。资金到位后，基建项目开工率提升，直接拉动了挖掘机的采购需求。第二是电动化挖掘机的替换需求。随着环保政策趋严和电动设备经济性改善，存量设备的更新换代正在形成独立的增量市场。

这两个因素叠加，使得国内需求在当前宏观环境下依然保持了强劲增长。5月销量环比下降14\%，这属于正常的季节性波动，不应被过度解读为趋势逆转。

这里的含义非常明确：国内工程机械需求已经不再完全依赖房地产周期。 对于关注中国基建投资和制造业升级的读者来说，这意味着需要调整对行业景气度的传统分析框架。过去“地产下行=工程机械承压”的逻辑，正在被“基建前置+电动化替换”的新组合替代。

2. 出口的韧性比表面数据更强：市场份额扩张正在抵消地缘风险

出口同比增长34\%，累计增长33\%——在红海局势紧张、燃料成本上升的背景下，这个数字本身就值得深究。但摩根斯坦利报告给出的解释才是真正的洞察：增长主要来自市场份额的持续提升。

报告提到，中国工程机械企业在欧洲、拉丁美洲、非洲等多个区域的市场份额正在扩大。这意味着，即便全球经济增速放缓或局部需求波动，中国企业通过产品竞争力、渠道布局和服务能力，正在从卡特彼勒、小松等国际对手手中夺取份额。

这是典型的“结构性增长”而非“周期性增长”。出口数据背后的驱动力，是中国制造业从“低成本制造”向“全球化运营”的升级。三一重工和中联重科的管理层在交流中均对海外需求表示乐观，这进一步印证了上述判断。

对于投资者而言，需要关注的不是单月出口增速的波动，而是中国企业在海外市场的市占率曲线。只要这条曲线保持上升趋势，出口增长就有望跑赢全球行业整体增速。

报告披露了一个重要细节：包括三一重工、徐工机械、柳工在内的主要主机厂，从5月中旬到6月陆续实施了提价。摩根斯坦利认为，这应支持行业利润率修复。

提价发生在销量高增长的背景下，说明供需格局对主机厂有利。过去几年，中国工程机械行业经历了激烈的价格战，利润率承压。当前的需求强度和竞争格局变化，给了企业重新定价的空间。

但这里有一个需要谨慎看待的问题：提价的可持续性。报告没有详细讨论竞争格局是否已经稳定到足以支撑长期提价。如果未来需求增速放缓，或者新进入者增加供给，提价能否维持就是未知数。摩根斯坦利的研报估值部分提到，三一重工的目标市盈率是23倍，这与2016-17年上行周期中的平均估值水平一致。这个估值隐含的假设，是利润率能够持续改善。

这是报告尚未完全回答的关键问题之一。对于希望深入理解行业定

[... middle omitted ...]

台等多个子行业，但不同领域的景气度可能存在显著差异。例如，电动化渗透率较高的品类，替换需求可能更强劲；而依赖大型基建项目的重型设备，则更容易受到全球利率环境的影响。

此外，摩根斯坦利在报告中提到了对恒立液压的估值方法，其中包含了对人形机器人零部件业务的单独DCF估值。这提示我们，部分工程机械产业链公司正在向机器人、自动化等新领域延伸，这些业务的增长逻辑与传统工程机械完全不同。行业上行周期的定义，可能需要从“设备销售周期”扩展为“制造业升级周期”。

对于产业决策者而言，这意味着不能简单地将“工程机械看好”等同于所有子行业和所有公司都受益。需要建立更精细的观察框架：关注企业在新领域的技术布局和收入占比，关注各子行业的供需关系和定价能力，关注海外各区域的增长差异。

这份报告提供了丰富的数据和判断，但有两个问题值得进一步探讨。

第一，国内需求的可持续性。提前发债的效果是一次性的还是可持续的？电动化替换需求的高峰何时到来？如果这两个因素消退，国内需求能否找到新的增长点？

\subsection{[R020] 摩根斯坦利：市场正在低估这一轮资本开支周期的宽度和持续性}
{\small\color{refgray}归类：AI/算力：资本开支周期扩散，供应链瓶颈转移 / 工业与制造业：资本开支周期超预期扩散}

这份报告的标题或许可以更直白：摩根斯坦利在近期对中国工业企业的调研后，得出的核心判断是，市场当前对AI驱动的资本开支周期的理解，仍然过于集中在半导体和算力硬件本身，而低估了这一轮周期向更广泛设备领域传导的力度、速度以及结构性含义。

这不是一份关于“AI服务器出货量”或“GPU需求”的报告。这是一份关于“当AI开始重塑制造业的每一个环节，谁在真正获益”的调研记录。摩根斯坦利的分析师团队在近期密集走访了超过20家中国工业领域的头部公司，涵盖自动化、机器人、电池设备、工程机械、阀门等多个细分赛道。调研反馈的强度，甚至超过了他们自己在今年3月的前一次调研。

报告传递的信号非常清晰：这一轮资本开支周期，正在从AI硬件设备，向通用设备、自动化、以及更广泛的先进制造领域扩散。而且，这种扩散不是线性、温和的，而是伴随着技术迭代、产品升级和供应链重构的复合型扩张。对于产业决策者和投资者而言，理解这一轮周期的宽度，比猜测某一季度的出货量数字重要得多。

1. AI需求正在从“算力基建”向“设备升级”传导，且传导效率超出预期

调研中最明确的信号来自AI设备供应链。摩根斯坦利发现，几乎所有与AI相关的设备供应商，对2026至2027年的订单和收入增长都给出了极为乐观的展望。这种乐观不是孤立的，而是基于三个结构性的驱动力。

第一，AI直接推动了产品和设备的升级，带来更高的单价和利润率。报告中提到的案例非常具体：大族激光用于AI服务器、光模块和TGV的超快激光钻孔设备，单价和毛利率显著高于传统产品；麦格米特在Rubin电源架中的潜在机会，价值含量和利润率更高；GKG精密应用于AI服务器的锡膏设备，毛利率达到约65\%，远高于其他下游应用。这些数字说明，AI对设备端的拉动不是简单的“多买几台”，而是“买更贵、更精密的设备”。

第二，技术迭代和量产化正在创造增量设备需求。调研发现，光模块组装自动化、X-ray CT检测设备、AOI检测设备等，都在从人工或半自动向全自动过渡。这种“替代人工”的自动化需求，本身就是AI产业化的副产品，却为设备公司创造了全新的市场空间。

第三，国内设备商正在快速扩大市场份额。在快速的产品升级和交付周期中，国内企业的响应速度和产能灵活性，正在成为相对于海外竞争对手的明显优势。报告明确指出，海外同行的产能瓶颈，正在为国内企业创造结构性机遇。

这三点叠加，意味着AI对设备的需求不是一次性的资本开支脉冲，而是持续的、升级驱动的需求曲线。摩根斯坦利将这一轮周期定义为“超级资本开支周期”，其核心特征就是规模大、迭代快、需求持续。

2. 苹果供应链的反馈比预期更强，2026至2027年的创新周期正在支撑设备需求

报告中的一个细节值得注意：苹果供应链相关公司的反馈，比分析师预期更为积极。这不仅仅是因为苹果本身的产品周期，而是因为苹果正在将设备需求从传统的3C自动化，扩展到3D打印、均温板、摄像头模组和UTG玻璃等全新领域。

报告提到，苹果将在今年推出折叠屏手机，并在明年推出六款新手机型号，其中包括一款纪念版机型。这意味着，2027年将成为重要的手机创新年份。对于设备商而言，这意味着什么？大族激光的管理层给出了一个具体的数字：2026年来自某客户的均温板激光焊接和自动化业务收入预计达到约10亿元人民币。这还只是单一客户、单一技术路线下的数字。

更重要的是，苹果在3D打印上的布局正在进入更大的产品周期窗口。报告指出，某客户计划在2026年下半年下达3D打印设备订单，首批订单（总共4000台需求中的一半）将优先给到大族激光。虽然3D打印的收入确认主要落在2027年，但这已经为设备商提供了一个明确的增长期权。

这一节的核心含义在于，消费电子领域的设备需求，正在从“跟随产品代际升级”转向“引领制造工艺变革”。当终端品牌

[... middle omitted ...]

。

与此同时，国内的人形机器人产业链正在变得更加活跃。减速器、轴承、电机等核心零部件的供应商，如绿的谐波、兆威机电、信捷电气、步科股份、华研精机等，都在应对高强度的出货预期，并持续扩张产能。报告特别指出，在谐波减速器领域，领先的国内企业应该能够脱颖而出。

恒立液压的情况也值得关注。作为美国领先人形机器人公司的丝杠供应商，恒立液压的份额被报告指出超过70\%。虽然当前人形机器人的产量有限，但恒立液压已经在墨西哥布局产能，目标是在年底前支持约10万台机器人的产能需求。更重要的是，恒立液压的工业丝杠业务正在成为第二个增长曲线，已有约2.5亿元的在手订单，月产能目标从目前的4000台提升到年底的15000台。

这里的关键洞察是：人形机器人的短期节奏放缓，并不改变国内供应链的中长期布局逻辑。恰恰相反，当海外领先企业的量产节奏出现不确定性时，国内供应链的自主推进和多元化客户开发，正在成为更重要的竞争要素。

4. 通用设备和工程机械的基本面同样扎实，全球化产能正在成为新护城河

\subsection{[R021] 摩根斯坦利：市场低估的不是AI需求，而是“旧存储”的供给再定价}
{\small\color{refgray}归类：AI/算力：资本开支周期扩散，供应链瓶颈转移}

摩根斯坦利：市场低估的不是AI需求，而是“旧存储”的供给再定价

Computex 2026刚刚落幕，当绝大多数市场目光聚焦于Nvidia的Vera CPU、Agentic PC以及云端AI芯片的算力竞赛时，摩根斯坦利半导体团队发布了一份视角独特的会后总结。这份报告的核心判断，与市场主流叙事形成微妙但重要的错位：真正值得关注的超额收益机会，可能不在云端芯片，也不在AI PC，而在被多数人视为“过时”的旧存储领域。

报告标题直接点明了这一排序——“Cloud, PC and old memory”，而最后的“old memory”恰恰是这份报告最想传达的增量信息。摩根斯坦利认为，市场对2027年CPU、GPU及周边芯片的强劲前景已有一定共识，对AI PC的谨慎也相对充分，但真正被低估的，是SLC NAND、高密度NOR等“旧存储”在AI推理时代的结构性短缺。这一判断，叠加近期板块股价波动，为投资者提供了一个具有吸引力的入场窗口。

关于云端半导体，市场的主要分歧从来不是需求强弱。GPU服务器持续 ramp up，CPU服务器因AI与agentic应用而需求旺盛，这些已是共识。摩根斯坦利在Computex上获得的关键增量信息，更多集中在供给侧的细节变化。

ARM服务器DRAM含量偏低的问题，报告明确指出这本质上是“supply issue”——不是需求不足，而是供给端尚未充分适配。这一判断意味着，一旦供应链完成调整，相关内存接口芯片的需求弹性可能超出当前市场预期。与此同时，CXL（Compute Express Link）正在成为扩大x86 CPU中RDIMM使用的关键推动力，这直接利好全球内存接口芯片的领导者。

更具时间信号意义的是，CPO相关的组件可能从原本的2027年一季度提前至2026年四季度开始拉货。提前一个季度的pull-in，在半导体产业的库存与订单周期中，往往意味着下游客户对供需紧张的提前预判。这不能简单视为短期波动，而应理解为整个产业链对2027年算力基础设施升级节奏的集体前置。

2. Agentic PC的定价揭示了“量价矛盾”的核心困境

AI PC是本次Computex的另一焦点。Nvidia CEO将未来AI PC描述为“始终运行、易于访问、具备强agentic能力的助手”，并认为这将从根本上重塑PC的形态与价值。然而，摩根斯坦利的渠道调研揭示了现实与愿景之间的巨大鸿沟。

搭载N1X的AI PC定价需达到2899美元，N1型号也要1799美元。这一价格带已经逼近甚至超过许多高端工作站，远远超出主流消费市场的承受范围。报告对此的表述克制但清晰：“AI PC is a positive driver, but could take time to materialize”——积极驱动因素存在，但需要时间兑现。

更值得关注的不是价格本身，而是价格背后的结构性矛盾。AI PC的高定价源于CPU、内存、屏幕等多重升级的叠加——128GB DRAM搭配4TB SSD、4K屏幕、RTX Spark GPU，这些组件成本的刚性上涨，使得终端品牌几乎没有降价空间。然而，消费者的支付意愿是否同步提升？摩根斯坦利的调研显示，下游客户“似乎愿意接受涨价”，但这更多是供应链层面的被动接受，而非终端需求的主动验证。

这里的核心洞察在于：AI PC的GM（毛利率）稳定性，更多依赖于下游客户对成本上涨的短期容忍，而非AI PC本身创造了足够的需求溢价。一旦晶圆代工厂在2026年三季度进一步提价，这种脆弱的平衡可能被打破。PC半导体公司的利润率韧性，远比当前股价所反映的更为脆弱。

这是整份报告最具差异化的判断。摩根斯坦利明确提出，SLC NAND（单层单元NAND闪存）在未来两年可能迎来更强的增长前景。这一

[... middle omitted ...]

AND淘汰的SLC，因为其在低延迟、高耐久性方面的不可替代性，重新获得了数据中心架构中的战略地位。

第二，SLC NAND被用于增强传统HDD的性能。这一应用场景虽然技术细节尚未完全公开，但其逻辑在于：在混合存储架构中，SLC可以作为HDD的高速缓存层，弥补机械硬盘在随机读写性能上的先天不足。

更令人惊讶的信号来自NOR Flash。报告披露，高密度NOR的渠道价格已从先前的1美元飙升至8美元。驱动因素是什么？摩根斯坦利认为，VR/AR设备中的NOR含量是传统游戏机的2倍。这一价格跃迁不是温和的周期性复苏，而是一种近乎断崖式的供给短缺信号。

综合来看，旧存储赛道的结构性机会来自三个叠加因素：AI推理对低延迟存储的新需求、GPU Direct Storage对SLC的重新定义、以及VR/AR对高密度NOR的需求爆发。这些都不是周期性因素，而是由AI基础设施架构演变驱动的长期变化。

4. 报告尚未完全回答的关键问题：GPU Direct Storage的商业化节奏

\subsection{[R022] NOM：市场低估了中国造船的结构性护城河，真正的风险不在需求侧}
{\small\color{refgray}归类：区域市场：亚洲结构性机会与分化}

NOM：市场低估了中国造船的结构性护城河，真正的风险不在需求侧

过去一年，关于中国造船业市场份额“腰斩”的叙事在媒体上反复出现。从2024年占全球新船订单约四分之三，到2025年中期一度被韩国反超，这组数据被广泛解读为华盛顿制裁政策的初步成效。然而，NOM最新发布的研报提出了一个更值得深思的判断：这种波动主要反映的是周期性因素和统计口径的扭曲，而非中国造船业竞争力的结构性衰退。真正需要关注的，不是中国份额的短期起伏，而是全球造船业供给格局正在经历的深层变化——以及美国在试图重建产能时所面临的硬约束。

这份报告的核心洞察在于：市场将注意力过多放在了需求端的地缘政治扰动上，却低估了中国造船业在供给端已经建立的、由成本、规模和绿色转型共同构筑的护城河。同时，华盛顿的“盟友造船”战略虽然声势浩大，但其推进速度恰恰暴露了美国自身工业能力的瓶颈。理解这两条线的交汇，才是把握未来几年全球航运与造船资产定价的关键。

1. 中国份额的V型反弹揭示的是“统计幻觉”，而非竞争力衰退

NOM的分析首先澄清了一个被广泛误读的关键事实。中国在全球新船订单中的份额确实经历了剧烈波动：从2024年的约70\%（甚至更高）降至2025年中期的低点，彼时韩国在月度订单上短暂领先。但到2025年全年，中国份额已回升至约63\%，并在2026年第一季度进一步恢复至接近70\%的水平。如果看中国工信部按载重吨统计的数据，这一数字更高，达到69\%。

这个V型走势并非政策奏效的证据，而是三种力量交织的结果，其中两种是暂时的或统计性的。第一种力量是政策驱动的订单分流。在2025年4月至10月美国贸易代表办公室（USTR）301条款停泊费从公布到实施之间的窗口期，船东为了规避不确定性，确实有动机将涉及美国贸易航线的订单转向非中国建造的船厂。这解释了韩国在2025年上半年订单激增的大部分原因。但NOM认为，这本质上是对政策不确定性的对冲，而非需求的持久再分配。

第二种力量，也是最被低估的因素，是周期性的结构性错配。2025年全球新船订单总量大幅下降约27\%，且收缩最严重的领域是油轮和散货船——这正是中国份额最高的船型。例如，2025年前十个月，成品油轮订单同比下降近85\%。当分母（全球总订单）在收缩且构成不利时，中国份额的下降在很大程度上是一个统计假象。一旦全球订单总量在2026年第一季度反弹（同比增长约40\%），且反弹主要由油轮和散货船引领，中国份额便迅速恢复。这意味着，将短期份额波动归因于竞争力的丧失，可能是一个严重的误判。

2025年11月，特朗普与习近平在釜山会晤后达成的12个月暂停期，为理解地缘政治对造船业的影响提供了一个近乎完美的自然实验。NOM的报告清晰地展示了这一机制：一旦暂停期确立，船东避免中国船厂的经济激励消失，中国份额在2025年末至2026年初迅速回升。

这一现象的含义是深远的。它意味着，2025年中期中国份额的低点，在很大程度上是市场情绪和订单组合共同作用下的一个“条件性”事件，而非对中国造船业成本、交付能力和技术水平的永久性否定。只要地缘政治紧张局势得到阶段性缓和，订单就会自然回流。这为投资者提供了一个重要的观察框架：中国造船业的核心竞争力并未因政策冲击而受损，其护城河依然稳固。真正的风险在于，2026年11月暂停期到期后，双方能否达成新的框架协议，但这更多是一个政治博弈问题，而非产业基本面问题。

如果说地缘政治和周期波动是短期扰动，那么绿色船舶转型则是决定中国造船业长期地位的结构性变量。NOM的数据显示，2026年第一季度，中国新接国际造船订单中，80.2\%为绿色燃料动力船舶（涵盖LNG、甲醇、LPG、乙烷及电力推进）。这并非简单的市场趋势跟随，而是中国在高附加值船型领域实现系统性突破的标志。

中国在这一转型中的领先地

[... middle omitted ...]

为全球油轮市场创造了一个强大的需求脉冲。这一事件的传导机制并非简单的货运量增加，而是通过吨海里（tonne-mile）的重新计算：中东原油被迫经由更长的替代航线运输，每艘船的单次航程时间延长；同时，被困在海峡内的数十艘油轮暂时退出运营船队。在船队运力结构性紧张（约20\%的原油油轮船龄超过20年）、运费同比近乎翻倍以及预期库存回补周期将维持高需求的背景下，油轮新船订单在2026年第一季度同比增长了三倍，占全球总订单量的32\%，为2017年第二季度以来最高。

这一轮订单激增的分配结果至关重要。根据中国工信部数据，中国船厂承接了超过90\%的超大型原油轮（VLCC）新船订单。这并非偶然，而是中国在散货船和油轮领域长期积累的成本和产能优势的集中体现。当突发事件催生紧急需求时，只有中国船厂拥有足够的船坞、熟练的劳动力以及具有竞争力的交付周期来承接这些订单。这进一步印证了NOM的核心观点：中国造船业的优势是系统性的，而非依赖于特定船型或市场环境。突发事件只会加速而非逆转这一趋势。

\subsection{[R023] NOM：工业AI的货币化，市场低估了2-3年过渡期的结构性摩擦}
{\small\color{refgray}归类：未被主题摘要引用}

NOM：工业AI的货币化，市场低估了2-3年过渡期的结构性摩擦

市场对工业AI的讨论，大多集中在“谁先落地”、“谁有模型”、“谁拿到订单”这些叙事层面。但NOM在2026年6月发布的最新研报中，提出了一个更值得深思的框架性判断：工业AI在短期内不是独立的增长引擎，而是附着在存量DCS（分布式控制系统）基础上的收入放大器。真正决定胜负的，不是AI模型的能力，而是谁能把AI嵌入到已经存在的工业控制基础设施中。

这份报告的核心洞察在于，它没有陷入“AI赋能一切”的乐观叙事，而是指出了市场普遍低估的一个关键变量——从当前状态到工业AI规模化落地，中间存在一个2到3年的“过渡痛苦期”。这不仅仅是技术成熟度的问题，而是涉及客户采购习惯、跨厂商兼容性、以及商业模式从项目制向订阅制切换的结构性障碍。

对于关注中国自动化赛道的投资者和产业决策者来说，这份报告提供了一个难得的冷思考视角：在工业AI的宏大叙事背后，真正的货币化路径远比想象中曲折。以下是我们从报告中提炼出的五个关键层次，每一个都指向一个尚未被充分定价的判断。

1. 工业AI的竞争，不是模型之争，而是存量工业基础设施的“附着权”

市场容易将工业AI与通用大模型混为一谈，认为谁的自然语言处理能力强、谁的参数规模大，谁就能在工业场景中胜出。NOM的调研给出了截然不同的结论：中国流程工业的AI堆栈建立在四个软件支柱之上——TPT（时序预训练Transformer）、UCS（通用控制系统）、AOP（自主运行工厂）和OMC（运营管理与控制系统），它们全部层叠在开源基础模型之上。这意味着，底层模型的同质化程度很高，真正的差异化来自垂直行业的工艺知识积累。

报告清晰地绘制了国内OT（运营技术）厂商的定位地图：中控技术锚定石化化工，和利时锚定轨交与电力，汇川技术则以新能源为核心。每一个厂商的护城河都不是模型参数，而是数十年积累的工艺know-how。这种“附着权”一旦建立，竞争对手很难通过AI技术突破来撬动客户关系。

这里有一个关键推论：工业AI的竞争本质上是存量DCS安装基础的争夺战。谁控制了DCS，谁就掌握了AI落地的入口。对于新进入者来说，即便AI技术再先进，如果没有与现有控制系统的深度绑定，也很难撬动客户的信任和采购预算。

NOM在报告中对比了两种截然不同的AI部署路径。中控技术的模式是典型的软件驱动：将其旗舰AOP与TPT和订阅制MaaS/RaaS捆绑，销售给石化客户。2026年第一季度，中控披露了1.84亿元人民币的工业AI收入，占其总收入的12\%。但NOM调研指出，其中至少20\%是捆绑的DCS硬件收入，这意味着纯粹的软件订阅收入可能比表面数字更少。

汇川的路径则完全不同。它不是卖软件，而是将AI和计算能力直接嵌入伺服驱动器、电机和专用PLC中，以“物理融合”的方式实现AI部署。在新能源和电池工厂，汇川的集成控制器-伺服解决方案搭配AI与实时力控，甚至直接服务于人形机器人供应链。NOM认为，汇川通过硬件拉动实现AI货币化，而非依赖订阅软件。

这种差异意味着什么？从商业模式的防御性来看，硬件的替换成本远高于软件。一旦汇川的伺服驱动器和PLC成为工厂自动化产线的物理组成部分，客户更换供应商的摩擦成本极高。而软件订阅模式虽然边际利润更高，但在当前阶段面临客户对数据安全、国产替代的抵触，以及“先试用再付款”的采购惯性。汇川的路径虽然看起来更“重”，但可能构建了更深的护城河。

3. 2-3年过渡期的核心障碍，不是技术，而是“协议不通”和“客户不信任”

NOM报告中最具洞察力的部分，是对过渡痛苦期的诊断。它列出了两个关键的结构性约束。

第一个是MaaS/RaaS订阅模式面临的阻力。国有企业的客户对数据上云和第三方模型服务存在天然的抵触，这种抵触不仅来自安全考量

[... middle omitted ...]

模化落地不仅仅是单个厂商的技术突破，而是整个工业通信生态的协调问题。在各方利益博弈尚未达成共识之前，项目级交付的摩擦成本将长期存在。

这两个约束叠加，使得NOM判断：工业AI的收入拐点不会在短期内出现，至少需要2到3年的过渡期。这期间，市场可能会看到AI相关的收入增长，但更多是存量DCS升级的附属品，而非独立的增量。

4. 报告尚未完全回答的关键问题：软件订阅的“粘性”能否在国企客户中真正建立？

NOM在报告中乐观地指出，工业AI软件的利润率（超过20\%）远超DCS不到10\%的利润率，一旦规模上量，将带动混合利润率扩张。同时，报告强调“订阅粘性，而非一次性项目，才是持久的拐点”。

但这里有一个尚未被充分验证的假设：在中国的国有企业客户群体中，软件订阅模式能否真正建立“粘性”？国企的采购流程通常以项目制为主，预算按年度审批，对持续性的服务费支出存在天然的抵触。即使AI软件能带来效率提升，客户的决策逻辑可能仍然是“这次买断，下次再说”，而非“长期订阅，持续优化”。

\subsection{[R024] NOM：人民币中间价定价模型正在释放比市场预期更复杂的信号}
{\small\color{refgray}归类：中国政策与金融：制度重塑进行时}

当多数市场参与者将注意力集中在人民币汇率的短期波动与政策干预时，一份来自NOM亚洲外汇策略团队的最新报告揭示了一个更值得关注的深层变化：人民币中间价的定价模型正在经历一次系统性的重新校准，而这种校准的幅度与方向，可能比任何一次单一的政策表态都更具长期含义。

这份于2026年6月8日发布的报告，核心在于其提供了一个可量化的中间价预测模型。模型显示，在不包含逆周期因子的情况下，当日预测值为6.7921，较前值大幅调低236个基点，同时较前一交易日官方收盘价高出209个基点。而纳入逆周期因子调整后，预测值变为6.8109，较前次定盘价仅低48个基点。

这组数字本身并不复杂，但它们叠加在一起，指向了一个关键的判断：市场对人民币汇率的定价逻辑，可能正在从“单边预期管理”转向“多目标动态平衡”。这不仅仅是技术层面的模型调整，更是政策框架演进的一个可观测窗口。

NOM报告中的第二张图表，展示了模型误差（不含逆周期因子调整）的长期趋势。从2025年初接近-1800个基点的巨大负向误差，到2026年4月已收窄至约600个基点。这条曲线的斜率变化，比任何绝对值都更有意义。

误差的快速收敛，意味着市场力量与官方定价之间的背离正在系统性缩小。2025年初，模型系统性地高估了人民币的贬值压力，实际中间价远强于模型预测。而到了2026年，两者的差距已大幅收窄。这背后可能的原因有二：一是市场本身的预期变得更加理性，二是官方定价机制对市场信号的响应更加灵敏。

对于投资者而言，这意味着过去依赖“中间价偏离度”来捕捉政策意图的策略，其有效性正在下降。当模型误差趋于零时，中间价本身所包含的信息量将显著增加，它不再是政策立场的扭曲信号，而更像是一个市场出清价格的近似值。这对所有人民币相关资产的定价基准，都是一个根本性的改变。

报告中的第一张图表，列出了对预测变化贡献最大的四种货币：欧元贡献67个基点，韩元34个基点，澳元27个基点，墨西哥比索11个基点。欧元的权重远超其他货币，这并非偶然。

这组数据清晰地表明，当前人民币中间价的定价，其外部驱动力已经从单一的中美双边关系，转向了更广泛的全球贸易加权篮子。欧元的高权重，反映的是中欧贸易体量以及欧元在全球支付体系中的地位。韩元的贡献则揭示了东亚供应链的深度绑定。

这意味着，分析人民币汇率不能再仅仅盯着美元指数。欧洲的货币政策、韩国的出口数据、澳大利亚的大宗商品价格，这些看似遥远的变量，现在都直接参与到人民币中间价的每日形成过程中。对于做跨资产配置的投资者，这提出了更高的要求：汇率分析框架必须从双边扩展至多边，从单一市场扩展至全球流动性网络。

这是报告中最具政策含义的一个细节。在不含逆周期因子的模型中，预测值为6.7921，较前值大幅调低。但加入逆周期因子后，预测值变为6.8109，仅较前次定盘价微降48个基点。两者的差距，从之前的236个基点缩小到了仅19个基点。

这传递了一个微妙但重要的信号：逆周期因子仍在发挥作用，但其干预力度已经显著下降。过去，逆周期因子是市场冲击的主要缓冲器，其调整幅度动辄数百个基点。现在，它更像是一个精细的微调工具，只在边际上平滑定价的波动。

这种变化，可能意味着监管层对市场自我稳定能力的信心在增强。也可能意味着，在更高维度的政策目标（如国际收支平衡、资本流动管理）面前，单纯的汇率稳定不再是压倒一切的优先事项。无论如何，逆周期因子的“淡出”，是人民币汇率形成机制市场化进程中的一个重要里程碑。投资者需要适应一个中间价波动性可能适度放大的新环境。

NOM报告提供了一份2026年下半年的中国关键事件日历。这份日历本身不是分析，但它为理解模型变化背后的政策逻辑提供了时间轴。

从7月的政治局会议，到10月的国庆长假与APEC深圳峰会，再到12月的中央经济工作会议与

[... middle omitted ...]

题：模型假设的稳定性与政策意图的边界

尽管NOM的模型提供了极具价值的量化视角，但它仍然留下了几个关键问题，这也是任何模型都无法完全回答的。

第一，模型预测的准确性高度依赖于其背后假设的稳定性。如果全球贸易格局、资本流动模式或地缘政治风险出现非线性变化，模型的预测误差可能重新扩大。报告展示了误差收敛的历史，但并未对未来可能出现的误差反弹给出足够的情景分析。

第二，逆周期因子的参数设定是不透明的。报告提供了包含和不包含逆周期因子的两种预测，但并未解释逆周期因子本身的调整规则。这意味着，在市场出现极端压力时，监管层仍有能力通过调整逆周期因子来改变定价方向，而这种调整的时机和幅度，是模型无法预测的。

第三，也是最核心的问题：模型所隐含的“均衡汇率”水平，是否与政策层的长期目标一致？模型输出的是一个技术性价格，但政策层是否愿意接受这个价格作为长期参考，仍然是一个开放性问题。尤其是在2026年下半年这个充满不确定性的时间窗口，政策意图的优先级可能随时高于模型的技术指引。

\subsection{[R025] NOM：人民币中间价模型正在释放一个被市场忽略的定价信号}
{\small\color{refgray}归类：中国政策与金融：制度重塑进行时}

市场对人民币汇率的讨论，大多集中在贸易摩擦、资本流动和央行意图上。但一份来自NOM亚洲外汇策略团队的模型报告，提供了一个更底层、更可量化的视角：人民币中间价的定价机制，正在发生结构性偏移，而这种偏移的幅度和方向，可能被多数交易员低估了。

这份报告的核心工具是一个基于一篮子货币的中间价预测模型。模型显示，在剔除逆周期因子后，最新的中间价预测值为 6.7812，较前一日的官方收盘价低了 32 个基点。更重要的是，这一预测值较上一期模型结果大幅下调了 386 个基点。这不仅仅是日常波动，而是一个值得关注的信号。

为什么现在重要？因为市场往往只关注最终公布的中间价，而忽略了模型本身所揭示的“无干预”状态下的均衡位置。当模型预测与官方定价之间的偏差持续扩大，或者模型本身出现趋势性偏移时，它往往意味着央行面临的定价压力正在积聚，也预示着未来政策调整的潜在方向。

NOM这份报告的价值，不在于给出了一个具体的数字，而在于提供了一个可跟踪的、透明的分析框架。它让投资者能够从“猜测央行意图”转向“观测模型压力”。

这份NOM模型的预测逻辑，并不是基于对中国经济基本面的直接判断，而是基于一篮子货币的加权变动。这恰恰是理解人民币中间价定价机制的关键：央行通过设定中间价，在管理预期和反映市场供求之间寻找平衡，而一篮子货币的波动是其中最重要的客观锚点。

报告中的图 1 清晰地展示了这一点。在贡献最大的四个货币中，韩元和欧元的权重贡献均为负的 13 个基点，而澳元为正的 3.5 个基点，俄罗斯卢布为负的 3.5 个基点。这种分布并非巧合。

韩元的大幅走弱，反映了全球半导体周期和韩国出口竞争力的波动。欧元的下行，则与欧洲经济疲软以及欧洲央行的政策路径有关。这两个货币的同步走弱，对人民币中间价形成了向下的牵引力。而澳元的正向贡献，则体现了大宗商品价格和澳大利亚经济相对韧性的支撑。

这个事实合在一起意味着：人民币中间价的定价压力，主要来自外部，而非内部。 当前人民币面临的压力，很大程度上是“被拖累”的，而非“主动贬值”。这为理解央行的政策空间提供了重要背景。如果外部压力是主导因素，那么央行干预的必要性和有效性，就会与外部环境的演变紧密挂钩。

2. 模型误差的长期趋势，暗示着定价机制可能正在经历一场“隐性校准”

NOM报告中的图 2 展示了剔除逆周期因子后的模型误差历史数据。这张图的信息量比表面看起来要大得多。

数据显示，在 2025 年 1 月，模型误差曾高达负的 1800 个基点。这意味着，当时的模型预测远低于实际公布的中间价，央行通过逆周期因子或其它方式，显著拉高了中间价，以稳定市场预期。此后，误差逐步收窄，在 2025 年 7 月附近趋近于零。但随后，误差再次扩大，在 2025 年 10 月达到负的 600 个基点，并在 2026 年初转为正的 600 个基点。

这个 V 型反转，是一个极其重要的信号。它表明，央行对中间价的“管理强度”并非一成不变，而是在周期性地调整。 当误差为负时，央行在“托市”；当误差为正时，央行在“放任”甚至“引导”贬值。

从 2025 年下半年到 2026 年初，模型误差从负值快速转为正值，这很可能意味着央行正在逐步减少对中间价的干预力度，允许其更真实地反映市场压力。这种“隐性校准”的过程，往往比一次性的、大幅度的政策调整更具研究价值。它反映了决策层对汇率弹性的容忍度正在发生变化，而这种变化，是市场参与者必须密切跟踪的。

图 3 展示了中间价的每日变化幅度。在 2025 年 11 月至 2026 年 5 月期间，单日变动幅度从 20 个基点逐步放大到 150 个基点。尤其是在 2026 年 2 月和 4 月，出现了连续的、方向性的较大变动。

这并非偶然。一个稳定的中间价，往往意味着央行在

[... middle omitted ...]

的一系列关键事件，将成为检验模型压力的“压力测试”

NOM报告中的图 4 列出了 2026 年下半年中国的一系列重要政治和经济事件。这些事件不是孤立的，它们共同构成了一个“压力测试”的窗口。

从 7 月底的政治局会议，到 10 月的国庆黄金周，再到 11 月的 APEC 会议和 12 月的中央经济工作会议，以及年底可能的中美领导人会晤。每一个事件，都可能成为市场情绪和汇率预期的催化剂。

报告没有完全展开的是，这些事件对中间价模型的影响路径。例如，APEC 会议期间，如果中美关系出现积极信号，可能会推动风险偏好改善，从而对韩元、澳元等货币形成支撑，进而通过模型传导至人民币。反之，如果会议未能达成任何共识，地缘政治风险上升，则可能加剧资本外流压力，迫使央行在中间价上做出回应。

*这些事件本身不是定价因素，但它们所引发的市场反应，会通过一篮子货币的波动，最终反映在模型预测中。** 投资者需要关注的，不是事件本身，而是事件发生后，模型预测相对于官方定价的偏差如何变化。

\subsection{[R026] GS：美国供应链的“新常态”比市场预期的更早到来}
{\small\color{refgray}归类：未被主题摘要引用}

市场观察者仍在等待新一轮供应链危机的爆发信号——关税冲击、地缘冲突、港口劳资谈判，每一个变量都被视为潜在的导火索。但GS最新的供应链拥堵追踪报告给出了一个值得认真对待的判断：美国供应链的拥堵水平已经回落至疫情前基线附近，并且这种状态可能比大多数人预期的更具持续性。

这份发布于6月初的报告显示，GS每周供应链拥堵指数在最近一周环比上升了7\%，但瓶颈规模评分仍维持在“2”的水平。这个数字意味着什么？在2021年底至2022年初的供应链危机高峰期，该评分曾达到“10”的顶峰。而当前“2”的水平，已经与2020年疫情前的流动性水平基本一致。更重要的是，GS5月的周度平均拥堵评分同样录得“2.0”，这一结果不仅远低于2021年12月至2022年1月的峰值，而且与疫情前基线几乎完全对齐。

GS在报告中明确指出，如果供应链压力继续缓解，该指数有望在2026年更持续地停留在“1”的区间。这绝不是一句轻描淡写的预测。它意味着，过去三年驱动全球物流成本飙升、港口拥堵、运价暴涨的核心变量，正在从“供给不足”转向“需求正常化”。

对于投资者和企业决策者来说，这一判断的含金量在于：它不是一个孤立的港口数据改善，而是整个物流体系的系统性恢复。从GS追踪的指标来看，西海岸集装箱船等待靠泊数量已经连续数周保持在“1”的水平，东海岸虽然从4上升至5，但相比2021年动辄数十艘的积压，已是天壤之别。铁路多式联运方面，西海岸一级铁路（联合太平洋和伯灵顿北方圣达菲）的周度平均运量增速从上周的+14\%大幅放缓至-2\%，这并非需求崩塌，而是高基数效应和正常化进程的自然体现。

这些数据合在一起意味着：过去两年支撑海运、仓储、卡车运输等环节定价权的“稀缺溢价”正在消失。企业不能再依赖供应链紧张带来的被动议价能力，而必须主动将规模优势转化为真正的成本结构和客户关系壁垒。

GS报告中最值得关注的细节之一，是东西海岸的差异。西海岸的集装箱船积压量已经稳定在“1”附近，几乎可以忽略不计。但东海岸的积压量从4上升至5，虽然绝对值不高，却是一个连续上升的趋势。

这个差异背后有两个相互交织的驱动因素。第一，东海岸和墨西哥湾沿岸的港口正在吸收原本通过西海岸进口的货物，这是过去几年供应链多元化战略的直接结果。第二，巴拿马运河的水位限制和劳资谈判的不确定性，正在促使部分货主将订单前置或调整航线。

从GS追踪的月度指标来看，4月“三大港”（洛杉矶、长滩、奥克兰）的进口重箱量同比增速从3月的+11.6\%降至+9.5\%，而5月的数据已经转为-10\%。这种剧烈波动不是需求崩溃，而是前置订单的消化和库存调整的结果。对于投资者而言，这意味着运输类股票的盈利可预测性正在下降，季度间的波动可能成为常态而非例外。

GS报告显示，5月西海岸一级铁路的多式联运平均运量同比增长约6\%，但周度数据波动剧烈。BNSF的周度运量从+18.7\%骤降至+3.3\%，UNP则从+8.4\%降至-7\%。这种大起大落很容易被解读为需求疲软，但GS的数据分析框架提供了另一种解释。

从更长的周期看，2024年铁路多式联运经历了强劲的增长，部分原因是2023年的低基数效应。进入2025年，基数效应消退，加上零售商和制造商已经完成了库存重建，运量增速自然回归到更温和的水平。GS的月度拥堵评分在2023年1月至2025年5月期间几乎稳定在1.8，这本身就说明，当前的状态不是“好转”，而是“稳定”。

对于投资者，这里的关键洞察是：铁路公司的盈利增长动力正在从“量”转向“价”和“效率”。UNP的系统中转停留时间最近一周为19.8小时，BNSF为22.6小时，均处于可控区间。这意味着铁路公司可以通过精细化运营来提升利润率，而不是依赖运量增长。那些能够持续降低中转时间、提高资产周转率的公司，将在下一阶段获得估值溢价。


[... middle omitted ...]

转正，既反映了2024年低基数的影响，也暗示着当前的有效运力供给已经不再过剩。船公司通过减速航行、拆船和闲置运力等手段，正在主动管理供给。

这个判断如果成立，意味着海运运价的下行风险已经显著减小。对于依赖海运的进口商和零售商，这不是好消息；但对于班轮公司和物流服务商，这可能是利润率稳定的信号。值得继续跟踪的是，GS报告中的运价数据是基于Freightos的即期运价，而合同运价的变化通常滞后2-3个月。如果即期运价持续正增长，合同运价在2025年下半年可能面临上行压力。

GS在报告中坦诚地提出了一个核心问题：“关税和地缘冲突将对货运需求和时间安排产生什么影响？”这是整份报告最大的不确定性来源。

报告中的拥堵指数和瓶颈评分都是基于高频的物理指标——船舶等待数量、铁路运量、底盘停留时间等。这些指标反映的是已经发生的物流活动，而不是未来的订单意愿。关税政策的变化可以通过两个渠道影响供应链：一是前置订单效应（在关税生效前抢运），二是需求抑制效应（关税推高成本后减少进口）。

\subsection{[R027] JPM：香港地产市场真正的考验不是政策阴云，而是数据能否持续验证复苏}
{\small\color{refgray}归类：地产与消费：制度优化与结构分化}

JPM：香港地产市场真正的考验不是政策阴云，而是数据能否持续验证复苏

过去两周，香港地产股跑输恒生指数7\%。一个旧隐忧尚未消散——内地资本外流管控可能收紧；一个新隐忧又浮出水面——美国强劲的就业数据重新点燃了加息预期。两个“悬顶之剑”同时压境，市场情绪迅速转冷。

但JPM在这份最新研报中给出了一个值得细品的判断：这两个隐忧本身，都不足以推翻香港住宅市场的复苏趋势。房价年初至今已上涨9.6\%，该机构维持2026财年全年10\%-15\%的涨幅预测。真正值得关注的，不是这些担忧会不会发生，而是市场需要什么样的证据才能重新建立信心。

这份报告的价值不在于罗列风险，而在于它提供了一个清晰的框架：在不确定性中，哪些变量真正决定方向，哪些公司会在震荡中相对坚挺，以及投资者应该在什么位置重新入场。

第一个隐忧是内地资本外流管控可能收紧。市场担心这会切断来自内地买家的需求。JPM在5月底和6月初的两份报告中已经详细拆解过这个问题：截至目前，没有官方禁令禁止内地居民在香港购房；资本外流管控并非新政策，历史上已有先例；即使在最极端情景下——所有“基于内地的内地买家”完全退出市场——他们只占成交量的5\%-10\%，按金额计算也仅为10\%-15\%。更重要的是，报告中区分了“基于内地的内地买家”和“基于香港的内地买家”，后者受到的影响应该有限。香港楼市的绝大多数买家仍然是本地居民。

第二个隐忧是加息预期重燃。美国劳动力数据超预期后，市场开始重新定价利率路径。JPM目前维持利率暂停至2027年一季度的预测。但如果美国真加息，香港最优惠利率可能跟随上调，按揭利率也会走高。不过报告指出一个关键细节：如果利率维持在现有水平，香港购房者仍然享有边际正carry——名义按揭利率3.25\%，扣除现金回赠后的实际按揭利率约2.9\%-3.0\%，而净租金收益率为3.0\%。这个正利差虽然微薄，但方向是积极的。

这两个隐忧的共同特点是：都还没有实质化。但市场已经提前做出了反应。JPM认为，这种担忧可能只会因为持续强劲的数据而消退——比如房价、新盘销售率、二手成交量等关键指标继续保持韧性。

市场对内地买家退场的担忧，部分源于数据误读。报告提供了两组关键数据来校正认知。

第一组来自税务局：在2024/25财年，非香港身份证持有者的个人买家占成交量5.5\%，按金额计占7.2\%。这个比例虽然较前几年有明显上升——2020/21财年仅为0.1\%——但绝对占比仍然有限。

第二组来自中原地产，以买家姓氏拼音定义“内地买家”，口径更宽。按这个标准，2026年第一季度，“内地买家”在整体市场占比约23\%，在初级市场更高达36\%。但报告特意指出，这个口径无法区分买家当前居住地或身份，因此长期居住在香港的内地背景人士，甚至姓氏拼音为普通话拼音的本地居民，都被归入此类。

真正值得关注的是区域集中度。报告列出了自2024年3月“撤辣”以来，内地买家最集中的五大区域：启德、黄竹坑、长沙湾/深水埗、将军澳、中西区。启德以3,050宗交易遥遥领先，其中初级市场占2,786宗。这意味着，如果资本管控真的收紧，启德高端住宅项目可能首当其冲。而启德跑道区的未售单位中，中国海外发展占21\%，恒基兆业占13\%，会德丰和长实集团各占12\%-13\%。这些开发商的短期销售节奏可能因此承压。

这里的含义很清楚：内地买家退场不是系统性风险，但对特定区域和特定开发商而言，这是一个需要密切跟踪的尾部风险。

3. 加息影响被过度渲染，历史经验表明利率上升不等于楼市下跌

市场对加息的恐惧，部分源于对历史经验的误读。JPM在报告中明确指出一个反直觉的事实： 年和 年，香港都经历了加息周期，但房价和地产股都同步上涨。

背后的逻辑并不复杂。加息通常发生在经济过热、就业强劲、收入增长的背景下。此时居民的购

[... middle omitted ...]

化，这会在一定程度上缓冲对当期利润表的冲击。报告没有完全展开这个细节，但它意味着实际影响可能比数字显示的温和。

当前香港地产板块的交易价格相对于资产净值折让47\%，接近历史均值下方一个标准差的位置。按股息率计算，板块收益率为4.4\%，高于历史均值0.5个标准差。这两个指标共同指向一个判断：估值不算便宜，但也没有到泡沫程度。

更重要的是，JPM认为只要房价保持稳定，板块出现大幅估值下修的可能性不大。但需要警惕的是，地产股年初至今已经上涨12\%，而恒生指数同期下跌3\%。这个相对表现意味着部分乐观预期已经被定价。报告特别指出，即使不考虑新增隐忧，板块下行空间也控制在10\%以内。

对于长期投资者，报告在回调时建议关注的标的有四个：新鸿基地产（未来2-3年每股盈利和股息有望实现中个位数增长）、信和置业（5\%股息率且派息确定性高）、太古地产（内地零售持续改善，香港写字楼触底，5.6\%股息率加5\%年化股息增长）、香港置地（内地租户销售有望超预期，6.6\%高股息率且确定性高）。

\subsection{[R028] Bernstein：Medicare Advantage的利润率正在触底，但市场低估了“公用事业化”的长期定价逻辑}
{\small\color{refgray}归类：医疗与创新药：化疗增敏价值重估}

Bernstein：Medicare Advantage的利润率正在触底，但市场低估了“公用事业化”的长期定价逻辑

市场对Medicare Advantage（MA）的定价，长期以来被一个经典叙事主导：这是一个高增长、高利润、且受政策风险扰动的周期性赛道。投资者习惯于将MA计划视为保险周期的放大器——费率高时利润爆发，费率低时利润崩塌。这份Bernstein的暑期教学系列报告，提供了一个比周期判断更值得关注的框架：MA正在经历结构性转型，其终局不是利润率的均值回归，而是向“公用事业型”市场的收敛。这意味着，当前市场对MA相关股票的估值，可能既低估了利润率修复的幅度，也误读了长期稳态利润率的水平。

报告的核心判断是，MA行业利润率已经在 年触底，预计2026年将回升约1个百分点，此后每年以约0.5个百分点的速度温和扩张，直到2030年。这一判断的基础，并非政策红利或需求爆发，而是竞争退出、费率改善和行业定价纪律的回归。但比短期复苏更重要的，是Bernstein提出的长期图景：MA和Medicaid市场都将趋于“公用事业化”，最终由4-5家规模化玩家主导，MA的稳态利润率维持在2-3\%之间。如果这一判断成立，那么当前市场对头部MA企业的估值折扣，可能恰恰提供了逆向布局的机会。

理解MA利润率为何在 年遭遇重创，是判断后续走势的前提。Bernstein的分析提供了一个清晰的周期框架： 年间，MA市场收入增速高达15-20\%，利润率稳定在5\%左右。这一高回报吸引了大量资本涌入，竞争加剧直接压低了利润率，也为后续的费率冲击埋下了伏笔。

当CMS在2024年和2025年连续给出低于2.5\%的有效增长率时，行业利润率的崩溃几乎不可避免。这并非MA商业模式本身的结构性缺陷，而是典型的保险周期——高利润吸引竞争，竞争推高成本，成本叠加费率收紧导致危机。Bernstein指出， 年的有效增长率实际上低估了医疗成本趋势，这进一步放大了利润压力。

但关键在于，这个周期已经进入出清阶段。竞争者在低利润环境下被迫退出，行业集中度提升，头部企业开始将定价权回归到以盈利为导向。从这个角度看，利润率危机恰恰是市场出清的必要代价，而非行业衰落的信号。对于投资者而言，这意味着当前的低利润时期可能正是布局的窗口，而不是需要回避的雷区。

2. 2026年利润率修复的驱动力是结构性的，而非政策偶发性

Bernstein预计2026年行业MA利润率将扩张约1个百分点，随后每年以0.5个百分点的速度增长。这个判断令人关注，因为它并非建立在某个单一政策利好的假设上，而是基于三个相互支撑的结构性因素。

首先，竞争退出正在实质性改善定价环境。报告指出，在利润率压缩期间，大量中小型参与者已经或正在退出市场。头部企业如UNH、HUM、CVS的市场份额合计约55\%，但市场整体仍相对分散。随着尾部玩家离场，剩余的规模化企业将获得更强的议价能力，这直接反映在2026年的定价上。

其次，费率环境正在从低点改善。CMS公布的2026年MA有效增长率高达9.04\%，这显著高于2024年的2.28\%和2025年的2.33\%。虽然这一数字包含了技术性调整，但它至少表明政策制定者意识到了费率过低对行业可持续性的威胁。Bernstein认为，这标志着费率周期的底部已经出现。

第三，也是最重要的，行业正在从“市场份额优先”转向“定价盈利优先”。这一转变并非自愿，而是在利润压力下的被迫选择。但一旦形成共识，它对利润率的支撑将是持续的。报告特别指出，MA计划在2024年向成员提供了平均每月194美元的额外福利，这相当于15\%的额外价值。在利润压力下，企业将不得不重新审视福利设计的成本效益，这本身就是利润率修复的微观基础。

3. “公用事业化”终局意味着利润率上限被锁定，但

[... middle omitted ...]

LP-1s可能贡献每年额外1年的生命预期增长），总市场规模仍在扩大。更重要的是，在寡头格局下，企业将不再需要通过价格战争夺市场份额，利润率将趋于稳定。

2-3\%的利润率看起来不高，但考虑到MA业务的规模效应和资本轻属性，这个利润率水平对应的资本回报率可能相当可观。更重要的是，公用事业化的估值逻辑不同：投资者愿意为稳定的现金流和可预测的回报支付溢价，而不是为高增长但波动巨大的利润支付折价。如果市场逐渐接受这一终局假设，当前处于历史低位的估值倍数可能面临重估。

当然，这个判断仍需要验证。报告没有完全展开的是，公用事业化过程中，哪些企业能够成为最终的4-5家赢家。规模是必要条件，但并非充分条件。能否将规模转化为对供应商的议价权、对CMS的谈判能力、以及对成员留存率的控制力，才是决定性因素。

Bernstein识别了三个关键颠覆因素：GLP-1药物、人工智能、以及基于价值的护理（VBC）。这些因素将如何影响MA的利润率和竞争格局，报告提供了初步框架，但许多问题仍待解答。

\subsection{[R029] 摩根斯坦利：ASCO数据揭示中国创新药进入“化疗增敏”价值重估期}
{\small\color{refgray}归类：医疗与创新药：化疗增敏价值重估}

摩根斯坦利：ASCO数据揭示中国创新药进入“化疗增敏”价值重估期

这不是又一轮“me-too”分子的数据更新。摩根斯坦利在ASCO 2026最后两天的核心观察指向一个更具结构性的判断：中国生物科技公司的竞争逻辑，正在从“能否做出同类最佳”转向“能否在特定化疗骨架之上，做出可证明的、有临床意义的增量”。这份报告最值得关注的，不是AK112或mesutoclax的单组数据有多漂亮，而是这些数据背后，一个被低估的估值驱动因子正在浮现——化疗联合疗法的风险收益比重构。

市场此前对中国创新药的定价，高度集中在PD-1/VEGF双抗这类大单品预期上。但摩根斯坦利分析师Jack Lin团队在ASCO现场捕捉到的信号是：真正可能改变竞争格局的，不是单一靶点的突破，而是“化疗增敏剂”这一临床策略的系统性验证。如果这一策略被证实，将直接改写多个大适应症的标准治疗路径，并对现有IO（免疫肿瘤）疗法的市场地位形成实质性冲击。

1. AK112在SCLC的数据挑战了IO联合化疗的既有天花板，但样本构成的差异仍是关键变量

摩根斯坦利重点分析了AK112联合脂质体伊立替康（lipo-iri）在2L SCLC中的数据。ORR 61.7\%、mPFS 8.1个月、6个月PFS率69.1\%，这些数字在与BNT327联合紫杉醇的IO经治人群对比中显得突出——后者ORR仅37.2\%、mPFS 5.4个月、6个月PFS率46.1\%。更重要的是，≥3级TRAE发生率AK112组为36.7\%，显著低于对照组的78.6\%。

但报告没有止步于“AK112更好”这个表层结论。分析师做了一个关键的数据拆解：单独使用lipo-iri的RESILIENT研究中，ORR为44.1\%、mPFS为4.0个月。这意味着化疗本身已经设定了响应率的“地板”，AK112的核心价值不在于拉高响应率，而在于延长响应的持久性——尤其是在CFI Personal reading notes and learning share only. Not investment advice.

*AK112组合：ORR 61.7\%，mPFS 8.1个月，6个月PFS率69.1\%** *对比同类方案：ORR 37.2\%，mPFS 5.4个月，6个月PFS率46.1\%**

数据差异很明显，但有个细节值得注意：AK112组里CFI＜90天的难治患者比例是36.7\%，而对照方案是50\%。这个基线差异可能会影响横向对比的可靠性。

脂质体伊立替康单药的历史数据ORR 44.1\%、mPFS 4.0个月，说明化疗本身建立了疗效底线，AK112的关键加分项是持久性——即使在化疗耐药人群里，mPFS也能到6.7个月。

不过出血事件发生率20\%（≥3级未单独披露），QoL数据也需要持续跟踪，这是风险收益比的核心变量。

HR-MDS初治队列（n=10）：ORR 100\%，IWG23 compCR 90\% 对比同类BCL-2抑制剂：ORR 64-74\%，compCR约70\%

AML初治队列（n=44）：cCR 81.8\%，MRD阴性率86.5\% 对比同类：CR/CRi 67.1\%，MRD阴性35.4\%，≥3级感染50.6\%

Mesutoclax的疗效信号确实强，但样本量小是硬伤。MDS需要等VERONA研究后的OS/持久性数据，AML要看感染和血象恢复情况。BCL-2赛道里，疗效和安全性的平衡才是最终决胜点。

两个靶点都在验证同一个逻辑：精准人群+差异化机制=临床价值。但数据越漂亮，后续大样本验证的压力也越大。

China Biotechnology | Asia Pacific

ASCO 2026 Takeaways - Days 4-5

Vs BNT327+paclitaxel IO-exp

[... middle omitted ...]

3\% Equal-weight/Hold 1571 43\% 369 40\% 23\% 723 44\% Not-Rated/Hold 3 0\% 0 0\% 0\% 1 0\% Underweight/Sell 551 15\% 86 9\% 16\% 201 12\% Total 3,

Data include common stock and ADRs currently assigned ratings. Investment Banking Clients are companies from whom MS received investment banking compensation in the last 12 months. Due to rounding off of decimals, the percentages provided in the "\% of total" column may not add up to exactly 100 percent.

\subsection{[R030] Citi：AI供应链的瓶颈正在从GPU转移到PCB和铜箔基板}
{\small\color{refgray}归类：AI/算力：资本开支周期扩散，供应链瓶颈转移}

Citi：AI供应链的瓶颈正在从GPU转移到PCB和铜箔基板

市场对AI硬件产业链的关注，长期以来集中在GPU供给、CoWoS封装产能、以及HBM存储器的供需缺口上。但一份来自Citi近期台湾科技大会的调研纪要，揭示了一个正在发生的结构性变化：AI服务器产业链的瓶颈，正在从先进封装和晶圆制造，向下游的PCB和铜箔基板（CCL）环节转移。 这不是一个短期扰动，而是供给端投资节奏长期错配的结果。

Citi分析师Jack Chen及其团队在2026年6月初的台湾科技大会上，与三家核心PCB和CCL厂商——金像电子（GCE）、台光电（EMC）和联茂电子（TUC）进行了深度交流。三家公司的反馈高度一致：AI需求强劲，产能扩张速度跟不上客户需求，涨价正在成为行业共识。这份报告的价值，在于它不只是给出了一个“景气向好”的判断，而是系统性地拆解了三个不同环节（PCB、CCL、上游材料）各自的供需格局、定价策略和竞争分化。

Citi报告中最具结构性意义的判断，来自于EMC和TUC两家CCL厂商的共同观察：CCL行业的产能扩张速度，系统性落后于PCB行业的扩张速度。这意味着，即便PCB厂有能力建新厂、扩产线，它们也可能面临上游CCL供应不足的瓶颈。

这个判断的逻辑链条是清晰的。PCB厂商如金像电子计划在2027、2028、2029年每年新建一座工厂，扩产意愿强烈。但CCL厂商的扩产节奏更为谨慎。EMC明确表示其未来产能扩张将主要集中在中国，因为“产能部署更高效、客户需求更充足”。TUC的泰国新厂则因设备交付延迟，营收贡献时间点从预期推迟到了2026年三季度末。

这两组信息合在一起，指向一个结论：在AI服务器的需求爆发周期中，PCB和CCL之间的产能扩张速度差，将导致CCL成为新的供给瓶颈。Citi报告原文的表述是“CCL capacity ramp would be lower than new PCB capacity ramp, resulting in CCL shortage and potential price hikes”。这不是猜测，而是两家CCL龙头在科技大会上给出的明确指引。

对于投资者而言，这意味着需要重新审视AI硬件产业链的“瓶颈”分布。过去两年，市场习惯性地将焦点放在GPU和先进封装上。但Citi这份调研提示，当GPU供给逐步缓解、CoWoS产能持续开出之后，下游的PCB和CCL环节可能成为新的约束条件，进而影响AI服务器的实际出货量。

Citi报告中最具说服力的部分，是三家公司在“涨价”问题上的表态高度一致。EMC和TUC都明确表示，将根据通胀或客户需求的溢价程度，持续上调CCL价格。金像电子则透露，其客户目前已经接受了PCB的涨价提议。

这意味着，涨价不是一家公司的个别行为，而是行业性的供需再定价。Citi报告给出的背景是：CCL供应紧张将“持续”，支持进一步的涨价。EMC的策略是继续向下游传递原材料成本上升的压力，同时强调在调价过程中维持与关键AI客户的长期战略关系。TUC则采取了更为激进的定价策略——对M7及以下级别的CCL产品主动提价，对M8级别则与EMC保持协同。

这里值得注意的细节是，两家CCL公司在定价策略上出现了微妙的分化。TUC对低规格产品采取了更激进的涨价姿态，而EMC更强调“长期关系”和“客户黏性”。这种分化可能反映了它们在不同客户和不同产品线上的竞争定位差异。对于投资者来说，需要关注的是：当行业整体涨价时，哪家公司的议价能力更强、涨价空间更大，而不是简单地将所有CCL公司视为同一类资产。

金像电子作为PCB端，其涨价逻辑略有不同。它拥有长期采购协议的CCL供应保障，但临时加单几乎无法被满足。这意味着，金像电子的产能扩张和营收增长，将受到其供应链管理能力的制约，

[... middle omitted ...]

提升、以及新客户导入带来的增长。

但这里有一个需要验证的问题：金像电子的市场份额提升，是来自竞争对手的失误，还是自身技术或产能优势的体现？报告没有完全展开。在PCB行业，客户导入周期长、认证门槛高，一旦进入供应链，切换成本也高。如果金像电子的份额提升是可持续的，那么它在AI ASIC PCB领域的护城河就在加深。反之，如果只是短期订单分配的结果，则需要警惕竞争格局的变化。

4. 上游材料的分化正在加剧：玻璃纤维、铜箔、PTFE的供给判断不一

Citi报告在讨论上游材料时，呈现出一种有趣的“分裂”状态。EMC和TUC对玻璃纤维、铜箔的供给判断存在明显分歧。

在玻璃纤维方面，EMC认为其不太可能受到Low-Dk2短缺的影响，而TUC则认为Low-Dk2材料将持续紧张。两家公司都认为E-glass的供给状况正在改善，2027年可能不再是主要瓶颈。在铜箔方面，EMC预计2026年下半年至2027年可能出现供给紧张，而TUC则认为不存在明显的短缺风险，铜箔价格保持稳定。

\subsection{[R031] Bernstein：市场低估了支付网络的“量价分离”能力}
{\small\color{refgray}归类：科技与软件：AI分选机效应加速分化}

全球支付市场正在进入一个微妙的分水岭。过去十年，现金数字化是行业增长的核心叙事，卡支付渗透率从47\%攀升至64\%，Visa和Mastercard的股价也在这条曲线上获得了丰厚回报。但一个尖锐的问题正在浮现：当美国卡渗透率已达72\%，当亚洲市场增长放缓，当稳定币和即时支付不断制造噪音——现金数字化的跑道究竟还剩多少？

Bernstein刚刚发布了第12版全球支付预测模型。这份报告的结论，比表面看起来要复杂得多。它没有简单地回答“增长还有多少”，而是揭示了一个更关键的判断：支付网络的价值增长，正在从“量”的扩张转向“质”的变现。 那些只盯着卡交易量增速放缓的投资者，很可能低估了网络效应在交易结构、增值服务和新兴流量的二次放大作用。

我们仔细拆解了这份长达数十页的研报，提炼出五个核心洞察。它们不仅关乎Visa和Mastercard的投资逻辑，更指向整个支付产业链的定价权迁移。

Bernstein预测， 年全球卡交易量（C2B purchase volume）的复合年增长率约为8\%（恒定汇率），低于 年的9\%和 年的11\%。减速的主要驱动力来自两个方向：美国市场的高渗透率（72\%）导致增量空间收窄，以及亚太地区因本土支付生态的崛起而出现结构性放缓。

但这里有一个容易被忽视的细节。当衡量标准从交易金额切换到交易笔数时，增长曲线明显更陡峭：Bernstein预测全球交易笔数在 年的复合增长率约为9\%，仅比 年的10\%略有下降。原因在于，卡渗透率在笔数维度上仍然偏低——很多低客单价的场景（如便利店、交通、自动售货机）尚未完全数字化，而这些场景正是接触式支付和未来代理型商务的天然土壤。

对Visa和Mastercard而言，这意味着一层重要的收入缓冲。根据Bernstein的估算，两家网络超过三分之一的毛收入直接与交易笔数挂钩，而笔数的增长通常快于金额增长。当市场过度聚焦于“量的减速”时，这层结构性的韧性往往被低估。

美国是Bernstein预测中最受关注的市场，也是争议最大的区域。2025年美国卡渗透率约为72\%，Bernstein预计到2030年也仅会小幅升至74\%。这意味着，未来五年的增长将几乎完全依赖个人消费支出（PCE）的增长，而非渗透率的提升。

但这并不意味着美国市场失去了吸引力。Bernstein明确指出，即便在高度渗透的市场，双网络（Visa和Mastercard）历史上仍能实现双位数的收入增长。秘密在于：收入来源正在从“交易量的线性增长”转向“交易结构的非线性变现”。

具体而言，三个杠杆正在发挥作用。第一，代币化（tokens）正在改变每笔交易的价值分配，网络可以基于代币的发行和管理收取额外费用。第二，增值服务（VAS）的渗透率持续提升，这些服务包括风险分析、数据洞察、身份验证等，其利润率往往高于核心交换费。第三，新流量（new flows）如B2B支付、跨境支付和商业支付，正在开辟传统C2B交易之外的增量池。

Bernstein的模型显示，即便卡交易量增速放缓至5-6\%，叠加交易笔数9\%的增长和上述结构性杠杆，Visa和Mastercard仍有可能实现双位数的收入增长。这是一种典型的“量价分离”——收入增长正在脱离交易量的线性束缚。

3. 欧洲正在成为下一轮卡渗透率提升的主战场，但格局正在重塑

如果说美国是存量博弈，那么欧洲就是增量战场。Bernstein的数据显示，欧洲的卡渗透率从2019年的44\%跃升至2025年的76\%，预计到2030年将达到92\%。这是一个惊人的跃迁——五年内渗透率几乎翻倍，背后是欧盟推动的数字支付政策、接触式支付的普及以及消费者习惯的加速迁移。

Bernstein对欧洲的预测尤为积极： 年卡交易量复合增长率预计为11\%（恒定汇率），远高于全球平均的8\%。增长

[... middle omitted ...]


过去几年，账户对账户（A2A）支付、本地支付方式和稳定币不断制造新闻，市场对卡网络的替代风险始终存在担忧。Bernstein的模型提供了一个更冷静的视角。

报告指出，即便在那些A2A和本地支付取得成功的市场（如印度的UPI、巴西的Pix），卡交易量仍然保持高个位数到低双位数的增长。这背后的逻辑是：替代支付方式更多是在蚕食现金和支票的份额，而非直接替代卡支付。在印度，UPI的崛起确实挤压了借记卡的使用，但信用卡交易量仍在快速增长。消费者在数字支付生态中的行为是“多钱包并存”，而非非此即彼。

至于稳定币，Bernstein的观点更加明确：在五年的时间维度上，稳定币在零售消费者支付领域仍是一个“寻找问题的解决方案”。尽管稳定币在跨境汇款和美元需求旺盛的市场找到了实用场景，但在日常零售支付中，其用户体验、监管不确定性和商户接受度仍然构成重大障碍。Bernstein在另一份专题报告中对此有更详细的拆解，但核心结论是：稳定币对卡网络的威胁，至少在零售层面，被市场过度定价了。

\subsection{[R032] 摩根斯坦利：亚洲市场真正的超额收益来自结构性的“非共识定价”}
{\small\color{refgray}归类：区域市场：亚洲结构性机会与分化}

摩根斯坦利：亚洲市场真正的超额收益来自结构性的“非共识定价”

当多数投资者仍在为宏观数据的短期波动焦虑时，一份来自摩根斯坦利的最新研报揭示了一个更值得关注的事实：在过去一年里，其“三个可操作想法”组合累计获得了超过9,300个基点的相对超额收益，平均12个月总回报率达到16.0\%，相对基准的超额回报为10.6\%。这些数字并非来自某个单一主题的押注，而是来自对亚洲市场三类结构性力量的持续识别。

这份报告的核心判断是：当前亚洲市场的超额收益机会，正从依赖宏观贝塔转向寻找微观层面的“非共识定价”。无论是日本光通信与水冷模块的爆发、澳大利亚住房信贷周期的拐点，还是印度能源安全政策的重塑，这些机会的共同特征在于——市场共识尚未充分定价供给端的结构性变化。

以下是我们从这份研报中提炼的三个关键洞察，以及它们对投资者观察框架的启示。

1. 古河电工的利润新高路径揭示了一个被低估的产业逻辑：光通信与水冷模块的“双轮驱动”正在改写日本电子材料公司的盈利天花板

摩根斯坦利对古河电工（5801.T）给出“超配”评级，并将其目标价上调至66,000日元。支撑这一判断的核心逻辑并非公司现有的光纤业务复苏，而是两个正在加速的结构性变量：光通信产品的需求升级，以及水冷模块在数据中心散热领域的渗透加速。

报告明确指出，古河电工有望实现“新的创纪录利润且增长强劲”。这意味着，市场此前对这家公司的估值框架仍停留在传统光纤周期上，而忽略了其在AI算力基础设施配套中的新角色。水冷模块作为高功率数据中心散热的关键组件，其需求增速正在显著超越传统风冷方案。如果这一趋势持续，古河电工的盈利中枢将上移至历史高位之上。

这一判断的隐含含义更为重要：在亚洲电子材料领域，那些能够同时受益于“光通信升级”与“散热技术迭代”两个结构性趋势的公司，其盈利弹性和估值体系可能正在被系统性低估。投资者需要重新审视日本电子材料板块的定价模型，不能再用过去的周期框架来判断未来的利润高点。

2. 西太平洋银行的下调揭示了澳大利亚住房信贷市场的“空气口袋”：市场的乐观预期与信贷基本面的实际走弱之间存在显著裂口

摩根斯坦利对西太平洋银行（WBC.AX）给出“低配”评级，并下调了盈利预测和目标价。研报将这一调整的背景概括为“澳大利亚住房和抵押贷款市场前景的转变”，并维持对该行业的“谨慎”看法。

报告使用了“Air Pocket”一词来描述当前市场的状态。这是一个非常精准的比喻——航空领域的“空气口袋”意味着飞行器在看似平稳的航道上突然遭遇失速。映射到信贷市场，这意味着尽管澳大利亚经济数据尚未全面恶化，但住房信贷的微观信号已经出现了明显的减速迹象。银行板块的估值可能尚未充分反映这一信贷周期的拐点。

这一判断对投资者的启示是：在利率环境看似趋于稳定的背景下，市场容易忽视信贷需求端的内生性放缓。澳大利亚银行股的估值能否维持，取决于市场是否愿意重新定价其资产质量的潜在风险。而摩根斯坦利的选择表明，该机构认为当前的价格并未完全反映这一风险。

3. 印度斯坦石油的“超配”评级背后是能源安全逻辑的范式转换：政策制定者正在以超出市场共识的方式行动，印度燃料零售商成为能源供应改善的最大杠杆

摩根斯坦利将印度斯坦石油（HPCL.NS）列为印度燃料零售商中的首选标的，核心判断是“能源安全至关重要，政策制定者正在以令市场共识意外的方向行动”。报告进一步指出，印度燃料零售商是“能源供应改善超出石油冲击”背景下最具杠杆效应的标的。

这一判断的关键在于，市场此前对印度能源板块的定价主要基于“石油冲击”的负面情景，即油价上涨会压缩零售利润。但摩根斯坦利认为，政策制定者正在主动重塑能源供应的格局，其行动力度和方向超出了市场预期。如果印度能够通过政策手段改善能源供应的稳定性和成本结构，那么燃料零

[... middle omitted ...]

这些“非共识定价”的可持续性如何？摩根斯坦利的“三个可操作想法”组合在过去一年取得了显著的超额收益，但历史数据显示其命中率约为54\%。这意味着，接近一半的投资想法并未实现正回报。

这引出了一个更深层的追问：当前这三个判断（古河电工的利润新高、西太平洋银行的信贷风险、印度斯坦石油的政策红利）中，哪些最有可能从“非共识”演变为“共识”，从而驱动股价的持续重估？哪些可能因为关键假设的失效而无法兑现？

例如，古河电工的水冷模块需求能否持续超预期，取决于AI数据中心的资本开支节奏是否维持高位。西太平洋银行的信贷风险是否被低估，取决于澳大利亚就业市场的韧性是否超预期。印度斯坦石油的政策红利能否兑现，取决于印度政府能源改革的执行力度。这些变量在研报中有所提及，但并未展开其二阶影响。

对于希望深入理解这些风险的读者，我们建议进一步研读摩根斯坦利关于这三家公司的完整报告，特别是其中的情景分析和关键假设敏感性测试。这些内容在本文中无法完全展开，但它们是判断投资逻辑可靠性的核心依据。

\subsection{[R033] NOM：印度消费复苏并非全面开花，真正的结构性机会在组织化企业的份额掠夺}
{\small\color{refgray}归类：区域市场：亚洲结构性机会与分化}

NOM：印度消费复苏并非全面开花，真正的结构性机会在组织化企业的份额掠夺

印度消费股在经历了一段时间的沉寂后，似乎正迎来一波新的关注。但这份NOM关于印度4QFY26（2026年第一季度）消费品行业的深度研报，揭示了一个比“复苏”二字更为复杂的图景。市场可能正在定价需求的回暖，但真正被低估的，是供给端正在发生的结构性洗牌。

这份报告的核心判断是：印度消费行业的增长驱动力正在从“普惠式复苏”切换为“组织化企业的份额掠夺”。GST降税带来的短期价格红利固然重要，但更关键的长期变量是，在地缘政治引发的成本通胀和渠道变革中，拥有品牌、规模和供应链优势的组织化企业，正在系统性地从非组织化（unorganized）玩家和地方性品牌手中夺取市场份额。这并非一场所有船只都会水涨船高的行情，而是一场赢家通吃的淘汰赛。

1. 数据印证：组织化企业正在以两倍于行业的速度增长，这是最值得关注的信号

研报中最具说服力的数据并非销售增速，而是组织化部门与非组织化部门之间的增速裂口。报告明确指出，虽然印度快速消费品（FMCG）行业的整体销量增长环比有所放缓，但组织化企业的销量增长却创下了过去四年来的最高纪录。在必需品（Staples）领域，组织化企业的4Q销量增速达到9\%，是其过去八个季度均值（6.7\%）的1.35倍；在食品领域，这一增速更是高达14\%，远超行业平均水平。

这意味着什么？这意味着在整体需求并未出现井喷式爆发的背景下，组织化企业正在以一种前所未有的速度“吃掉”非组织化玩家的市场。报告中的图表（Fig. 7和Fig. 17）清晰地展示了这一剪刀差。这并非简单的周期性回暖，而是一个结构性的市场份额再分配过程。对于投资者而言，关注的重点应从“消费行业好不好”转向“谁在吃谁的份额”。

2. 成本通胀是催化剂，而非灾难：它正在加速淘汰缺乏定价权的玩家

市场通常将地缘政治冲突导致的输入性成本通胀视为消费行业的利空。然而，这份NOM报告提供了另一个视角：通胀恰恰是组织化企业扩大优势的催化剂。报告反复强调，在原材料成本上涨的背景下，组织化企业凭借其强大的品牌议价能力和规模采购优势，能够以可控的方式（中个位数提价）向下游传导成本压力，而对销量几乎没有实质影响。

相反，非组织化的小厂商和地方品牌，由于缺乏品牌忠诚度和供应链韧性，在成本压力下要么难以提价，要么提价即导致客户流失。报告提到，HPC（家庭及个人护理）领域的竞争强度有所上升，但组织化企业依然通过成本效率措施维持了利润率。这背后揭示了一个残酷的商业逻辑：在通胀环境下，定价权是企业的生命线。拥有定价权的企业，通胀是利润的助推器；没有定价权的企业，通胀是生存的绞索。

3. 渠道变革的终局：快速商务正在重塑品牌格局，D2C的威胁已成过去时

报告用相当篇幅讨论了渠道变革，尤其是快速商务（Quick Commerce）的崛起。这并非新话题，但NOM给出了一个值得深思的判断：快速商务平台目前正在成为拥有高市场份额和强品牌认知度的大品牌的“顺风车”。原因在于，平台为了吸引用户，倾向于提供折扣，而大品牌最有能力参与这一游戏，同时其高端产品线（premium portfolio）在新渠道中获得了更多曝光。

更关键的是，报告明确指出，来自直接面向消费者（D2C）品牌的竞争压力正在显著减弱。数据显示，尽管印度D2C品牌数量从2018年的300个激增到2025年的1.1万个，但只有233个品牌（约2\%）的销售额超过了15亿印度卢比，且大部分尚未盈利。在需求疲软和成本压力加大的环境下，新品牌和新产品的生存空间被急剧压缩。这意味着，对于已经建立渠道壁垒和品牌心智的行业龙头而言，来自新进入者的“颠覆”威胁，至少在短期内已经大大降低。这为它们未来的利润率和市场份额提供了额外的安全垫。

4. 报告尚

[... middle omitted ...]

这样依赖农村批发渠道的企业，报告已经指出了“双重定价”对交易增长的短期冲击。农村市场的走向，将是验证当前“份额掠夺”叙事能否持续的核心试金石。

5. 一个值得投资者留意的观察框架：从“销量增长”转向“定价权与利润率韧性”

基于以上分析，投资者在审视印度消费股时，可能需要调整自己的观察框架。过去，市场可能更关注销量的同比增速。但在当前阶段，一个更有效的框架应该是关注两个维度：第一，企业在成本通胀环境下的定价能力，这体现在毛利率和净利率的稳定性上；第二，企业从非组织化市场抢夺份额的能力，这体现在其销量增速与行业平均增速的差值上。

NOM报告中对Marico的分析就是一个绝佳的例证。Marico在经历了此前60\%的大幅提价后，又因椰干（Copra）价格下跌而进行了约10\%的降价。这种灵活的价格调整能力，以及其在新兴渠道（电商、数字优先品牌）的布局，使其成为“兼具传统品牌护城河与未来增长赛道”的复合型标的。这种类型的公司，可能比单纯依赖销量增长的品类更值得给予估值溢价。

\subsection{[R034] NOM：韩国正极材料三巨头的分化，才是理解全球电池供应链重构的真正线索}
{\small\color{refgray}归类：区域市场：亚洲结构性机会与分化}

NOM：韩国正极材料三巨头的分化，才是理解全球电池供应链重构的真正线索

市场对电池材料行业的关注，长期集中在“谁拿到了多少订单”和“产能扩张速度”这两个维度上。但NOM在2026年6月新加坡亚洲投资论坛上与韩国三大正极材料公司——LG化学、L\&F和Ecopro BM——的交流揭示了一个更关键的信号：这三家公司在战略路径上出现了根本性的分化，而这种分化背后，是对全球电池供应链未来十年格局的不同押注。这份报告真正值得看的判断，不是某一家公司的出货量预测，而是这三家公司正在用完全不同的方式回答同一个问题：在欧美本土化生产、LFP与高镍路线并存、中国供应链主导地位被挑战的背景下，韩国正极材料企业的竞争优势究竟从哪里来？

1. Ecopro BM的战略锚点是欧洲高镍需求，但真正的胜负手在印尼镍矿整合

Ecopro BM是目前三家之中战略路径最清晰的一家。NOM报告显示，该公司现有正极材料产能244万吨，其中54万吨位于匈牙利，第一条产线计划于2026年6月投产，第二条在2026年下半年启动。更重要的是，Ecopro BM正将第三条产线转换为NCM生产，目标客户是欧洲的非韩国电池工厂。

这一布局的逻辑非常清晰。欧洲正在重新启动电动汽车消费补贴，英国和德国的政策信号已经明确，同时监管要求供应链本土化。Ecopro BM的高镍正极产品组合——而非中镍或LFP——恰好切中了欧洲市场对续航里程和能量密度的优先需求。但这里有一个容易被忽略的关键点：Ecopro BM的竞争优势并不完全来自其产品本身，而来自Ecopro集团在印尼镍冶炼厂中约10\%的持股，该冶炼厂总产能为15万吨，集团还在评估二期投资。

这意味着什么？Ecopro BM的原材料成本结构将显著优于那些依赖外部镍供应的竞争对手。在正极材料这样一个成本高度受原材料价格驱动的行业，垂直整合能力直接决定了利润率的稳定性。Ecopro BM的短期优先事项——匈牙利工厂爬坡、2027年固态电池正极商业化、印尼镍冶炼扩张——本质上都是在加固这一成本护城河。而该公司明确不优先投资LFP，这一判断本身就在押注高镍路线在中长期仍将是主流。

2. L\&F正在赌LFP的韩国本土化，但技术路线上的“无前驱体”才是真正值得跟踪的变量

L\&F的路径与Ecopro BM形成了鲜明对比。NOM预计L\&F 2026年出货量将同比增长19\%至8万吨，其中约85\%的收入来自LG能源解决方案，最终客户是特斯拉。但真正值得关注的不是这些存量业务，而是L\&F在LFP上的布局。

L\&F计划在韩国建设6万吨LFP正极产能，其中3万吨预计在2026年第三季度实现商业化，剩余3万吨在2027年上半年。首批产品将供应给三星SDI的美国ESS电池业务，这是基于2026年3月签署的长期供应协议。L\&F管理层对LFP正极的盈利前景表示乐观，假设售价为每公斤11美元——尽管他们也承认，整体盈利性仍不如高镍正极。

但这里有一个技术细节值得深入思考。L\&F正在推进“无前驱体”LFP技术。传统正极生产包含三个步骤：制造前驱体、将前驱体与锂混合、高温煅烧。无前驱体LFP意味着跳过混合原料的步骤，在煅烧过程中完成大部分生产。技术挑战在于，需要同时实现金属均匀混合、添加锂并形成晶体结构。L\&F和其他行业同行正在努力消除煅烧过程，目标是降低成本并减少对中国进口前驱体的依赖。

这个技术路线的战略含义远超L\&F一家公司。如果无前驱体LFP能够实现商业化，它将从根本上改变正极材料的成本结构和供应链依赖。目前，中国在前驱体领域占据主导地位，韩国和日本企业高度依赖进口。如果L\&F能够率先突破这一技术瓶颈，它不仅在LFP领域获得成本优势，更重要的是获得了供应链自主权。但NOM报告也暗示，这一技术仍处于开发阶段，商业化时间表尚未明确，这里仍需

[... middle omitted ...]

可能在 年。管理层承认，竞争对手在领先的电动汽车OEM客户方面获得了早期优势，但预计LG化学的地位将逐步改善，目标是实现与竞争对手约50:50的份额分配。

在化学业务方面，管理层提出了重组计划，核心是成立合资公司和优化石脑油裂解装置产能。目标是整合约400万吨的NCC产能，并关闭约100万吨，以应对行业供应过剩和乙烯经济性疲软的问题。管理层认为，通过重组可以实现显著的成本削减。

在股东回报方面，LG化学采取了三个具体举措：将管理层薪酬与ROE超过10\%等指标挂钩；在董事会层面监控NAV折价；重申在未来五年内将LG能源解决方案的持股比例降至70\%的计划。

这三个维度的行动之间存在内在联系。化学业务的重组释放的现金流，需要用于支持正极业务的扩张和股东回报。但问题在于，化学业务的重组效果何时能够体现在财务数据中？正极市占率的追赶需要多长时间？在竞争格局已经发生变化的情况下，50:50的目标是否过于乐观？NOM报告没有完全回答这些问题，这恰恰是投资者需要持续跟踪的核心。

\section{Disclaimer}
{\scriptsize\color{refgray}
This community is a paid learning and information-sharing community created for general educational purposes only. The membership fee is charged solely for access to community discussions, educational content, organization, and operational costs. It is not an advisory fee, management fee, performance fee, brokerage fee, or compensation for personalized financial, investment, legal, tax, or professional advice.\par
I participate and share information solely as an individual and for educational discussion among community members. I am not acting as a financial adviser, investment adviser, broker, dealer, portfolio manager, legal adviser, tax adviser, accountant, fiduciary, or any other licensed professional adviser. Nothing shared in this community constitutes, and should not be interpreted as, financial advice, investment advice, legal advice, tax advice, a recommendation, solicitation, offer, endorsement, instruction, or guarantee to buy, sell, hold, short, trade, or invest in any security, cryptocurrency, fund, derivative, strategy, or other financial product.\par
All content shared in this community, including messages, discussions, links, files, charts, examples, market commentary, watchlists, AI-generated or AI-polished summaries, and third-party materials, is general, impersonal, and educational in nature. It does not take into account any individual's financial situation, investment objectives, risk tolerance, investment horizon, liquidity needs, tax status, legal restrictions, or suitability requirements. No content should be relied upon as suitable for any specific person, account, portfolio, or financial decision.\par
Information shared in this community may be incomplete, inaccurate, outdated, speculative, AI-generated, AI-edited, or based on third-party internet sources that have not been independently verified. I make no representation or warranty as to the accuracy, completeness, timeliness, reliability, or suitability of any information shared.\par
Financial markets involve substantial risk, including the possible loss of principal. Past performance, historical data, hypothetical examples, simulated results, backtests, personal opinions, or educational examples do not guarantee future results. Each member is solely responsible for conducting their own research, exercising independent judgment, and making their own decisions. Before making any financial, investment, legal, or tax decision, members should consult qualified and licensed professionals.\par
Participation in this community, payment of any membership fee, or communication with me or other members does not create any adviser-client, fiduciary, professional, contractual, agency, partnership, employment, or other advisory relationship. By joining or participating in this community, you acknowledge and agree that you are solely responsible for your own decisions, actions, transactions, profits, losses, risks, and consequences, and that I shall not be liable for any direct, indirect, incidental, consequential, financial, legal, tax, or other loss or damage arising from or related to any content, discussion, information, or material shared in this community.\par
}
\end{document}
